% L'Etat et le pouvoir — Philosophie Tle Tech — Cours signé OnBuch+
% Programme officiel MINESEC. Compiler avec : tectonic N07-etat-pouvoir.tex
\newcommand{\DOCMATIERE}{Philosophie}
\newcommand{\DOCNIVEAU}{Tle Tech}
\newcommand{\DOCMODULE}{Philosophie – classe de Terminale technique}
\newcommand{\DOCLECON}{N07}
\newcommand{\DOCTITRE}{L'Etat et le pouvoir}
% OnBuch+ — préambule « cours signés » ENRICHI (compilé avec tectonic / XeLaTeX)
% Système visuel « Pop » d'OnBuch+ : fond crème (jamais blanc), bordures encre
% épaisses, ombres « dures » décalées, titres Archivo Black en capitales.
% Version MATHÉMATIQUES : graphes (pgfplots/tikz), boîtes pédagogiques
% (définition, propriété, méthode, attention, le savais-tu, exercice résolu,
% exercice, TP, bilan), page de garde et signature OnBuch+ ancrée en bas.
%
% Macros attendues dans main.tex AVANT \input{preamble} :
%   \DOCMATIERE  \DOCNIVEAU  \DOCMODULE  \DOCLECON (numéro)  \DOCTITRE
\documentclass[11pt]{article}

\usepackage{fontspec}
\usepackage[french]{babel}
\usepackage[a4paper,margin=2cm,top=3.4cm,bottom=2.8cm,headheight=1.4cm,footskip=1.3cm]{geometry}
\usepackage{amsmath,amssymb,amsfonts}
\usepackage{mathtools} % \xmapsto, \coloneqq... souvent utilisés par les modèles
\usepackage{cancel} % simplifications barrées (bilans de forces)
% \abs{z} (module, valeur absolue) et \norm{u} (norme), utilisés par les modèles
\providecommand{\abs}{}\let\abs\relax\DeclarePairedDelimiter{\abs}{\lvert}{\rvert}
\providecommand{\norm}{}\let\norm\relax\DeclarePairedDelimiter{\norm}{\lVert}{\rVert}
\usepackage[table]{xcolor} % [table] : fournit aussi \rowcolors (alternance de lignes)
\usepackage{pagecolor}
\usepackage{tikz}
\usetikzlibrary{arrows.meta,positioning,calc,shapes.geometric,shapes.misc,shapes.arrows,decorations.pathmorphing,decorations.pathreplacing,decorations.markings,patterns,shadows,mindmap,trees,fit,backgrounds,matrix,intersections,angles,quotes}
\usepackage{pgfplots}
\usepackage[version=4]{mhchem} % équations nucléaires \ce{^{235}_{92}U}
\usepackage[european,siunitx]{circuitikz} % schémas électriques
\pgfplotsset{compat=1.18}
\usepgfplotslibrary{fillbetween,groupplots} % aires entre courbes (calcul intégral)
\usepackage{enumitem}
\usepackage{fancyhdr}
% [most] charge toutes les bibliothèques tcolorbox (listings, minted, raster,
% documentation...) et gonfle inutilement la mémoire TeX sur un document long
% (~300 boîtes) : on ne charge que ce qui est réellement utilisé.
\usepackage[skins,breakable]{tcolorbox}
\usepackage{graphicx}
\usepackage{colortbl}
\usepackage{booktabs}
\usepackage{tabularx}
\usepackage{multirow}
\usepackage{diagbox} % case d'en-tête diagonale des tableaux à double entrée
% Colonnes extensibles centrée / gauche / droite (C, L, R), souvent utilisées par les modèles
\newcolumntype{C}{>{\centering\arraybackslash}X}
\newcolumntype{L}{>{\raggedright\arraybackslash}X}
\newcolumntype{R}{>{\raggedleft\arraybackslash}X}
\usepackage{array}
\usepackage{multicol}
\usepackage{siunitx}
% parse-numbers=false : accepte aussi \SI{2,5 \times 10^{-3}}{...}
\sisetup{locale=FR,output-decimal-marker={,},group-minimum-digits=4,parse-numbers=false}
\DeclareSIUnit\ohm{\ensuremath{\symup{\Omega}}} % U+2126 absent de la police
\DeclareSIUnit\division{div} % réglages d'oscilloscope (ms/div, V/div)
\DeclareSIUnit\tr{tr} % tours (vitesse de rotation, ex. \SI{1500}{\tr\per\minute})
\DeclareSIUnit\molar{M} % concentration molaire (ex. \SI{3}{\molar}, \SI{1}{\micro\molar})
\DeclareSIUnit\an{an} % année (le modèle confond parfois avec \year, un compteur TeX) — ex. \SI{5}{\kilo\meter\per\an}
\DeclareSIUnit\FCFA{FCFA} % franc CFA (ex. \SI{500}{\FCFA\per\kilo\gram})
\DeclareSIUnit\degreeBrix{\textdegree Brix} % teneur en sucre d'un sirop/jus (réfractométrie)
\DeclareSIUnit\month{mois} % le modèle utilise parfois \month, un compteur TeX, comme unité de durée
\DeclareSIUnit\microliter{\micro\liter} % le modèle écrit parfois \microliter en un seul token
\DeclareSIUnit\million{millions} % dénombrement (ex. \SI{15}{\million\per\mL} spermatozoïdes)
\usepackage{qrcode}
% \degree : commande fréquemment utilisée par les modèles pour un angle ou une
% température en dehors de \SI{}{} (non fournie par les paquets chargés).
\providecommand{\degree}{^{\circ}}
% \qed (fin de démonstration), utilisé par les modèles sans amsthm
\providecommand{\qed}{\hfill$\square$}
% \barre : double barre des valeurs interdites dans un tableau de variations (tkz-tab)
\providecommand{\barre}{\ensuremath{\|}}
% environnement proof (amsthm non chargé) utilisé par les modèles
\ifdefined\proof\else\newenvironment{proof}[1][Démonstration]{\par\noindent\textit{#1.}\ }{\qed\par}\fi
\usepackage{lastpage}
\usepackage{hyperref}
\hypersetup{hidelinks}

% ============================================================
% Palette « Pop » OnBuch+ — mêmes hex que la classe P (plus_ui.dart)
% ============================================================
\definecolor{poporange}{HTML}{F59321}
\definecolor{popdark}{HTML}{E07A0C}
\definecolor{popcream}{HTML}{FFF6E8}
\definecolor{popink}{HTML}{1C1714}
\definecolor{poppurple}{HTML}{7A5AE0}
\definecolor{popgreen}{HTML}{2E9E6A}
\definecolor{poppink}{HTML}{E0457B}
\definecolor{popblue}{HTML}{2F7DE1}
\definecolor{popgold}{HTML}{C98A12}
\definecolor{popmuted}{HTML}{8A7F76}
\definecolor{popline}{HTML}{EADFD2}
\definecolor{popgray}{HTML}{8A7F76}
\colorlet{popgrey}{popgray}
% Alias tolérants (le modèle nomme parfois les couleurs différemment)
\colorlet{popwhite}{white}
\colorlet{popblack}{popink}
\colorlet{popred}{poppink}
\colorlet{popyellow}{popgold}
% Teintes claires (fonds de tableaux/encadrés) — le modèle nomme parfois
% les couleurs avec un suffixe L (ex. popblueL) sans les avoir définies.
\colorlet{poporangeL}{poporange!20}
\colorlet{popdarkL}{popdark!20}
\colorlet{poppurpleL}{poppurple!20}
\colorlet{popgreenL}{popgreen!20}
\colorlet{poppinkL}{poppink!20}
\colorlet{popblueL}{popblue!20}
\colorlet{popgoldL}{popgold!20}
\colorlet{popredL}{popred!20}
\colorlet{popgrayL}{popgray!20}
% Teintes claires pour fonds de tableaux / zones
\colorlet{popblueL}{popblue!10!white}
\colorlet{poporangeL}{poporange!14!white}
\colorlet{popgreenL}{popgreen!12!white}
\colorlet{poppurpleL}{poppurple!12!white}
\colorlet{poppinkL}{poppink!12!white}
\colorlet{popgoldL}{popgold!14!white}

\pagecolor{popcream}

% ============================================================
% Typographie
% ============================================================
\defaultfontfeatures{Ligatures=TeX}
\setmainfont{PlusJakartaSans-Medium.ttf}[Path=../../../fonts/,BoldFont=PlusJakartaSans-Bold.ttf]
% Maths en sans-serif (Fira Math) avec chiffres et lettres droites de Plus
% Jakarta, pour que les formules se fondent dans le texte.
\usepackage[math-style=ISO,bold-style=ISO]{unicode-math}
\setmathfont{FiraMath-Regular.otf}[Path=../../../fonts/]
\setmathfont{PlusJakartaSans-Medium.ttf}[Path=../../../fonts/,range={up/num,up/{Latin,latin},\mathup}]
\usepackage{icomma} % virgule décimale sans espace en mode math
% Caractères mathématiques Unicode absents des polices de texte (Plus Jakarta,
% Archivo Black) : rendus par leur équivalent mathématique, en texte comme en
% formule — sinon ils s'affichent en carré vide (ex. « Arithmétique dans ℤ »).
\usepackage{newunicodechar}
\newunicodechar{ℝ}{\ensuremath{\mathbb{R}}}
\newunicodechar{ℤ}{\ensuremath{\mathbb{Z}}}
\newunicodechar{ℂ}{\ensuremath{\mathbb{C}}}
\newunicodechar{ℕ}{\ensuremath{\mathbb{N}}}
\newunicodechar{ℚ}{\ensuremath{\mathbb{Q}}}
\newunicodechar{𝔻}{\ensuremath{\mathbb{D}}}
\newunicodechar{α}{\ensuremath{\alpha}}
\newunicodechar{β}{\ensuremath{\beta}}
\newunicodechar{γ}{\ensuremath{\gamma}}
\newunicodechar{δ}{\ensuremath{\delta}}
\newunicodechar{ε}{\ensuremath{\varepsilon}}
\newunicodechar{θ}{\ensuremath{\theta}}
\newunicodechar{λ}{\ensuremath{\lambda}}
\newunicodechar{μ}{\ensuremath{\mu}}
\newunicodechar{σ}{\ensuremath{\sigma}}
\newunicodechar{τ}{\ensuremath{\tau}}
\newunicodechar{φ}{\ensuremath{\varphi}}
\newunicodechar{ψ}{\ensuremath{\psi}}
\newunicodechar{ω}{\ensuremath{\omega}}
\newunicodechar{Σ}{\ensuremath{\Sigma}}
% majuscules produites par \MakeUppercase dans les titres (α-aminés -> Α)
\newunicodechar{Α}{\ensuremath{\alpha}}
\newunicodechar{Β}{\ensuremath{\beta}}
\newunicodechar{Γ}{\ensuremath{\Gamma}}
\newunicodechar{Δ}{\ensuremath{\Delta}}
\newunicodechar{Ε}{\ensuremath{\varepsilon}}
\newunicodechar{Θ}{\ensuremath{\Theta}}
\newunicodechar{Λ}{\ensuremath{\Lambda}}
\newunicodechar{Μ}{\ensuremath{\mu}}
\newunicodechar{Τ}{\ensuremath{\tau}}
\newunicodechar{Φ}{\ensuremath{\Phi}}
\newunicodechar{Ψ}{\ensuremath{\Psi}}
\newunicodechar{Ω}{\ensuremath{\Omega}}
\newunicodechar{↦}{\ensuremath{\mapsto}}
\newunicodechar{∈}{\ensuremath{\in}}
\newunicodechar{∉}{\ensuremath{\notin}}
\newunicodechar{∧}{\ensuremath{\wedge}}
\newunicodechar{⇔}{\ensuremath{\Leftrightarrow}}
\newunicodechar{⟺}{\ensuremath{\Longleftrightarrow}}
\newunicodechar{⇒}{\ensuremath{\Rightarrow}}
\newunicodechar{⟹}{\ensuremath{\Longrightarrow}}
\newunicodechar{∘}{\ensuremath{\circ}}
\newunicodechar{∪}{\ensuremath{\cup}}
\newunicodechar{∩}{\ensuremath{\cap}}
\newunicodechar{≡}{\ensuremath{\equiv}}
\newunicodechar{⊂}{\ensuremath{\subset}}
\newunicodechar{⊥}{\ensuremath{\perp}}
\newunicodechar{∃}{\ensuremath{\exists}}
\newunicodechar{∀}{\ensuremath{\forall}}
\newunicodechar{∛}{\ensuremath{\sqrt[3]{\,}}}
\newunicodechar{∜}{\ensuremath{\sqrt[4]{\,}}}
\newunicodechar{⃗}{\ensuremath{{}^{\to}}}
\newunicodechar{ⁿ}{\ifmmode^{n}\else\textsuperscript{\ensuremath{n}}\fi}
\newunicodechar{ˣ}{\ifmmode^{x}\else\textsuperscript{\ensuremath{x}}\fi}
\newunicodechar{ᵇ}{\ifmmode^{b}\else\textsuperscript{\ensuremath{b}}\fi}
\newunicodechar{ʳ}{\ifmmode^{r}\else\textsuperscript{\ensuremath{r}}\fi}
\newunicodechar{ᵘ}{\ifmmode^{u}\else\textsuperscript{\ensuremath{u}}\fi}
\newunicodechar{ʸ}{\ifmmode^{y}\else\textsuperscript{\ensuremath{y}}\fi}
\newunicodechar{ᵃ}{\ifmmode^{a}\else\textsuperscript{\ensuremath{a}}\fi}
\newunicodechar{ᵉ}{\ifmmode^{e}\else\textsuperscript{\ensuremath{e}}\fi}
\newunicodechar{ᵏ}{\ifmmode^{k}\else\textsuperscript{\ensuremath{k}}\fi}
\newunicodechar{ᵖ}{\ifmmode^{p}\else\textsuperscript{\ensuremath{p}}\fi}
\newunicodechar{ᴺ}{\ifmmode^{N}\else\textsuperscript{\ensuremath{N}}\fi}
\newunicodechar{ᶜ}{\ifmmode^{c}\else\textsuperscript{\ensuremath{c}}\fi}
\newunicodechar{ʰ}{\ifmmode^{h}\else\textsuperscript{\ensuremath{h}}\fi}
\newunicodechar{ᵠ}{\ifmmode^{q}\else\textsuperscript{\ensuremath{q}}\fi}
\newunicodechar{ₙ}{\ifmmode_{n}\else\textsubscript{\ensuremath{n}}\fi}
\newunicodechar{ₐ}{\ifmmode_{a}\else\textsubscript{\ensuremath{a}}\fi}
\newunicodechar{ᵢ}{\ifmmode_{i}\else\textsubscript{\ensuremath{i}}\fi}
\newunicodechar{ᵣ}{\ifmmode_{r}\else\textsubscript{\ensuremath{r}}\fi}
\newunicodechar{ₖ}{\ifmmode_{k}\else\textsubscript{\ensuremath{k}}\fi}
\newunicodechar{ᵦ}{\ifmmode_{\beta}\else\textsubscript{\ensuremath{\beta}}\fi}
\newunicodechar{ₓ}{\ifmmode_{x}\else\textsubscript{\ensuremath{x}}\fi}
\newfontfamily\popdisplay{ArchivoBlack-Regular.ttf}[Path=../../../fonts/]
\newfontfamily\poplogo{SpaceGrotesk-Bold.ttf}[Path=../../../fonts/]
\newfontfamily\popsemi{PlusJakartaSans-SemiBold.ttf}[Path=../../../fonts/]
\newfontfamily\popxbold{PlusJakartaSans-ExtraBold.ttf}[Path=../../../fonts/]
\newfontfamily\popxboldls{PlusJakartaSans-ExtraBold.ttf}[Path=../../../fonts/,LetterSpace=3.0]

\setlist[enumerate]{leftmargin=1.7em,itemsep=5pt,topsep=4pt}
\setlist[itemize]{leftmargin=1.7em,itemsep=4pt,topsep=4pt,label={\color{poporange}\textbullet}}
\setlength{\parindent}{0pt}
\setlength{\parskip}{5pt}
\renewcommand{\arraystretch}{1.25}

% pgfplots : style Pop par défaut
\pgfplotsset{
  every axis/.append style={axis line style={popink,line width=1pt},tick style={popink},
    grid style={popline,line width=0.5pt},label style={font=\small},tick label style={font=\footnotesize}},
  popcurve/.style={poporange,line width=1.8pt,no markers,smooth},
}
% Même style disponible dans un \draw TikZ simple (hors pgfplots)
\tikzset{popcurve/.style={poporange,line width=1.8pt}}
\tikzset{popgrid/.style={popline,very thin}}
\colorlet{popcurve}{poporange} % le nom du style est parfois utilisé comme couleur

% ============================================================
% Pill / squiggle
% ============================================================
\newcommand{\poppill}[3]{%
  \tikz[baseline=-0.55ex]\node[draw=popink,line width=1.2pt,rounded corners=2.6pt,%
    fill=#2,inner xsep=6.5pt,inner ysep=3pt]{%
    \popxboldls\fontsize{7}{7}\selectfont\color{#3}\MakeUppercase{#1}};%
}
\newcommand{\popsquiggle}[1]{%
  \tikz[baseline=-0.4ex]{
    \draw[#1,line width=2.6pt,line cap=round]
      plot [smooth] coordinates {(0,0.09) (0.34,0.01) (0.68,0.21) (1.02,0.01) (1.36,0.10)};
  }%
}
% Mot-clé surligné (marqueur orange sous le texte)
\newcommand{\cle}[1]{{\popxbold\color{popdark}#1}}

% ============================================================
% En-tête / pied
% ============================================================
\pagestyle{fancy}
\fancyhf{}
\renewcommand{\headrulewidth}{0pt}
\renewcommand{\footrulewidth}{0pt}
\lhead{\raisebox{-0.15ex}{\poplogo\fontsize{13}{13}\selectfont\color{popink}OnBuch\color{poporange}+}}
\rhead{\poppill{\DOCMATIERE\ \textbullet\ \DOCNIVEAU\ \textbullet\ Leçon \DOCLECON}{white}{popink}}
\lfoot{\popxbold\fontsize{7.6}{7.6}\selectfont\color{popmuted}\MakeUppercase{Cours signé OnBuch+}\ \textbullet\ {\popsemi\DOCTITRE}}
\rfoot{\tikz[baseline=-0.6ex]\node[rounded corners=8.5pt,fill=popink,minimum height=17pt,minimum width=17pt,inner xsep=5pt,inner sep=0]{\popxbold\fontsize{8}{8}\selectfont\color{popcream}\thepage\,/\,\pageref*{LastPage}};}

% ============================================================
% Titres
% ============================================================
\newcommand{\chaptitre}[2]{%
  \begin{tcolorbox}[enhanced,colback=poporange,colframe=popink,boxrule=2.6pt,arc=11pt,
      drop shadow=popink,left=15pt,right=15pt,top=12pt,bottom=11pt,before skip=3pt,after skip=18pt]
    \poppill{#2}{popink}{popcream}\par\vspace{9pt}
    {\raggedright\hyphenpenalty=10000\exhyphenpenalty=10000\popdisplay\fontsize{22}{24}\selectfont\color{popcream}\MakeUppercase{#1}\par}
  \end{tcolorbox}
}
% Section numérotée : pastille orange avec le numéro + titre Archivo
\newcounter{popsec}
\newcommand{\coursec}[1]{%
  \stepcounter{popsec}\setcounter{popsub}{0}%
  \par\vspace{16pt}\noindent
  \begin{minipage}{\linewidth}
  \tikz[baseline=-0.7ex]\node[draw=popink,line width=1.6pt,fill=poporange,rounded corners=4pt,minimum size=22pt,inner sep=0]{\popdisplay\fontsize{12}{12}\selectfont\color{popcream}\thepopsec};%
  \hspace{8pt}{\popdisplay\fontsize{15}{17}\selectfont\color{popink}\MakeUppercase{#1}}\par
  \vspace{4pt}\noindent{\color{poporange}\rule{36pt}{3.6pt}}
  \end{minipage}\par\nopagebreak\vspace{8pt}
}
\newcounter{popsub}
\newcommand{\courssub}[1]{%
  \stepcounter{popsub}%
  \par\vspace{9pt}\noindent{\popxbold\fontsize{11.5}{13}\selectfont\color{popink}{\color{poporange}\thepopsec.\thepopsub}\hspace{6pt}#1}\par\nopagebreak\vspace{4pt}
}

% ============================================================
% Boîtes pédagogiques — même recette : fond blanc, bordure épaisse colorée,
% ombre dure encre, badge plein. Titre optionnel [#1].
% ============================================================
\tcbset{popbase/.style={breakable,enhanced,colback=white,boxrule=2.3pt,arc=9pt,
  shadow={3pt}{-3pt}{0pt}{popink},left=14pt,right=14pt,top=11pt,bottom=11pt,before skip=11pt,after skip=15pt,
  fontupper=\popsemi\fontsize{10.6}{14.5}\selectfont\color{popink},
  fontlower=\popsemi\fontsize{10.6}{14.5}\selectfont\color{popink}}}
% Coche dessinée (le glyphe ✓ n'existe pas dans les polices du document)
\newcommand{\popcheck}[1]{\tikz[baseline=-0.5ex]\draw[#1,line width=1.6pt,line cap=round,line join=round](-0.13,0.0)--(-0.04,-0.09)--(0.13,0.1);}
\providecommand{\coche}{\popcheck{popgreen}}
\providecommand{\resultat}[1]{\par\smallskip\noindent\fcolorbox{popgreen}{popgreenL}{\parbox{\dimexpr\linewidth-2\fboxsep-2\fboxrule\relax}{\textbf{Résultat.} #1}}\par\smallskip}
\newcommand{\popboxhead}[3]{\poppill{#1}{#2}{white}\ifx\relax#3\relax\else\hspace{6pt}{\popxbold\fontsize{10.6}{12}\selectfont\color{popink}#3}\fi\par\vspace{7pt}}

\newtcolorbox{aretenir}[1][]{popbase,colframe=popgreen,before upper={\popboxhead{À retenir}{popgreen}{#1}}}
\newtcolorbox{exemplebox}[1][]{popbase,colframe=poporange,before upper={\popboxhead{Exemple}{poporange}{#1}}}
\newtcolorbox{definition}[1][]{popbase,colframe=popblue,before upper={\popboxhead{Définition}{popblue}{#1}}}
\newtcolorbox{propriete}[1][]{popbase,colframe=poppurple,before upper={\popboxhead{Propriété}{poppurple}{#1}}}
\newtcolorbox{methode}[1][]{popbase,colframe=popgold,before upper={\popboxhead{Méthode}{popgold}{#1}}}
\newtcolorbox{attention}[1][]{popbase,colframe=poppink,before upper={\popboxhead{Attention}{poppink}{#1}}}
\newtcolorbox{experience}[1][]{popbase,colframe=popblue,colback=popblueL,before upper={\popboxhead{Expérience}{popblue}{#1}}}
% « Le savais-tu ? » : carte violette pleine (contexte, culture, Cameroun)
\newtcolorbox{savaistu}[1][]{popbase,colback=poppurple,colframe=popink,
  fontupper=\popsemi\fontsize{10.4}{14.2}\selectfont\color{white},
  before upper={\poppill{Le savais-tu\,?}{white}{poppurple}\ifx\relax#1\relax\else\hspace{6pt}{\popxbold\color{white}#1}\fi\par\vspace{7pt}}}
% Exercice résolu : énoncé en haut, \tcblower puis la solution sur fond crème
\newcounter{popexo}\newcounter{popexr}
\newtcolorbox{exoresolu}[1][]{popbase,colframe=poporange,colbacklower=popcream,
  segmentation style={popink,line width=1.2pt,dashed},
  before upper={\stepcounter{popexr}\popboxhead{Exercice résolu \thepopexr}{poporange}{#1}},
  before lower={\poppill{Solution}{popink}{popcream}\par\vspace{6pt}}}
% Exercice (non résolu en place) — la correction est dans la partie Corrigés
\newtcolorbox{exercice}[1][]{popbase,colframe=popink,
  before upper={\stepcounter{popexo}\popboxhead{Exercice \thepopexo}{popink}{#1}}}
% Corrigé
\newtcolorbox{corrige}[1][]{popbase,colframe=popgreen,colback=popgreenL,
  before upper={\popboxhead{Corrigé}{popgreen}{#1}}}
% Objectifs / prérequis / bilan
\newtcolorbox{objectifs}{popbase,colframe=popink,colback=poporangeL,
  before upper={\popboxhead{Objectifs de la leçon}{popink}{}}}
\newtcolorbox{prerequis}{popbase,colframe=popink,colback=popblueL,
  before upper={\popboxhead{Prérequis}{popblue}{}}}
\newtcolorbox{bilan}{popbase,colframe=popink,colback=white,boxrule=2.8pt,
  before upper={\popboxhead{Fiche bilan}{poporange}{L'essentiel à connaître pour l'examen}}}
% Figure encadrée avec légende
\newtcolorbox{popfigure}[1][]{enhanced,breakable=false,colback=white,colframe=popline,boxrule=1.4pt,arc=8pt,
  left=10pt,right=10pt,top=10pt,bottom=8pt,before skip=10pt,after skip=12pt,halign=center,
  lower separated=false,halign lower=center,
  fontlower=\popsemi\fontsize{9}{11.5}\selectfont\color{popmuted},
  #1}
\newcommand{\legende}[1]{\par\vspace{6pt}{\popsemi\fontsize{9}{11.5}\selectfont\color{popmuted}#1}}

% ============================================================
% Signature OnBuch+
% ============================================================
\newcommand{\poplogoinline}[1]{{\poplogo\fontsize{#1}{#1}\selectfont\color{popink}OnBuch\color{poporange}+}}
\newcommand{\signatureonbuch}{%
  \begin{tcolorbox}[enhanced,colback=popink,colframe=popink,boxrule=2.2pt,arc=10pt,
      drop shadow=poporange,left=14pt,right=14pt,top=10pt,bottom=10pt,before skip=0pt,after skip=0pt]
    \tikz[baseline=-0.7ex]\node[circle,fill=popgreen,draw=popcream,line width=1.3pt,minimum size=22pt,inner sep=0]{\popcheck{white}};%
    \hspace{9pt}\begin{minipage}[c]{0.62\linewidth}
      {\popxbold\fontsize{12}{13}\selectfont\color{popcream}Signé\ \textbullet\ {\poplogo OnBuch\color{poporange}+}}\par\vspace{2pt}
      {\raggedright\popsemi\fontsize{8}{10.5}\selectfont\color{popline}\MakeUppercase{Équipe pédagogique OnBuch+}\\\MakeUppercase{Conforme au programme officiel MINESEC}\par}
    \end{minipage}\hfill
    {\poplogo\fontsize{22}{22}\selectfont\color{popcream}OnBuch\color{poporange}+}
  \end{tcolorbox}%
}

% ============================================================
% Page de garde
% #1 titre · #2 matière · #3 niveau · #4 module · #5 n° leçon · #6 durée ·
% #7 description / objectifs (texte ou itemize)
% ============================================================
\newcommand{\pagegarde}[7]{%
  \thispagestyle{empty}
  \newgeometry{a4paper,margin=1.7cm,top=1.6cm,bottom=1.4cm}%
  \begin{tikzpicture}[remember picture,overlay]
    \fill[poporange!18] ([xshift=-3.2cm,yshift=-3.6cm]current page.north east) circle (4.6cm);
    \fill[poppurple!14] ([xshift=2.4cm,yshift=7.6cm]current page.south west) circle (3.8cm);
  \end{tikzpicture}%
  {\poplogo\fontsize{30}{30}\selectfont\color{popink}OnBuch\color{poporange}+}\hfill
  \raisebox{0.6ex}{\poppill{Cours signé \textbullet\ Premium}{popink}{popcream}}\par
  \vspace{1pt}\popsquiggle{poppurple}\par
  \vspace{8pt}
  \begin{tcolorbox}[enhanced,colback=poporange,colframe=popink,boxrule=2.8pt,arc=13pt,
      drop shadow=popink,left=18pt,right=18pt,top=12pt,bottom=13pt,after skip=10pt]
    \poppill{#2 \textbullet\ #3}{popink}{popcream}\hspace{5pt}\poppill{#4}{white}{popink}\par\vspace{8pt}
    {\popdisplay\fontsize{14}{14}\selectfont\color{popink}LEÇON #5}\par\vspace{4pt}
    {\raggedright\hyphenpenalty=10000\exhyphenpenalty=10000\popdisplay\fontsize{25}{27}\selectfont\color{popcream}\MakeUppercase{#1}\par}
    \vspace{8pt}
    {\popxbold\fontsize{9.5}{9.5}\selectfont\color{popink}\MakeUppercase{Durée conseillée : #6}}
  \end{tcolorbox}
  \begin{tcolorbox}[enhanced,colback=white,colframe=popink,boxrule=2.2pt,arc=10pt,
      drop shadow=popink,left=16pt,right=16pt,top=10pt,bottom=10pt,after skip=8pt]
    \poppill{Dans cette leçon}{poppurple}{white}\par\vspace{5pt}
    \popsemi\fontsize{9.3}{12.6}\selectfont\color{popink}#7
  \end{tcolorbox}
  \noindent\begin{minipage}[t]{0.487\linewidth}
    \begin{tcolorbox}[enhanced,colback=popgreen,colframe=popink,boxrule=2.2pt,arc=10pt,
        drop shadow=popink,left=11pt,right=11pt,top=10pt,bottom=10pt,height=3.1cm,valign=center]
      \tcbox[colback=white,colframe=popink,boxrule=1.3pt,arc=5pt,boxsep=0pt,nobeforeafter,
        left=3pt,right=3pt,top=3pt,bottom=3pt,valign=center]{\qrcode[height=1.9cm]{https://play.google.com/store/apps/details?id=com.luvvix.onbuch&hl=fr}}%
      \hspace{8pt}\begin{minipage}[c]{0.5\linewidth}
        {\popdisplay\fontsize{11}{12.5}\selectfont\color{white}\MakeUppercase{Télécharge OnBuch}}\par\vspace{4pt}
        {\raggedright\popsemi\fontsize{8.4}{11}\selectfont\color{white}Gratuit \textbullet\ Google Play\\Tuteur IA, annales, concours, résultats\par}
      \end{minipage}
    \end{tcolorbox}
  \end{minipage}\hfill
  \begin{minipage}[t]{0.487\linewidth}
    \begin{tcolorbox}[enhanced,colback=poppurple,colframe=popink,boxrule=2.2pt,arc=10pt,
        drop shadow=popink,left=11pt,right=11pt,top=10pt,bottom=10pt,height=3.1cm,valign=center]
      \tikz[baseline=-0.7ex]\node[circle,fill=white,draw=popink,line width=1.3pt,minimum size=1.6cm,inner sep=0]{\popdisplay\fontsize{24}{24}\selectfont\color{poppurple}+};%
      \hspace{8pt}\begin{minipage}[c]{0.6\linewidth}
        {\popdisplay\fontsize{11}{12.5}\selectfont\color{white}\MakeUppercase{Passe à OnBuch+}}\par\vspace{4pt}
        {\raggedright\popsemi\fontsize{8.4}{11}\selectfont\color{white}Tous les cours signés, Tuteur IA illimité, fiches PDF — onglet OnBuch+ de l'app\par}
      \end{minipage}
    \end{tcolorbox}
  \end{minipage}
  \vspace{8pt}
  \popcontenu
  \vfill
  \signatureonbuch
  \restoregeometry
  \newpage
}

% Rangée « Ce cours contient » (page de garde)
\newcommand{\poptile}[3]{% #1 couleur #2 gros texte #3 légende
  \begin{tcolorbox}[enhanced,colback=#1,colframe=popink,boxrule=2pt,arc=9pt,shadow={2.5pt}{-2.5pt}{0pt}{popink},
      left=6pt,right=6pt,top=9pt,bottom=9pt,halign=center,valign=center,height=1.75cm,width=\linewidth,nobeforeafter]
    {\popdisplay\fontsize{12.5}{13}\selectfont\color{white}\MakeUppercase{#2}}\par\vspace{3pt}
    {\centering\popsemi\fontsize{7.8}{9.5}\selectfont\color{white}#3\par}
  \end{tcolorbox}}
\newcommand{\popcontenu}{%
  \par\noindent{\popxboldls\fontsize{7.5}{7.5}\selectfont\color{popmuted}CE COURS SIGNÉ CONTIENT}\par\vspace{7pt}
  \noindent\begin{minipage}[t]{0.235\linewidth}\poptile{popblue}{Cours}{complet, illustré, pas à pas}\end{minipage}\hfill
  \begin{minipage}[t]{0.235\linewidth}\poptile{poporange}{Méthodes}{et exercices résolus}\end{minipage}\hfill
  \begin{minipage}[t]{0.235\linewidth}\poptile{poppink}{Exercices}{type examen + corrigés}\end{minipage}\hfill
  \begin{minipage}[t]{0.235\linewidth}\poptile{popgreen}{Fiche bilan}{l'essentiel à retenir}\end{minipage}\par}

% Sommaire
\newtcolorbox{sommaire}{enhanced,colback=white,colframe=popink,boxrule=2.2pt,arc=10pt,
  drop shadow=popink,left=18pt,right=18pt,top=13pt,bottom=13pt,before skip=6pt,after skip=20pt,
  before upper={\poppill{Sommaire}{poppurple}{white}\par\vspace{9pt}\popsemi\fontsize{10.6}{14.5}\selectfont\color{popink}}}

% Signature de fin de cours — poussée tout en bas de la dernière page
\newcommand{\findecours}{%
  \par\vfill\nopagebreak
  \begin{center}\popsquiggle{poporange}\end{center}\vspace{4pt}
  \signatureonbuch
}
\AtBeginDocument{\def\llbracket{\mathopen{[\mkern-3.5mu[}}\def\rrbracket{\mathclose{]\mkern-3.5mu]}}} % intervalles d'entiers (glyphe ⟦ absent de la police)
% Glyphes absents de Fira Math : redéfinis à partir de glyphes disponibles
\AtBeginDocument{%
  \def\setminus{\mathbin{\backslash}}\let\smallsetminus\setminus
  \def\vdots{\mathord{\vbox{\baselineskip3.2pt\lineskiplimit0pt\kern4pt\hbox{.}\hbox{.}\hbox{.}}}}%
  \def\ddots{\mathinner{\mkern1mu\raise6.5pt\hbox{.}\mkern2mu\raise3.6pt\hbox{.}\mkern2mu\raise0.7pt\hbox{.}\mkern1mu}}%
  \def\triangle{\mathord{\scalebox{0.9}{$\symup{\Delta}$}}}%
  \def\nabla{\mathord{\raisebox{\depth}{\scalebox{1}[-1]{$\symup{\Delta}$}}}}%
  \def\checkmark{\popcheck{popdark}}%
  \def\star{\mathbin{\ast}}\def\bigstar{\mathord{\ast}}%
  \def\diamond{\mathbin{\scalebox{0.6}{$\Diamond$}}}%
  \def\bigcup{\mathop{\mathchoice{\raisebox{-0.3em}{\scalebox{1.7}{$\cup$}}}{\scalebox{1.25}{$\cup$}}{\cup}{\cup}}\displaylimits}%
  \def\bigcap{\mathop{\mathchoice{\raisebox{-0.3em}{\scalebox{1.7}{$\cap$}}}{\scalebox{1.25}{$\cap$}}{\cap}{\cap}}\displaylimits}%
  \def\lesseqqgtr{\mathrel{\lessgtr}}\def\bigoplus{\mathop{\oplus}\displaylimits}%
}
\providecommand{\ssi}{si et seulement si} % abréviation fréquente
\AtBeginDocument{\providecommand{\wideparen}{\overparen}} % arc de cercle

\begin{document}

\pagegarde{L'Etat et le pouvoir}{Philosophie}{Tle Tech}{Philosophie – classe de Terminale technique}{N07}{3 séances (fiche de progression harmonisée)}{Cette leçon interroge la notion d'État : sa définition, son fondement (contrat social, force, tradition), les formes de gouvernements qu'il peut prendre, ainsi que son rôle et sa finalité. À travers des exemples camerounais et africains, elle conduit l'élève à distinguer État, gouvernement et nation, et à justifier une position argumentée sur la légitimité du pouvoir politique.
\vspace{4pt}
\begin{multicols}{2}\small
\begin{itemize}
\item À la fin de cette leçon, je sais définir l'État et identifier ses éléments constitutifs (territoire, population, souveraineté)
\item À la fin de cette leçon, je sais distinguer État, gouvernement, nation et régime politique
\item À la fin de cette leçon, je sais expliquer les principaux fondements de l'État (force, droit divin, contrat social) à partir des textes de Hobbes, Locke et Rousseau
\item À la fin de cette leçon, je sais classifier les formes de gouvernements selon le nombre de gouvernants et selon l'exercice du pouvoir
\item À la fin de cette leçon, je sais analyser un document (texte philosophique, article de la Constitution camerounaise) pour en dégager la conception de l'État qu'il défend
\item À la fin de cette leçon, je sais appliquer les notions de l'unité à des situations concrètes de la vie politique camerounaise et africaine
\end{itemize}\end{multicols}}

\begin{sommaire}
\begin{multicols}{2}
\begin{enumerate}
\item Qu'est-ce que l'État ? Définition et éléments constitutifs
\item Le fondement de l'État : d'où vient son autorité ?
\item Les formes de gouvernements
\item Rôle et finalité de l'État
\item Atelier d'application : analyser des situations concrètes
\item Activité d'intégration
\item Méthodes et exercices résolus
\item Exercices
\item Corrigés des exercices
\item Fiche bilan
\end{enumerate}
\end{multicols}
\end{sommaire}

\begin{objectifs}
\begin{itemize}
\item À la fin de cette leçon, je sais définir l'État et identifier ses éléments constitutifs (territoire, population, souveraineté)
\item À la fin de cette leçon, je sais distinguer État, gouvernement, nation et régime politique
\item À la fin de cette leçon, je sais expliquer les principaux fondements de l'État (force, droit divin, contrat social) à partir des textes de Hobbes, Locke et Rousseau
\item À la fin de cette leçon, je sais classifier les formes de gouvernements selon le nombre de gouvernants et selon l'exercice du pouvoir
\item À la fin de cette leçon, je sais analyser un document (texte philosophique, article de la Constitution camerounaise) pour en dégager la conception de l'État qu'il défend
\item À la fin de cette leçon, je sais appliquer les notions de l'unité à des situations concrètes de la vie politique camerounaise et africaine
\item À la fin de cette leçon, je sais dégager le rôle et la finalité de l'État (ordre, justice, bien commun, protection des libertés)
\item À la fin de cette leçon, je sais rédiger et justifier une démarche argumentée, une réponse ou une conclusion en dissertation ou en explication de texte
\end{itemize}
\end{objectifs}

\begin{prerequis}
\begin{itemize}
\item Notion de pouvoir et d'autorité (vie de classe, famille, chef traditionnel)
\item Notion de liberté et de loi abordée dans les leçons précédentes de Terminale
\item Notions de citoyenneté et de droits de l'homme vues en Éducation civique (collège)
\item Connaissance générale de l'organisation politique du Cameroun (Président, Assemblée, Constitution de 1996)
\item Notion de justice (leçon antérieure du programme de Terminale)
\end{itemize}
\end{prerequis}

\newpage

\coursec{Qu'est-ce que l'État ? Définition et éléments constitutifs}

\courssub{Situation d'accroche et problématique}

À Yaoundé comme à Bafoussam, chaque élève possède une carte nationale d'identité, paie parfois des taxes via ses parents, circule sur des routes construites par l'État et respecte des lois votées en son nom. Le matin, l'enseignant qui fait l'appel est un fonctionnaire de l'État ; le soir, le policier qui règle la circulation au carrefour agit au nom de l'État ; le certificat de naissance conservé précieusement à la maison a été délivré par l'État. Pourtant, lors des élections ou des débats sur les réseaux sociaux, certains jeunes Camerounais affirment : « L'État ne sert à rien, il ne fait que nous opprimer. » D'autres répondent : « Sans État, ce serait la loi du plus fort. » Qui a raison ? Qu'est-ce qui fonde l'autorité de l'État, et à quoi sert-il réellement ?

Cette leçon nous invite à interroger philosophiquement cette réalité que nous côtoyons chaque jour sans jamais la questionner. Avant de pouvoir juger l'État — le défendre ou le critiquer —, il faut d'abord répondre à une question préalable, en apparence simple mais en réalité redoutable : \emph{qu'est-ce que l'État ?} De quels éléments est-il constitué ? En quoi se distingue-t-il du gouvernement, de la nation, de la société ? Telle est la problématique de cette première section : \textbf{qu'est-ce qui fait qu'un État est un État ?}

\courssub{Le sens du mot « État » : étymologie et usage courant}

Commençons, comme tout bon philosophe, par interroger le mot lui-même. Le terme « État » vient du latin \emph{status}, qui signifie « manière d'être », « position stable », « situation ». Ce mot désigne à l'origine ce qui \emph{se tient debout} (\emph{stare} : être debout), ce qui demeure, ce qui dure. L'étymologie est déjà philosophiquement riche : l'État est pensé, dès l'origine du mot, comme ce qui apporte la \cle{stabilité} et la \cle{permanence} dans la vie collective. Là où les individus naissent et meurent, là où les modes changent et les gouvernants se succèdent, l'État, lui, \emph{demeure}.

Dans le langage courant, le mot « État » possède plusieurs sens qu'il faut savoir distinguer pour éviter les confusions :

\begin{itemize}
\item On parle de l'« état » de quelque chose pour désigner sa condition, sa manière d'être momentanée : « l'état des routes après la saison des pluies », « l'état de santé d'un malade ». C'est le sens ancien, encore vivant.
\item On parle de « l'État » (avec une majuscule) pour désigner l'organisation politique qui gouverne un pays : « l'État camerounais », « l'État français ».
\item On dit familièrement « l'État » pour désigner l'administration ou les pouvoirs publics : « c'est l'État qui doit payer », « les employés de l'État ».
\end{itemize}

C'est évidemment le deuxième sens qui nous intéresse en philosophie politique. Mais remarquons que les trois sens partagent une même idée de fond : celle d'une \emph{situation établie, ordonnée et durable}. Dire qu'un pays « a un État », c'est dire qu'il possède une organisation stable qui ordonne la vie collective.

\begin{attention}[Ne pas confondre les sens du mot]
Dans une dissertation, employer le mot « État » sans préciser son sens expose au hors-sujet. Quand le sujet porte sur « l'État », il s'agit toujours de l'organisation politique souveraine, jamais de l'« état » d'une chose. Prenez l'habitude, dans votre introduction, de définir clairement le terme dès la première définition du sujet.
\end{attention}

\courssub{Définition de l'État}

Après l'intuition étymologique, venons-en à la définition rigoureuse. Observons d'abord ce qui distingue un pays doté d'un État d'une simple foule d'individus vivant sur un même sol. Des milliers de personnes peuvent occuper un même territoire sans former un État : il suffit de penser aux populations en guerre civile où plus aucune autorité ne s'impose. Inversement, un État existe dès lors qu'une autorité unique s'exerce effectivement sur un groupe humain installé sur un territoire déterminé. C'est cette observation qui conduit à la définition classique.

\begin{definition}[L'État]
L'\cle{État} est une \textbf{organisation politique souveraine} qui exerce l'\textbf{autorité suprême} sur une \textbf{population} installée sur un \textbf{territoire déterminé}. Il se manifeste par un ensemble d'institutions (gouvernement, administration, armée, police, justice) chargées d'édicter et de faire respecter les lois, de prélever les impôts et d'assurer les services publics.
\end{definition}

Analysons cette définition terme à terme, car chaque mot compte :

\begin{itemize}
\item \textbf{« Organisation politique »} : l'État n'est ni une famille, ni une entreprise, ni une association. C'est une structure instituée dont la fonction est de gouverner, c'est-à-dire de diriger la vie collective en prenant des décisions obligatoires pour tous.
\item \textbf{« Souveraine »} : l'État ne reconnaît au-dessus de lui aucune autorité supérieure. Nous y reviendrons en détail.
\item \textbf{« Autorité suprême »} : les décisions de l'État s'imposent à tous ceux qui se trouvent sur son territoire, qu'ils soient citoyens ou étrangers de passage.
\item \textbf{« Population » et « territoire déterminé »} : l'État n'exerce pas son pouvoir dans le vide ni sur n'importe qui : il a des assises concrètes, humaines et géographiques.
\end{itemize}

\begin{aretenir}[L'essentiel de la définition]
Retenez la formule en trois points : l'État, c'est \textbf{un pouvoir souverain} exercé sur \textbf{une population} et \textbf{un territoire}. Ces trois éléments sont indissociables : si l'un manque, il n'y a plus d'État au sens propre.
\end{aretenir}

\courssub{Les trois éléments constitutifs de l'État}

La tradition juridique et politique identifie trois éléments sans lesquels aucun État ne peut exister. Examinons-les successivement.

\textbf{Premier élément : un territoire délimité.} L'État exerce son autorité sur un espace géographique précis, borné par des \cle{frontières}. Le territoire comprend le sol, le sous-sol (richesses minières), l'espace aérien situé au-dessus et, pour les États côtiers, les eaux territoriales. Les frontières ne sont pas de simples lignes : elles marquent la limite où cesse la loi d'un État et où commence celle d'un autre. Quand un voyageur passe du Cameroun au Nigeria à Mfum, il change littéralement de monde juridique : les lois, la monnaie, la langue officielle et les autorités ne sont plus les mêmes.

\textbf{Deuxième élément : une population.} Un territoire vide ne fait pas un État. Il faut un groupe humain qui y réside de manière stable et qui est soumis à l'autorité de l'État. Cette population n'a pas besoin d'être homogène : le Cameroun rassemble plus de 250 groupes ethniques et deux langues officielles, et il forme pourtant un seul État. Ce qui compte, c'est le lien juridique qui relie les personnes à l'État : la \cle{nationalité} ou la simple résidence sur le territoire.

\textbf{Troisième élément : un pouvoir souverain.} C'est l'élément décisif. Il faut une autorité organisée capable d'édicter des lois, de les faire appliquer et de trancher les conflits, sans dépendre d'aucun pouvoir supérieur.

\begin{definition}[La souveraineté]
La \cle{souveraineté} est le caractère d'un pouvoir qui ne reconnaît \textbf{aucun pouvoir supérieur}, ni à l'intérieur ni à l'extérieur. \emph{À l'intérieur} (souveraineté interne), l'État est seul habilité à faire les lois, à lever l'impôt et à rendre la justice sur son territoire. \emph{À l'extérieur} (souveraineté externe), l'État est indépendant : il n'obéit à aucun autre État et traite avec eux d'égal à égal par la diplomatie et les traités.
\end{definition}

\begin{popfigure}
\begin{center}
\begin{tikzpicture}[scale=1]
% Triangle des trois éléments
\coordinate (T) at (-3.4,-2.2);
\coordinate (P) at (3.4,-2.2);
\coordinate (S) at (0,3.2);
\draw[line width=1.2pt,popblue] (T) -- (P) -- (S) -- cycle;
% Nœuds
\node[circle,fill=popblueL,draw=popblue,line width=1pt,minimum size=2.5cm,align=center,font=\small\bfseries] at (T) {TERRITOIRE\\ \footnotesize\mdseries frontières, sol,\\ \footnotesize\mdseries espace aérien};
\node[circle,fill=popgreenL,draw=popgreen,line width=1pt,minimum size=2.5cm,align=center,font=\small\bfseries] at (P) {POPULATION\\ \footnotesize\mdseries habitants,\\ \footnotesize\mdseries nation};
\node[circle,fill=poporangeL,draw=poporange,line width=1pt,minimum size=2.5cm,align=center,font=\small\bfseries] at (S) {SOUVERAINETÉ\\ \footnotesize\mdseries autorité suprême\\ \footnotesize\mdseries interne et externe};
% Centre
\node[star,star points=8,fill=poppink,draw=poppink!70!black,minimum size=2.1cm,align=center,font=\small\bfseries,text=white] at (0,-0.2) {L'ÉTAT};
\end{tikzpicture}
\end{center}
\legende{Le triangle des trois éléments constitutifs de l'État. Retirez un seul sommet — territoire, population ou souveraineté — et l'édifice entier s'effondre : il n'y a plus d'État.}
\end{popfigure}

\courssub{Le monopole de la violence légitime selon Max Weber}

Mais comment l'État fait-il respecter ses décisions ? Qu'est-ce qui distingue sa loi d'une simple recommandation ? Le sociologue et philosophe allemand \textbf{Max Weber} (1864-1922) a donné à cette question une réponse devenue classique : l'État est « une communauté humaine qui revendique, avec succès, le \emph{monopole de la violence physique légitime} » à l'intérieur d'un territoire donné.

Décortiquons cette formule célèbre :

\begin{itemize}
\item \textbf{« Violence physique »} : l'État dispose de moyens de contrainte — la police qui arrête, la justice qui condamne, la prison qui enferme, l'armée qui défend. Contrairement à une association ou à une famille, l'État peut user de la force pour faire respecter ses lois.
\item \textbf{« Monopole »} : l'État est le \emph{seul} à pouvoir légitimement user de cette force. Si un particulier se fait justice lui-même — par exemple en frappant le voleur qui a dérobé sa moto au marché Mokolo —, il commet lui-même un délit. La vengeance privée est interdite ; seule la justice de l'État peut punir.
\item \textbf{« Légitime »} : cette violence n'est pas arbitraire. Elle est encadrée par le droit et reconnue comme juste parce qu'elle s'exerce selon des procédures légales (arrestation, procès, jugement contradictoire).
\end{itemize}

L'intuition de Weber est profonde : l'État ne supprime pas la violence, il la \emph{canalise} et la \emph{régule}. En interdisant aux individus de se faire justice eux-mêmes et en concentrant la force entre les mains d'institutions contrôlées par le droit, l'État fait sortir la société de la « guerre de tous contre tous » que décrivait déjà le philosophe anglais Thomas Hobbes. Sans ce monopole, chaque conflit de voisinage, chaque dette impayée, chaque accident de la route dégénérerait en vengeance privée sans fin.

\begin{exemplebox}[Police, armée, justice : les bras armés de l'État]
Au Cameroun, ce monopole s'incarne concrètement dans trois institutions : la \textbf{police} et la \textbf{gendarmerie}, qui assurent l'ordre intérieur ; l'\textbf{armée}, qui garantit l'intégrité du territoire face aux menaces extérieures ; la \textbf{justice}, qui tranche les litiges et punit les infractions au nom du peuple camerounais. Aucune autre organisation — ni chef traditionnel, ni milice privée, ni groupe religieux — ne peut légalement emprisonner ou punir.
\end{exemplebox}

\courssub{Application au Cameroun : un État concret}

La théorie prend chair lorsqu'on l'applique à notre propre pays. Vérifions que le Cameroun satisfait bien aux trois critères constitutifs.

\begin{exemplebox}[Le Cameroun, un État en chiffres]
Le Cameroun compte environ \textbf{28 millions d'habitants}, répartis sur un territoire de \SI{475650}{km^2}, organisé en \textbf{10 régions} et \textbf{360 communes}. Autant d'échelons par lesquels l'État est présent dans la vie quotidienne : la mairie qui délivre l'acte de naissance, le lycée public qui scolarise, l'hôpital de district qui soigne, le poste de gendarmerie qui sécurise. Ses institutions républicaines sont issues de la \textbf{Constitution du 18 janvier 1996}, qui fait du Cameroun une « République décentralisée, une et indivisible ».
\end{exemplebox}

\begin{popfigure}
\begin{center}
\begin{tikzpicture}[scale=0.62]
% Contour très simplifié du Cameroun
\draw[line width=1.3pt,popblue,fill=popblueL!50] 
  (0,0) -- (1.2,-1.6) -- (3.2,-1.8) -- (4.6,-0.8) -- (5.0,1.0) 
  -- (4.4,2.6) -- (5.2,4.2) -- (4.6,6.0) -- (3.4,7.6) -- (2.2,9.4) 
  -- (1.6,8.2) -- (0.8,6.6) -- (1.4,5.0) -- (0.4,3.4) -- (-0.4,1.8) -- cycle;
% Capitale
\node[star,star points=5,fill=poppink,minimum size=0.45cm,inner sep=0pt] at (2.3,2.2) {};
\node[font=\small\bfseries,anchor=west] at (2.6,2.2) {Yaoundé (capitale)};
% Autres villes
\fill[popdark] (0.6,1.4) circle (0.09); \node[font=\footnotesize,anchor=east] at (0.5,1.4) {Douala};
\fill[popdark] (1.2,4.6) circle (0.09); \node[font=\footnotesize,anchor=east] at (1.1,4.6) {Bafoussam};
\fill[popdark] (1.9,8.6) circle (0.09); \node[font=\footnotesize,anchor=west] at (2.0,8.6) {Maroua};
\fill[popdark] (3.0,6.6) circle (0.09); \node[font=\footnotesize,anchor=west] at (3.1,6.6) {Garoua};
% Nord
\draw[-{Stealth},line width=1pt,popdark] (5.8,8.6) -- (5.8,9.8);
\node[font=\small\bfseries] at (5.8,10.2) {N};
% Échelle indicative
\draw[line width=1pt] (0,-2.6) -- (2,-2.6);
\draw[line width=1pt] (0,-2.7) -- (0,-2.5);
\draw[line width=1pt] (2,-2.7) -- (2,-2.5);
\node[font=\footnotesize] at (1,-3.0) {$\approx$ 300 km};
\end{tikzpicture}
\end{center}
\legende{Carte très simplifiée du Cameroun : le territoire (environ \SI{475650}{km^2}), délimité par des frontières internationales, avec sa capitale Yaoundé et quelques grandes villes. Le territoire est le premier élément constitutif de l'État.}
\end{popfigure}

\courssub{Distinctions fondamentales : État, gouvernement, nation, société civile}

Pour penser clairement, il faut maintenant distinguer l'État de trois réalités voisines avec lesquelles on le confond trop souvent.

\textbf{État et gouvernement.} Le \cle{gouvernement} est l'équipe dirigeante qui administre l'État à un moment donné : président, ministres, préfets. Or les gouvernants \emph{passent} — ils sont élus, nommés, remplacés, démissionnent —, tandis que l'État \emph{demeure}. Quand un nouveau gouvernement est formé à Yaoundé, les lois, les frontières, les traités internationaux et l'administration restent en place. Le gouvernement est le pilote ; l'État est le navire.

\begin{attention}[L'erreur fréquente : confondre État et gouvernement]
C'est la confusion la plus répandue, y compris dans les copies d'examen. Dire « je déteste l'État » parce qu'on critique le gouvernement en place est un contresens : on peut s'opposer à un gouvernement tout en respectant l'État et ses institutions. Retenez la formule : \textbf{le gouvernement est l'équipe dirigeante temporaire, l'État est l'institution permanente}. Critiquer un ministre, c'est critiquer un gouvernant ; remettre en cause l'existence même de l'organisation étatique, c'est une tout autre position (l'anarchisme), qu'il faut alors assumer et argumenter.
\end{attention}

\textbf{État et nation.} La \cle{nation} est une communauté humaine unie par une culture, une histoire, une langue ou une volonté de vivre ensemble. C'est une réalité \emph{culturelle et historique}, tandis que l'État est une réalité \emph{juridique et politique}. Les deux ne coïncident pas toujours : il existe des nations sans État (les Kurdes, répartis entre plusieurs pays) et des États multinationaux. Le Cameroun illustre la construction d'une nation unique à partir d'une grande diversité ethnique : plus de 250 groupes, mais un seul État et une devise commune, « Paix – Travail – Patrie ».

\textbf{État et société civile.} La \cle{société civile} désigne l'ensemble des groupements que les citoyens forment librement en dehors de l'État : familles, associations, syndicats, églises, entreprises, tontines, clubs sportifs, organisations non gouvernementales. La société civile est le domaine de l'initiative privée et de la liberté d'association ; l'État est le domaine de l'autorité publique et de la loi obligatoire. Une démocratie saine suppose un équilibre : une société civile vigoureuse qui exprime les besoins des citoyens, et un État qui encadre sans étouffer.

\begin{popfigure}
\begin{center}
\begin{tikzpicture}[scale=0.95]
% Société : grand ensemble
\draw[line width=1.2pt,popmuted,fill=popcream] (0,0) ellipse (5.6 and 3.4);
\node[font=\small\bfseries,text=popmuted] at (0,3.0) {LA SOCIÉTÉ};
% État
\draw[line width=1.2pt,popblue,fill=popblueL!60] (-2.4,0.4) ellipse (2.1 and 1.5);
\node[font=\small\bfseries,text=popblue] at (-2.4,1.35) {L'ÉTAT};
\node[font=\footnotesize,align=center] at (-2.4,0.7) {institutions permanentes :\\ lois, armée, justice,\\ administration};
% Gouvernement : sous-ensemble de l'État
\draw[line width=1pt,poporange,fill=poporangeL] (-2.4,-0.35) ellipse (1.15 and 0.55);
\node[font=\footnotesize\bfseries,text=poporange!80!black] at (-2.4,-0.35) {Gouvernement};
% Nation
\draw[line width=1.2pt,popgreen,fill=popgreenL!60] (2.6,0.4) ellipse (2.1 and 1.5);
\node[font=\small\bfseries,text=popgreen!70!black] at (2.6,1.35) {LA NATION};
\node[font=\footnotesize,align=center] at (2.6,0.5) {communauté culturelle\\ et historique : langues,\\ traditions, mémoire};
% Société civile
\node[font=\footnotesize,align=center,text=popmuted!30!black] at (0,-2.3) {SOCIÉTÉ CIVILE : familles, associations,\\ syndicats, tontines, églises, entreprises};
\end{tikzpicture}
\end{center}
\legende{Quatre réalités à ne pas confondre. Le gouvernement est une partie \emph{temporaire} de l'État ; la nation est une communauté culturelle ; la société civile regroupe les initiatives libres des citoyens. Ces ensembles se recouvrent partiellement sans jamais se confondre.}
\end{popfigure}

\courssub{L'État de droit : un État soumis au droit}

Reste une dernière précision, capitale. Dire que l'État est souverain ne signifie pas qu'il peut tout faire n'importe comment. Dans un \cle{État de droit}, l'État est \emph{lui-même soumis au droit} : il n'agit que dans les limites fixées par les lois, et ses agents peuvent être poursuivis s'ils abusent de leur pouvoir. C'est la différence décisive entre un État de droit et un État arbitraire, où le souverain décide selon son bon plaisir.

Au sommet des normes se trouve la \textbf{Constitution} : au Cameroun, celle du 18 janvier 1996. En dessous viennent les lois votées par le Parlement, puis les décrets et arrêtés pris par le pouvoir exécutif, enfin les décisions administratives. Cette \emph{hiérarchie des normes} garantit qu'aucune décision d'un sous-préfet ne peut contredire une loi, et qu'aucune loi ne peut contredire la Constitution. Le citoyen peut ainsi invoquer le droit contre l'administration elle-même : c'est là le propre de l'État de droit.

\begin{savaistu}[« L'État, c'est moi »]
L'expression « L'État, c'est moi », attribuée au roi de France Louis XIV (1638-1715), illustre la conception \emph{absolutiste} du pouvoir : le souverain s'identifie totalement à l'État, concentre tous les pouvoirs et ne rend de comptes à personne. C'est exactement l'inverse de l'État de droit moderne, où le dirigeant n'est qu'un serviteur temporaire de l'État, soumis comme tout citoyen à la Constitution. Le passage de l'absolutisme à l'État de droit est l'une des grandes conquêtes de l'histoire politique.
\end{savaistu}

\courssub{Exercice résolu et entraînement}

\begin{exoresolu}[Identifier les éléments constitutifs]
\textbf{Énoncé.} Un groupe de réfugiés traverse la frontière et s'installe temporairement dans un camp au Cameroun. Ce groupe forme-t-il un État ? Justifiez votre réponse en mobilisant les trois éléments constitutifs.
\tcblower
\textbf{Solution détaillée.} Non, ce groupe ne forme pas un État, car il ne réunit aucun des trois éléments constitutifs de manière propre. \textbf{Premièrement}, le groupe n'a pas de territoire qui lui appartienne : il séjourne sur le territoire camerounais, qui demeure sous la souveraineté de l'État camerounais. \textbf{Deuxièmement}, la population du camp est certes un groupe humain, mais elle n'est pas organisée en communauté politique stable dotée d'une nationalité propre. \textbf{Troisièmement}, et c'est le point décisif, ce groupe ne dispose d'aucun pouvoir souverain : il ne peut ni édicter des lois, ni lever l'impôt, ni rendre la justice ; il relève au contraire de l'autorité du Cameroun et du droit international. La réponse doit donc articuler les trois critères (territoire, population, souveraineté) et montrer que l'absence de souveraineté suffit à elle seule à exclure l'existence d'un État.
\end{exoresolu}

\begin{exercice}[À vous de jouer]
\begin{enumerate}
\item Définissez l'État et la souveraineté en une phrase chacun, sans recopier le cours mot pour mot.
\item Expliquez, en cinq lignes maximum, la formule de Max Weber : « l'État détient le monopole de la violence légitime ».
\item Un élève affirme : « Puisque le gouvernement a changé, l'État camerounais n'existe plus. » Corrigez cette erreur en vous appuyant sur la distinction État/gouvernement.
\item Pourquoi dit-on que le Cameroun est un État de droit ? Donnez deux arguments.
\end{enumerate}
\end{exercice}

\begin{corrige}[Exercice : À vous de jouer]
\begin{enumerate}
\item L'État est une organisation politique souveraine exerçant l'autorité suprême sur une population et un territoire déterminés. La souveraineté est le caractère d'un pouvoir qui ne reconnaît aucun pouvoir supérieur, ni à l'intérieur ni à l'extérieur.
\item L'État est le seul à pouvoir user légitimement de la force (police, justice, armée) pour faire respecter les lois ; toute violence privée, comme la vengeance ou le châtiment par soi-même, est interdite et punie.
\item L'élève confond État et gouvernement. Le gouvernement est l'équipe dirigeante temporaire : il change sans affecter l'existence de l'État. L'État, institution permanente (lois, territoire, administration, Constitution), demeure à travers les changements de gouvernants.
\item Le Cameroun est un État de droit parce que, d'une part, il dispose d'une Constitution (celle du 18 janvier 1996) placée au sommet d'une hiérarchie des normes à laquelle toutes les autorités sont soumises ; d'autre part, ses agents, y compris les plus hauts dirigeants, doivent respecter la loi et peuvent répondre de leurs actes devant la justice.
\end{enumerate}
\end{corrige}

\begin{aretenir}[Bilan de la section]
L'État est une organisation politique souveraine qui exerce l'autorité suprême sur une population et un territoire déterminés. Il repose sur trois éléments indissociables : le territoire, la population et la souveraineté. Selon Max Weber, il se caractérise par le monopole de la violence légitime. Il faut le distinguer du gouvernement (temporaire), de la nation (culturelle) et de la société civile (libre). Enfin, dans un État de droit, l'État est lui-même soumis au droit, la Constitution occupant le sommet de la hiérarchie des normes. Reste à savoir d'où vient cette autorité : c'est l'objet de la section suivante.
\end{aretenir}

\coursec{Le fondement de l'État : d'où vient son autorité ?}

Dans la section précédente, nous avons défini l'État par ses trois éléments constitutifs : un \cle{territoire}, une \cle{population} et un \cle{pouvoir souverain} qui s'exerce sur les deux premiers. Mais cette définition soulève aussitôt une question plus profonde et plus inquiétante : pourquoi ce pouvoir a-t-il le \emph{droit} de commander ? Un gendarme peut arrêter un chauffard à Yaoundé, un percepteur peut exiger l'impôt, un juge peut condamner un citoyen à la prison. Qu'est-ce qui distingue ces actes de ceux d'un bandit qui, lui aussi, use de la force ? Autrement dit : qu'est-ce qui rend le pouvoir de l'État \cle{légitime} ? C'est le problème du \cle{fondement} de l'État : d'où vient son autorité ?

\courssub{Le problème : pourquoi obéissons-nous à l'État ?}

Commençons par une expérience de pensée simple. Imaginez qu'un inconnu armé vous arrête dans la rue et vous réclame de l'argent : vous lui obéissez par peur, mais vous estimez qu'il n'a \emph{aucun droit} sur vous. Imaginez maintenant qu'un agent de l'État vous réclame le paiement d'une taxe : vous obéissez aussi, mais cette fois la plupart des citoyens reconnaissent qu'il est \emph{dans son droit}. La différence ne tient donc pas à la force — les deux peuvent contraindre — mais à la \cle{légitimité}, c'est-à-dire à la justification morale et juridique du pouvoir.

\begin{definition}[Légitimité]
La \cle{légitimité} est le caractère d'un pouvoir qui est reconnu comme conforme au droit et à la justice par ceux qui lui obéissent. Un pouvoir légitime n'a pas besoin de recourir constamment à la violence : il est obéi parce qu'on le juge fondé en droit.
\end{definition}

Le problème philosophique est donc le suivant : \emph{quelle est la source de la légitimité de l'État ?} La tradition philosophique a donné trois grandes réponses : la force, le droit divin (ou la tradition), et le contrat social. Examinons-les successivement.

\courssub{Première réponse : la force et la conquête}

La réponse la plus ancienne et la plus brutale est celle-ci : l'État tient son pouvoir de la \cle{force}. Celui qui commande est celui qui a vaincu. Les grands empires de l'histoire sont nés de conquêtes militaires : le vainqueur impose sa loi au vaincu, et son autorité repose sur sa capacité à punir les désobéissants. On parle alors de pouvoir fondé sur la \cle{domination}.

Cette thèse a l'avantage du réalisme : il est vrai que beaucoup d'États sont nés de la guerre. Mais elle se heurte à une objection décisive, formulée avec une grande clarté par \textbf{Jean-Jacques Rousseau} dans le \emph{Du contrat social} (1762) : \og la force ne crée pas le droit \fg. En effet :

\begin{itemize}
\item Si la force fondait le droit, alors tout pouvoir fort serait légitime, y compris celui du voleur qui vous dévalise sous la menace.
\item Obéir par peur n'est pas obéir par devoir : dès que la force faiblit, l'obéissance cesse. Le pouvoir fondé sur la seule force est donc fragile et instable.
\item La force est un \emph{fait} ; le droit est une \emph{norme}. On ne peut pas déduire logiquement ce qui \emph{doit être} de ce qui \emph{est}.
\end{itemize}

\begin{attention}[Piège à éviter]
Ne confondez pas \emph{la force comme moyen} et \emph{la force comme fondement}. Tout État, même légitime, dispose de la force (police, armée) pour faire respecter la loi. Mais la force n'est alors qu'un \emph{instrument} au service du droit, non la \emph{source} du droit. Dire « l'État a la force » est vrai ; dire « l'État est légitime parce qu'il est fort » est faux.
\end{attention}

\courssub{Deuxième réponse : le droit divin et la tradition}

Une deuxième réponse affirme que le pouvoir vient de \cle{Dieu} ou de la \cle{tradition}. Dans l'Europe monarchique, les rois régnaient \og de droit divin \fg : ils tenaient leur couronne de Dieu, et non des hommes. Désobéir au roi était donc un péché autant qu'un crime. Le roi de France Louis XIV pouvait ainsi déclarer que les rois sont les lieutenants de Dieu sur terre.

Cette conception n'est pas propre à l'Europe. Au Cameroun, les \cle{chefferies traditionnelles} (chefferies bamiléké et bamoun à l'Ouest, lamidats du Nord, sultanats) fondent leur autorité sur la \cle{coutume} : le chef est légitime parce qu'il est l'héritier des ancêtres, le gardien des traditions et l'intermédiaire entre les vivants et les esprits. Son autorité n'a pas besoin d'élections : elle repose sur la lignée et le sacré.

Cette réponse a une force réelle : elle enracine le pouvoir dans des croyances partagées et lui donne une grande stabilité. Mais elle comporte des limites décisives :

\begin{itemize}
\item Elle rend toute contestation impossible : si le pouvoir vient de Dieu, critiquer le souverain devient un sacrilège, ce qui ouvre la voie à la tyrannie.
\item Elle repose sur une foi que tous ne partagent pas : dans une société pluraliste, on ne peut fonder le droit de tous sur la religion de quelques-uns.
\item Elle confond l'ancienneté avec la justice : une coutume ancienne n'est pas nécessairement juste.
\end{itemize}

\courssub{Troisième réponse : le contrat social}

La réponse moderne, élaborée aux XVII\textsuperscript{e} et XVIII\textsuperscript{e} siècles, renverse la perspective : l'État ne vient ni de la force, ni de Dieu, mais du \cle{consentement} des individus eux-mêmes. C'est la théorie du \cle{contrat social}.

\begin{definition}[Le contrat social]
Le \cle{contrat social} est la convention par laquelle les individus renoncent à une partie de leur liberté naturelle en échange de la sécurité et des droits garantis par l'État. L'autorité politique est donc fondée sur un accord, réel ou supposé, entre les membres de la société.
\end{definition}

Pour penser ce contrat, les philosophes imaginent ce qu'était la vie des hommes \emph{avant} l'État : c'est ce qu'ils appellent l'\cle{état de nature}. Ils comparent ensuite cette situation à celle de l'\cle{état civil}, la vie sous l'autorité de l'État. Mais les trois grands théoriciens du contrat — Hobbes, Locke et Rousseau — ne décrivent ni le même état de nature, ni le même contrat, ni le même pouvoir.

\begin{attention}[Ne pas confondre]
L'\cle{état de nature} n'est pas un état sauvage ou primitif que l'on aurait observé historiquement. C'est une \emph{fiction théorique}, une expérience de pensée : on imagine l'homme privé d'État, de lois et de juges, afin de comprendre ce que l'État nous apporte et pourquoi nous l'acceptons. Ne dites jamais dans une copie que Hobbes ou Rousseau « décrivent la préhistoire ».
\end{attention}

\textbf{Hobbes : la sécurité au prix de la liberté.} Dans le \emph{Léviathan} (1651), \textbf{Thomas Hobbes} décrit l'état de nature comme une \og guerre de tous contre tous \fg. Sans autorité commune pour trancher les conflits, chacun a droit sur tout, chacun craint l'autre, et la vie de l'homme est, selon sa formule célèbre, \og solitaire, misérable, dangereuse, animale et brève \fg. Pour échapper à cette insécurité mortelle, les hommes concluent un pacte : ils \emph{aliènent} (abandonnent) leur liberté à un \cle{souverain} tout-puissant, le Léviathan, qui garantit la paix par la crainte qu'il inspire. Chez Hobbes, le contrat produit donc un pouvoir \emph{absolu} : le prix de la sécurité est l'obéissance quasi totale.

\begin{savaistu}[Pourquoi « Léviathan » ?]
Le titre de l'ouvrage de Hobbes renvoie au \emph{Léviathan}, un monstre marin gigantesque de la Bible (livre de Job). Hobbes fait de ce monstre l'image de l'État : un être artificiel immense, composé de tous les citoyens, dont la puissance impressionnante protège chacun en décourageant toute agression. La couverture originale du livre montre d'ailleurs un géant dont le corps est formé d'une multitude de petits personnages : le peuple lui-même.
\end{savaistu}

\textbf{Locke : un contrat pour protéger les droits naturels.} Dans le \emph{Traité du gouvernement civil} (1690), \textbf{John Locke} propose une vision moins sombre. L'état de nature n'est pas une guerre permanente, mais il reste \emph{incommode} : les hommes y possèdent des \cle{droits naturels} — la \emph{vie}, la \emph{liberté} et la \emph{propriété} — que personne n'est chargé de protéger efficacement. Le contrat social a donc pour but de garantir ces droits en créant un juge impartial : l'État. Conséquence capitale : si le gouvernement viole les droits qu'il était chargé de protéger, le contrat est rompu et le peuple a le \cle{droit de résistance}, voire de renverser ce gouvernement. Le pouvoir est donc \emph{limité} et \emph{conditionnel}.

\textbf{Rousseau : la souveraineté du peuple.} Dans le \emph{Du contrat social} (1762), \textbf{Jean-Jacques Rousseau} radicalise la logique du consentement. Pour lui, le contrat authentique ne transfère pas la souveraineté à un monarque : la \cle{souveraineté appartient au peuple}, qui ne peut ni la céder ni se faire représenter. Les lois sont légitimes lorsqu'elles expriment la \cle{volonté générale}.

\begin{definition}[La volonté générale]
La \cle{volonté générale} est, selon Rousseau, la volonté du corps politique visant l'intérêt commun. Elle se distingue de la simple somme des volontés particulières (que Rousseau appelle \og volonté de tous \fg) : ce n'est pas l'addition des égoïsmes, mais ce que veulent les citoyens en tant qu'ils pensent au bien de la cité tout entière.
\end{definition}

\begin{exemplebox}[La Constitution camerounaise, héritière de Rousseau]
L'article 2 de la Constitution de la République du Cameroun du 18 janvier 1996 affirme : \og la souveraineté nationale appartient au peuple camerounais qui l'exerce par voie de référendum ou par l'intermédiaire de ses représentants \fg. Cette formule est un héritage direct de la théorie rousseauiste du contrat social : le pouvoir des gouvernants n'est qu'une délégation du peuple, seul souverain. Les élections et le référendum sont les procédures concrètes par lesquelles ce consentement s'exprime.
\end{exemplebox}

\begin{popfigure}
\begin{center}
\begin{tikzpicture}[
  box/.style={draw=popblue, thick, rounded corners, fill=popblueL, text width=3.1cm, align=center, font=\small, inner sep=4pt},
  head/.style={draw=popdark, thick, rounded corners, fill=popblue!45, text width=3.1cm, align=center, font=\small\bfseries, inner sep=4pt},
  lab/.style={font=\footnotesize\itshape, text=popmuted, align=center},
  node distance=0.35cm]
\node[head, align=center] (h1){Hobbes\\ \emph{Léviathan} (1651)};
\node[head, right=0.6cm of h1, align=center] (h2){Locke\\ \emph{Traité du gouvernement} (1690)};
\node[head, right=0.6cm of h2, align=center] (h3){Rousseau\\ \emph{Du contrat social} (1762)};
\node[box, below=0.5cm of h1] (n1) {État de nature : guerre de tous contre tous, vie insécure};
\node[box, below=0.5cm of h2] (n2) {État de nature : liberté avec droits naturels (vie, liberté, propriété) mal protégés};
\node[box, below=0.5cm of h3] (n3) {État de nature : indépendance naturelle, mais insécurité des biens et conflits latents};
\node[box, below=0.35cm of n1] (c1) {Contrat : aliénation totale de la liberté au souverain};
\node[box, below=0.35cm of n2] (c2) {Contrat : délégation limitée pour protéger les droits naturels};
\node[box, below=0.35cm of n3] (c3) {Contrat : chacun s'unit à tous et n'obéit qu'à la volonté générale};
\node[box, below=0.35cm of c1, fill=poporangeL, draw=poporange] (p1) {Pouvoir absolu : sécurité contre obéissance};
\node[box, below=0.35cm of c2, fill=poporangeL, draw=poporange] (p2) {Pouvoir limité : droit de résistance si les droits sont violés};
\node[box, below=0.35cm of c3, fill=poporangeL, draw=poporange] (p3) {Souveraineté populaire : le peuple fait les lois};
\draw[->, thick, popblue] (h1) -- (n1); \draw[->, thick, popblue] (h2) -- (n2); \draw[->, thick, popblue] (h3) -- (n3);
\draw[->, thick, popblue] (n1) -- (c1); \draw[->, thick, popblue] (n2) -- (c2); \draw[->, thick, popblue] (n3) -- (c3);
\draw[->, thick, poporange] (c1) -- (p1); \draw[->, thick, poporange] (c2) -- (p2); \draw[->, thick, poporange] (c3) -- (p3);
\end{tikzpicture}
\end{center}
\legende{Comparaison des trois théories du contrat social : état de nature, contenu du contrat et forme de pouvoir qui en résulte.}
\end{popfigure}

\courssub{Bilan : trois fondements, une légitimité moderne}

\begin{popfigure}
\begin{center}
\begin{tikzpicture}[
  grow=right,
  level 1/.style={sibling distance=3.2cm, level distance=3.6cm},
  level 2/.style={sibling distance=1.5cm, level distance=3.8cm},
  racine/.style={draw=popdark, very thick, rounded corners, fill=popblue!45, align=center, font=\small\bfseries, text width=2.6cm, inner sep=4pt},
  branche/.style={draw=popblue, thick, rounded corners, fill=popblueL, align=center, font=\small, text width=3cm, inner sep=3pt},
  feuille/.style={draw=popmuted, rounded corners, fill=popcream, align=center, font=\footnotesize, text width=3.4cm, inner sep=3pt},
  edge from parent/.style={draw=popblue, thick, ->}]
\node[racine] {Fondements de l'État}
  child { node[branche, align=center]{Contrat social\\ (consentement)}
    child { node[feuille] {Rousseau : souveraineté populaire, volonté générale} }
    child { node[feuille] {Locke : protection des droits naturels} }
    child { node[feuille] {Hobbes : sécurité par le souverain} } }
  child { node[branche] {Droit divin et tradition}
    child { node[feuille] {Rois de droit divin ; chefferies traditionnelles (coutume, lignée)} } }
  child { node[branche] {Force et conquête}
    child { node[feuille] {Domination du vainqueur ; critiquée par Rousseau : la force ne crée pas le droit} } };
\end{tikzpicture}
\end{center}
\legende{Arbre des fondements de l'État : trois grandes réponses au problème de la légitimité, la théorie du contrat social se subdivisant en trois versions.}
\end{popfigure}

Tirons le bilan de notre enquête. La \emph{force} ne peut fonder le droit : elle contraint sans légitimer. Le \emph{droit divin} et la \emph{tradition} ont longtemps justifié les pouvoirs, et les chefferies camerounaises montrent que cette légitimité coutumière reste vivante ; mais ils interdisent la contestation et reposent sur des croyances non partagées par tous. Seule la théorie du \emph{contrat social} fonde l'autorité de l'État sur ce qui peut valoir pour tout homme raisonnable : le \cle{consentement}.

\begin{aretenir}[L'essentiel]
\begin{itemize}
\item La \cle{légitimité} est ce qui distingue le pouvoir de l'État de la pure violence.
\item La force ne crée pas le droit (Rousseau) : obéir par peur n'est pas obéir par devoir.
\item Le droit divin et la tradition fondent l'autorité sur Dieu ou la coutume (monarchies sacrées, chefferies traditionnelles).
\item Le \cle{contrat social} fonde l'État sur le consentement : Hobbes (sécurité, pouvoir absolu), Locke (droits naturels, pouvoir limité), Rousseau (souveraineté populaire, volonté générale).
\item La légitimité moderne repose sur le consentement exprimé par les \emph{élections} et la \emph{Constitution} (art. 2 de la Constitution camerounaise de 1996).
\end{itemize}
\end{aretenir}

\begin{methode}[Analyser un texte philosophique sur l'État]
Face à un extrait de Hobbes, Locke ou Rousseau au commentaire de texte :
\begin{enumerate}
\item \textbf{Identifier la thèse} : quelle réponse l'auteur donne-t-il à la question du fondement de l'État ?
\item \textbf{Repérer les arguments} : quelles raisons, quels exemples, quelles expériences de pensée (état de nature) justifient cette thèse ?
\item \textbf{Dégager la conception de l'État} défendue : pouvoir absolu, limité ou populaire ? Quel rapport entre liberté et sécurité ?
\item \textbf{Confronter l'auteur aux autres} : en quoi sa position diffère-t-elle des deux autres théories du contrat, et des fondements par la force ou la tradition ?
\end{enumerate}
\end{methode}

\begin{exoresolu}[La force fonde-t-elle le droit ?]
\textbf{Énoncé.} Un usurpateur s'empare du pouvoir par les armes et déclare : « Je suis le plus fort, donc j'ai le droit de gouverner. » En vous appuyant sur Rousseau, montrez en quelques lignes pourquoi ce raisonnement est invalide, puis indiquez ce qui, selon la théorie du contrat social, pourrait rendre ce pouvoir légitime.
\tcblower
\textbf{Solution détaillée.} Le raisonnement de l'usurpateur confond un fait (être le plus fort) avec un droit (avoir le droit de commander). Or, comme le montre Rousseau, la force ne crée pas le droit : si la force fondait la légitimité, alors tout brigand serait légitime dès qu'il est le plus fort, et l'obéissance cesserait dès que la force faiblirait. Obéir par peur n'est pas obéir par devoir ; c'est céder à une contrainte, non reconnaître une autorité. Un pouvoir fondé sur la seule force est donc à la fois injuste et instable. Selon la théorie du contrat social, ce pouvoir ne deviendrait légitime que s'il reposait sur le \emph{consentement} des citoyens : par exemple s'il était approuvé par une constitution librement adoptée et confirmé par des élections, de sorte que les lois expriment la volonté générale (Rousseau) et protègent les droits naturels des citoyens (Locke).
\end{exoresolu}

\begin{exercice}[À vous de jouer]
\begin{enumerate}
\item Définissez la légitimité et expliquez en quoi elle diffère de la force.
\item Pourquoi Hobbes juge-t-il nécessaire un souverain absolu ? Quel inconvénient cette solution présente-t-elle par rapport à celle de Locke ?
\item Distinguez volonté générale et volonté de tous chez Rousseau, puis montrez comment l'article 2 de la Constitution camerounaise de 1996 s'inspire de cette conception.
\end{enumerate}
\end{exercice}

\begin{corrige}[Exercice]
\begin{enumerate}
\item La légitimité est le caractère d'un pouvoir reconnu comme conforme au droit et à la justice par ceux qui lui obéissent. La force contraint les corps par la peur ; la légitimité obtient l'adhésion des esprits : on obéit au pouvoir légitime parce qu'on juge qu'il a le \emph{droit} de commander, et non seulement parce qu'on craint sa punition.
\item Hobbes décrit l'état de nature comme une guerre de tous contre tous : sans autorité commune, la vie est insécure. Seul un souverain tout-puissant, le Léviathan, peut imposer la paix par la crainte. Mais cette solution sacrifie entièrement la liberté : le citoyen ne dispose d'aucun recours contre un souverain injuste. Locke, au contraire, limite le pouvoir : l'État n'a pour mission que de protéger les droits naturels (vie, liberté, propriété), et le peuple garde le droit de renverser un gouvernement qui les viole.
\item La volonté générale est la volonté du corps politique visant l'intérêt commun ; la volonté de tous n'est que la somme des volontés particulières, c'est-à-dire l'addition des intérêts égoïstes. L'article 2 de la Constitution de 1996 affirme que la souveraineté nationale appartient au peuple camerounais, qui l'exerce par voie de référendum ou par ses représentants : c'est l'héritage direct de Rousseau, pour qui le peuple est le seul souverain et la loi l'expression de la volonté générale.
\end{enumerate}
\end{corrige}

Ainsi, la légitimité moderne de l'État repose sur le consentement des citoyens, exprimé par la Constitution et les élections. Mais une fois fondé, ce pouvoir peut prendre des formes très diverses : monarchie ou république, démocratie ou dictature. C'est ce que nous examinerons dans la section suivante, consacrée aux formes de gouvernements.

\coursec{Les formes de gouvernements}

\courssub{Introduction : du fondement à l'exercice du pouvoir}

Dans la section précédente, nous avons vu que l'État tire son autorité de fondements divers : la force, le contrat social, la tradition ou le droit. Mais une fois l'autorité légitimée, encore faut-il l'organiser. \textbf{Qui gouverne ?} et \textbf{Comment le pouvoir s'exerce-t-il ?} Ces deux questions commandent la classification des régimes politiques. Depuis l'Antiquité, les philosophes tentent de ranger les États en « formes de gouvernement » pour distinguer le bon du mauvais, le juste de l'injuste, le stable de l'instable. Cette classification n'est pas un exercice de bibliothèque : elle nous donne des lunettes pour lire l'actualité camerounaise et mondiale.

\courssub{La classification d'Aristote : le critère du nombre et de la finalité}

\emph{Intuition.} Imaginez trois situations : un chef de village décide seul pour tous ; un conseil d'anciens se réunit pour trancher ; l'assemblée de tous les habitants vote à main levée. Le nombre de gouvernants change, mais aussi l'esprit de la décision.

\begin{experience}[La grille de lecture d'Aristote (Politique, Livre III)]
\textbf{Observation.} Aristote croise deux critères : \textbf{le nombre de gouvernants} (un, quelques-uns, plusieurs/tous) et \textbf{la finalité du pouvoir} (l'intérêt général / commun vs l'intérêt privé / particulier des gouvernants).
\textbf{Résultat.} Cela donne six formes, rangées par paires (forme « juste » / forme « dégénérée »).
\end{experience}

\begin{popfigure}
\begin{tikzpicture}[scale=0.9, every node/.style={font=\small}]
\tikzset{
  header/.style={fill=popblue, text=white, font=\bfseries, align=center, minimum height=1.2cm},
  just/.style={fill=popgreenL, align=center, minimum height=1.5cm},
  deg/.style={fill=poppinkL, align=center, minimum height=1.5cm},
  crit/.style={fill=poporangeL, font=\bfseries, align=center, minimum height=1.2cm},
  grid/.style={draw=popline, thick}
}
\matrix (tab) [matrix of nodes, nodes={grid, anchor=center, align=center}, column sep=-\pgflinewidth, row sep=-\pgflinewidth,
  column 1/.style={crit, text width=2.5cm},
  column 2/.style={just, text width=3.5cm},
  column 3/.style={deg, text width=3.5cm},
  row 1/.style={header, minimum height=1.2cm}
] {
  \textbf{Nombre de gouvernants} & \textbf{Forme juste (Intérêt général)} & \textbf{Forme dégénérée (Intérêt privé)} \\
  \textbf{Un seul} & \textbf{Monarchie}\newline(Roi vertueux, père du peuple) & \textbf{Tyrannie}\newline(Despote, arbitraire, peur) \\
  \textbf{Quelques-uns} & \textbf{Aristocratie}\newline(Les meilleurs, les plus sages) & \textbf{Oligarchie}\newline(Riches, clan, népotisme) \\
  \textbf{Le plus grand nombre} & \textbf{République / Politéia}\newline(Citoyens vertueux, loi commune) & \textbf{Démagogie}\newline(Foule manipulée, ochlocratie) \\
};
\end{tikzpicture}
\legende{Tableau de la classification aristotélicienne des six formes de gouvernement. La « Politéia » (République) est pour Aristote la meilleure forme réaliste, mélange équilibré.}
\end{popfigure}

\begin{definition}[Formes justes et formes dégénérées]
Une \textbf{forme juste} de gouvernement est celle où le pouvoir s'exerce \emph{pour l'intérêt commun} (le bien de la cité entière). Une \textbf{forme dégénérée} (ou perverse) est celle où le pouvoir s'exerce \emph{pour l'intérêt privé} des gouvernants (un seul, une classe, une faction).
\end{definition}

\begin{attention}[Piège classique : confondre « démocratie » aristotélicienne et moderne]
Pour Aristote, la « démocratie » (pouvoir du \emph{dêmos}, peuple pauvre et majoritaire) est une forme \textbf{dégénérée} (démagogie), car le peuple gouverne pour ses intérêts de classe. La forme juste du gouvernement de plusieurs est la \textbf{Politéia} (République), où les citoyens vertueux visent le bien commun. Le sens moderne de « démocratie » (régime libre, élections, droits) correspond plutôt à la \emph{Politéia} ou à la « démocratie constitutionnelle ».
\end{attention}

\begin{exemplebox}[Illustrations historiques]
\begin{itemize}
\item \textbf{Monarchie juste :} Louis IX (Saint Louis) rendant la justice sous son chêne, cherchant l'équité pour tous.
\item \textbf{Tyrannie :} Denys l'Ancien à Syracuse, se méfiant de ses sujets et gouvernant par la terreur (l'épée de Damoclès).
\item \textbf{Aristocratie :} la République de Venise à son apogée, dirigée par un conseil de patriciens soucieux de la prospérité commune.
\item \textbf{Oligarchie :} la Rome républicaine tardive, où quelques familles patriciennes (les \emph{optimates}) accaparent les magistratures et les terres publiques.
\item \textbf{Politéia :} la Constitution américaine de 1787 (séparation des pouvoirs, fédéralisme, déclaration des droits) vise cet équilibre.
\item \textbf{Démagogie :} Athènes après la bataille des Arginuses (406 av. J.-C.), où l'assemblée, sous le coup de l'émotion, condamne à mort des généraux pourtant victorieux.
\end{itemize}
\end{exemplebox}

\courssub{La classification moderne : le critère de l'exercice du pouvoir}

\emph{Intuition.} Aujourd'hui, on regarde moins « qui » gouverne (souvent des élus) que « comment » le pouvoir s'exerce : y a-t-il des freins ? Le peuple peut-il sanctionner les gouvernants ?

\begin{definition}[Démocratie]
Forme de gouvernement où la \textbf{souveraineté appartient au peuple}, qui l'exerce \textbf{soit directement} (démocratie directe), \textbf{soit par l'intermédiaire de représentants élus} (démocratie représentative). Elle suppose : le pluralisme politique, les libertés publiques (expression, presse, association), l'État de droit, des élections libres, régulières et sincères, et la séparation des pouvoirs.
\end{definition}

\begin{definition}[Dictature]
Régime où le \textbf{pouvoir est concentré} entre les mains d'un individu ou d'un petit groupe, \textbf{sans contrôle effectif} ni reddition de comptes, souvent maintenu par la force (armée, police politique). Absence de séparation des pouvoirs, pas d'élections libres.
\end{definition}

\begin{definition}[Totalitarisme]
Régime extrême où un \textbf{parti unique} (ou un leader charismatique) contrôle \textbf{l'État, l'économie, l'information, la culture} et cherche à \textbf{encadrer toute la vie des individus} (pensée, loisirs, vie privée) par une idéologie officielle et la terreur policière (ex. : URSS stalinienne, Allemagne nazie, Cambodge des Khmers rouges).
\end{definition}

\begin{aretenir}[Différences clés]
\begin{itemize}
\item \textbf{Dictature autoritaire :} « Ne vous occupez pas de politique, obéissez, taisez-vous, nous gérons l'ordre et l'économie. » (Ex. : Pinochet au Chili, Franco en Espagne.)
\item \textbf{Totalitarisme :} « Vous \textbf{devez} adhérer au parti, penser comme nous : votre âme nous appartient. » (Ex. : Hitler, Staline.)
\item \textbf{Démocratie :} « Vous \textbf{avez le droit et le devoir} de participer, critiquer, choisir, contrôler. »
\end{itemize}
\end{aretenir}

\courssub{Les variantes de la démocratie : directe, représentative, semi-directe}

\begin{exemplebox}[Athènes, V\textsuperscript{e} siècle av. J.-C. : le modèle de la démocratie directe]
\textbf{Dispositif.} L'\textbf{Écclésia} (Assemblée du peuple) : tous les citoyens (hommes nés de parents citoyens, âgés de plus de 20 ans, soit une minorité de la population) se réunissent régulièrement sur la colline de la Pnyx. Ils votent \textbf{les lois}, \textbf{la guerre et la paix}, l'\textbf{ostracisme} (bannissement de dix ans) et élisent les stratèges (généraux).
\textbf{Exclusions.} Femmes, métèques (étrangers résidents), esclaves : \emph{non-citoyens}.
\textbf{Leçon.} La démocratie directe suppose une petite cité, une forte cohésion et du temps libre. Elle est difficilement transposable telle quelle aux États modernes de millions d'habitants.
\end{exemplebox}

\begin{definition}[Démocratie représentative (ou indirecte)]
Le peuple élit des \textbf{représentants} (députés, sénateurs, président) qui \textbf{votent les lois et gouvernent} en son nom. Le mandat est en général \textbf{représentatif} : l'élu vote librement et rend compte aux électeurs lors du renouvellement de son mandat. C'est le modèle dominant (France, États-Unis, Cameroun, Allemagne...).
\end{definition}

\begin{definition}[Démocratie semi-directe (ou participative)]
Système mixte : représentation élue + \textbf{instruments de démocratie directe} (référendum, initiative populaire, révocation des élus). Ex. : la \textbf{Suisse} (référendums fréquents), la \textbf{France} (référendum prévu par les articles 11 et 89 de la Constitution), le \textbf{Cameroun} (article 28 de la Constitution : référendum possible sur proposition du Président de la République).
\end{definition}

\begin{savaistu}[Le référendum au Cameroun]
L'article 28 de la Constitution de 1996 (révisée en 2008) prévoit que le Président de la République peut, après consultation des Bureaux de l'Assemblée nationale et du Sénat, soumettre un projet de loi au référendum. Le référendum de mai 1972 (passage à l'État unitaire) illustre cette procédure d'initiative présidentielle. Aucun référendum d'initiative populaire n'est prévu par la Constitution camerounaise.
\end{savaistu}

\courssub{La séparation des pouvoirs : le rempart contre l'arbitraire (Montesquieu)}

\emph{Problème.} Si le même organe fait les lois, les applique et les juge, il n'y a plus de loi, seulement la volonté du plus fort. \textbf{« Tout homme qui a du pouvoir est porté à en abuser ; il va jusqu'à ce qu'il trouve des limites. »} (Montesquieu, \emph{De l'esprit des lois}, 1748, XI, 4).

\begin{experience}[L'enquête de Montesquieu en Angleterre (1729-1731)]
\textbf{Observation.} Montesquieu observe un régime où le roi (exécutif) ne fait pas seul les lois (Parlement : Lords et Communes) et où les juges (judiciaire) sont indépendants. Le roi dispose d'un veto, le Parlement vote l'impôt, les juges appliquent la loi.
\textbf{Théorisation.} Trois pouvoirs \textbf{distincts} (séparation des organes) et \textbf{séparés} (personnes différentes), mais \textbf{équilibrés} par des \textbf{freins et contrepoids} (\emph{checks and balances}) : veto, dissolution, contrôle juridictionnel, responsabilité ministérielle.
\end{experience}

\begin{popfigure}
\begin{tikzpicture}[scale=1.05, every node/.style={font=\small}]
\tikzset{
  pouvoir/.style={circle, draw=popblue, thick, fill=popblueL, text width=2.6cm, align=center, minimum size=2.6cm, font=\bfseries},
  arrow/.style={->, >=Stealth, thick, popdark},
  lbl/.style={font=\scriptsize, text=popdark, align=center, text width=3cm, fill=white, inner sep=1pt}
}
\node[pouvoir] (exec) at (0,0) {EXÉCUTIF\newline \mdseries Président, Gouvernement, Administration};
\node[pouvoir] (leg) at (6,0) {LÉGISLATIF\newline \mdseries Parlement : Assemblée, Sénat};
\node[pouvoir] (jud) at (3,-4.6) {JUDICIAIRE\newline \mdseries Tribunaux, Cours, Conseil constitutionnel};

\draw[arrow, bend left=20] (exec) to node[lbl, above] {Veto, dissolution, nomination des juges} (leg);
\draw[arrow, bend left=20] (leg) to node[lbl, below] {Vote des lois et de l'impôt, contrôle du gouvernement} (exec);
\draw[arrow, bend right=15] (exec) to node[lbl, left] {Nomination des magistrats, droit de grâce} (jud);
\draw[arrow, bend right=15] (jud) to node[lbl, right] {Contrôle de constitutionnalité, indépendance du juge} (exec);
\draw[arrow, bend left=15] (leg) to node[lbl, right] {Vote du budget de la justice} (jud);
\draw[arrow, bend left=15] (jud) to node[lbl, left] {Interprétation des lois} (leg);
\end{tikzpicture}
\legende{Schéma de la séparation des pouvoirs (Montesquieu) avec les freins et contrepoids (\emph{checks and balances}). Chaque pouvoir limite les deux autres.}
\end{popfigure}

\begin{propriete}[Conditions d'une séparation effective (à savoir citer au Bac)]
\begin{enumerate}
\item \textbf{Séparation organique :} des personnes distinctes (pas de cumul strict : ministre $\neq$ député, juge $\neq$ ministre).
\item \textbf{Séparation fonctionnelle :} des fonctions distinctes (légiférer $\neq$ exécuter $\neq$ juger).
\item \textbf{Indépendance du judiciaire :} inamovibilité des magistrats du siège, Conseil Supérieur de la Magistrature (CSM) garantissant les carrières.
\item \textbf{Contrôles réciproques (\emph{checks and balances}) :} veto, dissolution, responsabilité politique, contrôle de constitutionnalité.
\end{enumerate}
\end{propriete}

\begin{attention}[Séparation formelle vs séparation réelle]
Une Constitution peut \emph{écrire} la séparation des pouvoirs (formelle) mais la \emph{vider de sa substance} par la pratique : parti unique hégémonique, nomination des juges par l'exécutif sans garde-fous, Parlement votant sans débat. \textbf{Le critère démocratique, c'est l'effectivité des contrepoids.}
\end{attention}

\courssub{Les grands régimes contemporains : organisation du pouvoir exécutif}

La séparation des pouvoirs s'incarne différemment selon le rapport entre l'exécutif et le législatif.

\begin{center}
\begin{tabularx}{\linewidth}{|l|X|X|X|}\hline
\rowcolor{popblue} \textbf{\color{white}Critère} & \textbf{\color{white}Régime présidentiel} & \textbf{\color{white}Régime parlementaire} & \textbf{\color{white}Monarchie constitutionnelle} \\ \hline
\textbf{Chef de l'État} & Président élu au suffrage universel & Président ou monarque (rôle surtout symbolique) & Roi ou reine (héréditaire, symbolique) \\ \hline
\textbf{Chef du Gouvernement} & \textbf{= Président} (fusion) & Premier ministre (issu de la majorité parlementaire) & Premier ministre (responsable devant le Parlement) \\ \hline
\textbf{Séparation} & \textbf{Rigide} : le Parlement ne peut renverser le Président (sauf procédure d'\emph{impeachment}) & \textbf{Souple / collaboration} : gouvernement responsable devant le Parlement (motion de censure), dissolution possible & Identique au régime parlementaire \\ \hline
\textbf{Exemples} & \textbf{USA, Brésil, Cameroun, Sénégal, Côte d'Ivoire} & \textbf{Allemagne, Italie, Espagne, Inde} & \textbf{Royaume-Uni, Suède, Japon, Maroc} \\ \hline
\textbf{Avantage} & Stabilité du mandat, clarté du choix populaire & Flexibilité, responsabilité continue du gouvernement & Continuité de l'État, arbitrage neutre \\ \hline
\textbf{Risque} & Blocage institutionnel, personnalisation du pouvoir & Instabilité gouvernementale, pouvoir excessif des partis & Absence de légitimité élective du monarque \\ \hline
\end{tabularx}
\end{center}

\begin{savaistu}[La « cohabitation » française : l'exception qui éclaire la règle]
La France (V\textsuperscript{e} République) est un \textbf{régime semi-présidentiel} (hybride) : un président élu au suffrage universel direct et un Premier ministre responsable devant l'Assemblée nationale. Si la majorité parlementaire est opposée au président, on parle de \textbf{cohabitation} (1986-1988, 1993-1995, 1997-2002) : le président garde less domaines « réservés » (défense, diplomatie), le Premier ministre gère la politique intérieure. Depuis le quinquennat (2000) et l'inversion du calendrier électoral (législatives juste après la présidentielle), la cohabitation est devenue improbable.
\end{savaistu}

\courssub{Application au Cameroun : République, régime présidentiel, pluralisme et enjeux}

\begin{popfigure}
\begin{tikzpicture}[scale=0.85, every node/.style={font=\small, align=center}]
\tikzset{
  inst/.style={rectangle, rounded corners, draw=popblue, thick, fill=popblueL, text width=3.4cm, minimum height=1.5cm, font=\bfseries},
  subinst/.style={rectangle, rounded corners, draw=popgreen, fill=popgreenL, text width=3.2cm, minimum height=1cm},
  link/.style={->, >=Stealth, thick, popdark},
  linkd/.style={-, dashed, popmuted},
  lbl/.style={font=\scriptsize, text=popdark, align=center, text width=2.6cm, fill=white, inner sep=1pt}
}
\node[inst] (peuple) at (0,5) {PEUPLE SOUVERAIN\newline \mdseries (suffrage universel)};

\node[inst] (pres) at (0,2.5) {PRÉSIDENT DE LA RÉPUBLIQUE\newline \mdseries Chef de l'État, chef de l'exécutif, élu pour 7 ans};
\node[subinst] (pm) at (-4.5,0.8) {PREMIER MINISTRE\newline Chef du Gouvernement, nommé par le Président};
\node[subinst] (gov) at (4.5,0.8) {GOUVERNEMENT\newline Ministres, Secrétaires d'État};

\node[inst] (parl) at (0,-1.2) {PARLEMENT (bicaméral)};
\node[subinst] (an) at (-4.5,-3) {ASSEMBLÉE NATIONALE\newline 180 députés, 5 ans, suffrage direct};
\node[subinst] (sen) at (4.5,-3) {SÉNAT\newline 100 sénateurs : 70 élus régionaux, 30 nommés par le Président};

\node[inst] (jud) at (0,-5.4) {POUVOIR JUDICIAIRE};
\node[subinst] (cc) at (-4.5,-7) {CONSEIL CONSTITUTIONNEL\newline Régularité des élections, constitutionnalité des lois};
\node[subinst] (csm) at (0,-7) {CSM\newline Gestion des carrières des magistrats};
\node[subinst] (cour) at (4.5,-7) {COUR SUPRÊME\newline Judiciaire, administratif, comptes};

\draw[link] (peuple) -- node[lbl, left] {Élit} (pres);
\draw[link] (peuple) to[bend right=10] node[lbl, left] {Élit les députés} (an);
\draw[link] (an) to node[lbl, above] {Les conseillers régionaux élisent 70 sénateurs} (sen);

\draw[link] (pres) -- node[lbl, left] {Nomination} (pm);
\draw[link] (pres) -- node[lbl, right] {Nomination sur proposition du PM} (gov);
\draw[link] (pres) to[bend left=15] node[lbl, right] {Nomination de 30 sénateurs} (sen);
\draw[link] (pres) to[bend right=15] node[lbl, left] {Nomination des membres} (cc);
\draw[link] (pres) to[bend left=25] node[lbl, right] {Nomination des magistrats (sur avis du CSM)} (cour);

\draw[linkd] (pm) -- (gov);
\draw[linkd] (parl) -- (an);
\draw[linkd] (parl) -- (sen);
\draw[linkd] (jud) -- (cc);
\draw[linkd] (jud) -- (csm);
\draw[linkd] (jud) -- (cour);

\draw[link, bend left=35, poporange] (an) to node[lbl, left] {Motion de censure, vote du budget} (gov);
\draw[link, bend left=20, poppurple] (cc) to node[lbl, above] {Contrôle de constitutionnalité des lois} (parl);
\draw[link, bend right=60, poppurple] (cc) to node[lbl, below] {Proclamation des résultats de l'élection présidentielle} (pres);
\end{tikzpicture}
\legende{Organigramme simplifié des institutions de la République du Cameroun (Constitution de 1996, révisée en 2008). Flèches pleines : nomination ou élection. Pointillés : rattachement. Flèches orange : responsabilité politique. Flèches violettes : contrôle juridictionnel.}
\end{popfigure}

\begin{definition}[Régime politique du Cameroun (art. 1\textsuperscript{er} et 5 de la Constitution)]
Le Cameroun est une \textbf{République} (chef d'État élu, non héréditaire), \textbf{décentralisée} (collectivités territoriales), \textbf{une et indivisible}, \textbf{démocratique} (pluralisme, suffrage universel), \textbf{sociale} (droits économiques et sociaux) et \textbf{laïque} (neutralité religieuse de l'État). C'est un \textbf{régime présidentiel} : le Président de la République est le \textbf{chef de l'État et le chef effectif de l'exécutif}, élu au suffrage universel direct pour sept ans (la limite du nombre de mandats, prévue en 1996, a été supprimée par la révision de 2008). Il nomme le Premier ministre et les ministres, peut dissoudre l'Assemblée nationale, nomme 30 des 100 sénateurs ainsi que les membres du Conseil constitutionnel.
\end{definition}

\begin{exoresolu}[Analyse critique : la séparation des pouvoirs au Cameroun est-elle effective ?]
\textbf{Énoncé.} À partir de l'organigramme et de la Constitution, discutez l'effectivité de la séparation des pouvoirs au Cameroun.
\tcblower
\textbf{Solution détaillée.}
\textbf{1. Séparation formelle (les textes).} La Constitution distingue clairement les trois pouvoirs. Le Président ne siège pas au Parlement. Les magistrats du siège sont en principe inamovibles. Le Conseil constitutionnel juge de la régularité des élections et de la conformité des lois à la Constitution. Il existe une Cour suprême (judiciaire, administratif, comptes). \textbf{Formellement, la séparation est acquise.}

\textbf{2. Porosité entre exécutif et législatif.}
\begin{itemize}
\item Le Président nomme le Premier ministre et le Gouvernement sans investiture parlementaire.
\item L'Assemblée nationale ne peut renverser le Gouvernement que par une \textbf{motion de censure} adoptée à une majorité qualifiée (les deux tiers des députés) : seuil très difficile à atteindre.
\item Le Président peut dissoudre l'Assemblée nationale et nomme 30\,\% des sénateurs.
\item \textbf{Résultat :} exécutif dominant, Parlement affaibli (majorité présidentielle quasi automatique depuis le retour du multipartisme). Aucune cohabitation n'est possible.
\end{itemize}

\textbf{3. Indépendance du judiciaire (le problème central).}
\begin{itemize}
\item \textbf{Nomination :} le Président nomme les magistrats \emph{sur avis du CSM}, mais il \textbf{préside le CSM}, dont le ministre de la Justice est vice-président. Le CSM n'est donc pas indépendant de l'exécutif.
\item \textbf{Conseil constitutionnel :} ses membres sont nommés par le Président de la République et par les présidents des assemblées, organes dominés par la majorité présidentielle.
\item \textbf{Carrières :} mutations, promotions et discipline dépendent du CSM, donc indirectement de l'exécutif.
\item \textbf{Conséquence :} le juge peine à sanctionner l'arbitraire administratif ou électoral (contentieux électoraux contestés à plusieurs reprises).
\end{itemize}

\textbf{4. Pluralisme politique (depuis les lois de décembre 1990 sur les libertés publiques).}
\begin{itemize}
\item \textbf{Acquis :} des centaines de partis légaux, des élections multipartites régulières, une presse privée active, des ONG de défense des droits.
\item \textbf{Limites :} code électoral critiqué (organisation d'ELECAM, découpage des circonscriptions, financement inégal, accès inégal aux médias publics), restrictions des manifestations, poursuites d'opposants.
\end{itemize}

\textbf{Conclusion.} Le Cameroun est une \textbf{République présidentielle formelle} marquée par un \textbf{hyper-présidentialisme réel}. La séparation des pouvoirs y est \textbf{formelle plus qu'effective} : l'exécutif domine le législatif (majorité mécanique, nomination de sénateurs) et influence fortement le judiciaire (nomination via le CSM et le Conseil constitutionnel). Le pluralisme existe, mais le terrain de jeu est inégal. De nombreux analystes parlent de \textbf{régime hybride} (élections pluralistes combinées à une hégémonie du parti au pouvoir), et non de démocratie consolidée au sens strict (État de droit, alternance effective, contrepoids réels).
\end{exoresolu}

\begin{methode}[Analyser un régime politique (fiche méthode Bac)]
Pour qualifier un régime (dissertation ou commentaire de texte) :
\begin{enumerate}
\item \textbf{Identifiez la source de légitimité :} le peuple (élections) ? la tradition ? la force ? la religion ?
\item \textbf{Comptez les gouvernants :} un seul (monarchie, pouvoir personnel) ? quelques-uns (oligarchie, junte) ? le plus grand nombre (démocratie) ?
\item \textbf{Examinez la finalité :} intérêt général (droits respectés, services publics) ou intérêt privé (prédation, clan, corruption) ?
\item \textbf{Testez la séparation des pouvoirs :} organique ? fonctionnelle ? indépendance du juge ? contrôles réels (veto, censure, contrôle de constitutionnalité) ?
\item \textbf{Évaluez le pluralisme :} partis libres ? médias libres ? élections sincères ? alternance possible ?
\item \textbf{Concluez :} démocratie (libérale, illibérale, déléguée) ? dictature (autoritaire, totalitaire, militaire, personnelle) ? régime hybride ?
\end{enumerate}
\end{methode}

\begin{exercice}[Entraînement : qualification de régimes]
Qualifiez les régimes suivants en utilisant la grille d'Aristote (juste/dégénéré) et la grille moderne (démocratie/dictature/totalitarisme). Justifiez par trois arguments à chaque fois.
\begin{enumerate}
\item La Suède actuelle (roi Carl XVI Gustave, élections législatives de 2022, forte transparence, faible corruption).
\item La Corée du Nord (Kim Jong-un, parti unique des Travailleurs, idéologie du \emph{Juche}, camps de rééducation, pas d'élections libres).
\item Le Chili sous Pinochet (1973-1990) : junte militaire, Constitution de 1980, répression, libéralisme économique, plébiscite de 1988 perdu par le régime.
\item La France de la III\textsuperscript{e} République (1870-1940) : Parlement puissant, président faible, ministères instables, lois laïques, affaire Dreyfus.
\end{enumerate}
\end{exercice}

\begin{corrige}[Exercice n°1 - Éléments de réponse]
\begin{enumerate}
\item \textbf{Suède :} Aristote $\to$ \textbf{Politéia} (gouvernants élus visant le bien commun, État-providence). Moderne $\to$ \textbf{démocratie parlementaire et monarchie constitutionnelle} (séparation souple, Premier ministre responsable devant le Parlement, roi purement symbolique, libertés pleines et entières ; la Suède figure parmi les premiers pays des indices internationaux de démocratie).
\item \textbf{Corée du Nord :} Aristote $\to$ \textbf{tyrannie} (un seul homme, intérêt du clan régnant, culte de la personnalité). Moderne $\to$ \textbf{totalitarisme} (parti unique, chef suprême, idéologie officielle encadrant toute la vie sociale, terreur policière, aucune séparation des pouvoirs ni liberté publique).
\item \textbf{Chili de Pinochet :} Aristote $\to$ \textbf{oligarchie militaire} (quelques généraux gouvernant pour l'« ordre » et leurs intérêts). Moderne $\to$ \textbf{dictature autoritaire militaire} (et non totalitarisme : pas de parti unique de masse ni d'idéologie encadrant toute la vie privée, économie libérale, et le régime accepte finalement le verdict du plébiscite de 1988).
\item \textbf{France de la III\textsuperscript{e} République :} Aristote $\to$ \textbf{Politéia instable} (gouvernement du plus grand nombre via le Parlement, visée de l'intérêt général mais instabilité ministérielle chronique). Moderne $\to$ \textbf{démocratie parlementaire instable} (séparation souple poussée à l'extrême, chutes fréquentes des gouvernements, mais libertés réelles, élections régulières et alternances pacifiques).
\end{enumerate}
\end{corrige}

\begin{aretenir}[Synthèse pour le Bac : les formes de gouvernements]
\begin{itemize}
\item \textbf{Aristote :} 6 formes = (nombre : un / quelques-uns / tous) $\times$ (finalité : intérêt général / intérêt privé). Monarchie/Tyrannie ; Aristocratie/Oligarchie ; Politéia/Démagogie.
\item \textbf{Classification moderne :} démocratie (souveraineté du peuple, libertés, État de droit, séparation des pouvoirs) vs dictature (pouvoir concentré, sans contrôle) vs totalitarisme (contrôle total de la société + idéologie + terreur).
\item \textbf{Types de démocratie :} directe (Athènes, partiellement la Suisse), représentative (modèle dominant), semi-directe (référendums).
\item \textbf{Séparation des pouvoirs (Montesquieu) :} exécutif / législatif / judiciaire + freins et contrepoids. Condition de la liberté politique.
\item \textbf{Régimes types :} présidentiel (USA, Cameroun), parlementaire (Allemagne), semi-présidentiel (France), monarchie constitutionnelle (Royaume-Uni, Suède).
\item \textbf{Cameroun :} République à régime présidentiel, pluralisme formel depuis 1990, mais \textbf{hyper-présidentialisme} ; séparation des pouvoirs \textbf{formelle plus qu'effective} ; régime souvent qualifié d'\textbf{hybride}.
\end{itemize}
\end{aretenir}

\coursec{Rôle et finalité de l'État}

Nous venons de voir que le pouvoir politique peut prendre des formes très diverses : monarchie, république, démocratie, régime autoritaire. Mais une question demeure, plus fondamentale encore : peu importe \emph{qui} gouverne et \emph{comment} — \emph{à quoi} l'État sert-il ? Un lycéen de Yaoundé qui passe son baccalauréat dans un lycée public, soigné à l'hôpital central, circulant sur une route goudronnée et protégé par la police, bénéficie chaque jour des services de l'État sans même y penser. Cette section se demande donc quelles sont les \cle{fonctions} de l'État et quelle est sa \cle{finalité}, c'est-à-dire le but ultime pour lequel il existe.

\courssub{Problématique : à quoi sert l'État ? Jusqu'où doit-il intervenir ?}

Imaginez un village sans chef, sans tribunal, sans école, sans police : chacun se fait justice soi-même, le plus fort l'emporte, les enfants ne sont pas instruits, les malades ne sont pas soignés. Cette image, proche de l'\emph{état de nature} décrit par Hobbes, montre par l'absurde l'utilité de l'État. Mais l'expérience inverse existe aussi : un État qui contrôle tout — ce que les citoyens lisent, disent, vendent, pensent — devient étouffant et tyrannique.

Le problème philosophique peut donc se formuler ainsi : \textbf{l'État doit-il se limiter à maintenir l'ordre, ou doit-il intervenir dans tous les domaines de la vie sociale ?} Autrement dit : où s'arrête la légitimité de son intervention ? Toute la réflexion sur le rôle et la finalité de l'État se joue dans cette tension entre un État trop faible (incapable de protéger) et un État trop fort (qui écrase les libertés).

\begin{attention}[Piège de dissertation : finalité et moyens]
Il faut distinguer soigneusement deux choses que les élèves confondent souvent :
\begin{itemize}
\item la \cle{finalité} de l'État : \emph{ce pour quoi} il existe — l'ordre, la justice, la sécurité, le bien commun ;
\item les \cle{moyens} de l'État : \emph{par quels instruments} il agit — les lois, l'impôt, la force publique (police, armée), l'administration.
\end{itemize}
Dire « la finalité de l'État est l'impôt » est une erreur : l'impôt n'est qu'un moyen au service de finalités (financer l'école, la santé, la sécurité). Dans une dissertation, commencez toujours par identifier la finalité avant de discuter les moyens.
\end{attention}

\courssub{La fonction régalienne : assurer l'ordre et la sécurité}

La fonction la plus ancienne et la plus fondamentale de l'État est la fonction dite \cle{régalienne} (du latin \emph{rex}, roi : ce qui appartient au souverain). Elle comprend trois domaines inséparables :

\begin{itemize}
\item \textbf{la police} : maintenir l'ordre intérieur, protéger les personnes et les biens ;
\item \textbf{l'armée} : défendre le territoire contre les menaces extérieures ;
\item \textbf{la justice} : trancher les conflits selon la loi et punir les délinquants.
\end{itemize}

Pourquoi cette fonction est-elle première ? Hobbes le montre dans le \emph{Léviathan} (1651) : sans un pouvoir capable de faire respecter les règles, la vie des hommes serait « solitaire, misérable, dangereuse, animale et brève ». La sécurité est la condition de tout le reste : on ne peut ni étudier, ni commercer, ni élever ses enfants dans l'insécurité. Un État qui se contente de cette seule fonction est appelé \textbf{État gendarme} : il veille, comme un gendarme au carrefour, mais n'intervient pas davantage dans la vie sociale.

\begin{exemplebox}[L'État gendarme au quotidien]
Lorsque des bandits attaquent un quartier de Douala la nuit, les habitants appellent la police ; lorsqu'un différend frontalier oppose le Cameroun à un pays voisin, c'est l'armée qui monte la garde ; lorsque deux commerçants du marché Mokolo se disputent un étal, c'est le tribunal qui tranche. Sans ces trois services, chacun devrait assurer seul sa propre défense : retour à la loi du plus fort.
\end{exemplebox}

\courssub{La fonction de justice : garantir les droits et l'égalité devant la loi}

L'État ne se contente pas de punir : il doit aussi \emph{protéger}. Sa deuxième fonction est de garantir les \cle{droits} et les \cle{libertés} des citoyens : liberté de circuler, de s'exprimer, de pratiquer sa religion, droit de propriété, droit à un procès équitable. Cela suppose que tous les citoyens soient égaux devant la loi, sans distinction d'ethnie, de religion, de fortune ou de sexe.

Cette exigence conduit à une notion capitale, celle d'État de droit.

\begin{definition}[État de droit]
Un \cle{État de droit} est un État dans lequel le pouvoir lui-même est soumis au droit : les gouvernants, comme les gouvernés, doivent respecter la Constitution et les lois. L'État de droit garantit les libertés fondamentales des citoyens et organise le contrôle de la légalité des actes du pouvoir (justice constitutionnelle, séparation des pouvoirs).
\end{definition}

L'idée est révolutionnaire : ce n'est pas seulement le citoyen qui obéit à la loi, c'est aussi — et d'abord — \emph{l'État lui-même}. Un ministre qui viole la loi doit pouvoir être jugé comme n'importe quel citoyen. Montesquieu (\emph{De l'esprit des lois}, 1748) en donne la condition : la séparation des pouvoirs législatif, exécutif et judiciaire, car « il n'y a point de liberté si la puissance de juger n'est pas séparée de la puissance législative et de l'exécutrice ».

\courssub{La fonction économique et sociale : l'État providence}

Depuis le XX\textsuperscript{e} siècle, l'État a pris en charge des tâches nouvelles qui dépassent largement le rôle de gendarme :

\begin{itemize}
\item \textbf{l'éducation} : construire des écoles, former et payer les enseignants, organiser les examens nationaux ;
\item \textbf{la santé} : hôpitaux publics, campagnes de vaccination, centres de santé intégrés ;
\item \textbf{les infrastructures} : routes, ponts, ports, barrages, stades, électrification rurale ;
\item \textbf{la redistribution} : prélever l'impôt sur tous pour financer des services accessibles à tous, y compris aux plus pauvres (bourses d'étudiants, soins gratuits pour certaines catégories).
\end{itemize}

Un tel État, qui « pourvoit » aux besoins essentiels de la population, est appelé \cle{État providence}. Son principe est la \emph{solidarité} : l'impôt payé par les plus riches finance des services dont bénéficient aussi les plus démunis, ce qui corrige les inégalités que le marché, livré à lui-même, ne corrige pas.

\begin{exemplebox}[Exemple chiffré : l'État social au Cameroun]
Au Cameroun, l'État est le premier employeur public du pays : la fonction publique compte environ \num{300000} agents (enseignants, infirmiers, policiers, administratifs). L'État finance et gère des milliers d'écoles primaires et secondaires publiques, des dizaines d'hôpitaux généraux et centraux, et des milliers de kilomètres de routes nationales. Le budget de l'État, de l'ordre de plusieurs milliers de milliards de FCFA, est alimenté essentiellement par les impôts, les taxes et les revenus pétroliers : c'est l'illustration concrète de la fonction sociale et redistributive de l'État. Le simple fait que tu passes le baccalauréat dans un lycée public, avec des enseignants payés par l'État, est une manifestation de l'État providence.
\end{exemplebox}

\begin{popfigure}
\begin{center}
\begin{tikzpicture}[scale=1.05]
% Camembert des fonctions de l'État
\fill[popblue] (0,0) -- (0:2.2) arc (0:72:2.2) -- cycle;
\fill[poporange] (0,0) -- (72:2.2) arc (72:144:2.2) -- cycle;
\fill[popgreen] (0,0) -- (144:2.2) arc (144:216:2.2) -- cycle;
\fill[poppurple] (0,0) -- (216:2.2) arc (216:288:2.2) -- cycle;
\fill[popgold] (0,0) -- (288:2.2) arc (288:360:2.2) -- cycle;
\draw[white,thick] (0,0) -- (0:2.2);
\draw[white,thick] (0,0) -- (72:2.2);
\draw[white,thick] (0,0) -- (144:2.2);
\draw[white,thick] (0,0) -- (216:2.2);
\draw[white,thick] (0,0) -- (288:2.2);
\node[white,font=\small\bfseries] at (36:1.35) {Sécurité};
\node[white,font=\small\bfseries] at (108:1.35) {Justice};
\node[white,font=\small\bfseries] at (180:1.35) {Éducation};
\node[white,font=\small\bfseries] at (252:1.35) {Santé};
\node[popdark,font=\small\bfseries] at (324:1.35) {Économie};
% Légende
\node[anchor=west,font=\small] at (2.9,1.6) {\textcolor{popblue}{\rule{0.35cm}{0.35cm}}~Ordre, police, armée};
\node[anchor=west,font=\small] at (2.9,1.0) {\textcolor{poporange}{\rule{0.35cm}{0.35cm}}~Égalité devant la loi, droits};
\node[anchor=west,font=\small] at (2.9,0.4) {\textcolor{popgreen}{\rule{0.35cm}{0.35cm}}~Écoles, formation, examens};
\node[anchor=west,font=\small] at (2.9,-0.2) {\textcolor{poppurple}{\rule{0.35cm}{0.35cm}}~Hôpitaux, protection sociale};
\node[anchor=west,font=\small] at (2.9,-0.8) {\textcolor{popgold}{\rule{0.35cm}{0.35cm}}~Infrastructures, redistribution};
\end{tikzpicture}
\end{center}
\legende{Les grandes fonctions de l'État : fonctions régaliennes (sécurité, justice) et fonctions sociales (éducation, santé, économie). Schéma indicatif : les parts ne représentent pas des données budgétaires exactes.}
\end{popfigure}

\courssub{Quelle finalité ? Les grandes doctrines en débat}

Sur la finalité de l'État, les philosophes ne sont pas d'accord. On peut ordonner leurs positions le long d'un axe, de l'État minimal à l'État maximal.

\textbf{Le libéralisme : l'État minimal.} Pour les libéraux (Locke, Adam Smith), la finalité de l'État est de protéger les \emph{libertés individuelles} et la propriété — rien de plus. L'État doit rester un « veilleur de nuit » : sécurité, justice, armée, et c'est tout. Toute intervention supplémentaire (économie, école, santé) menace la liberté et l'initiative individuelles. La devise libérale pourrait être : « l'État le meilleur est celui qui gouverne le moins ».

\textbf{Le socialisme et l'État providence : l'État intervenant.} Pour les socialistes et les partisans de l'État social, la liberté formelle ne suffit pas : à quoi sert la liberté d'étudier à qui n'a pas d'école ? La finalité de l'État est de réaliser l'\emph{égalité réelle} en intervenant dans l'économie et la société : services publics, redistribution, protection des plus faibles. L'État devient un instrument de justice sociale.

\textbf{Le bien commun : la finalité aristotélicienne.} Plus ancienne et plus profonde, la position d'Aristote (\emph{Politique}) dépasse l'alternative : « toute cité vise un bien », et la cité, communauté la plus haute, vise le bien le plus haut. L'État n'existe pas seulement pour survivre ou pour s'enrichir, mais pour permettre la \emph{vie bonne} : une vie vertueuse, juste et épanouie. C'est l'idée de \cle{bien commun}.

\begin{definition}[Bien commun]
Le \cle{bien commun} désigne ce qui profite à l'ensemble de la communauté politique, au-delà des intérêts particuliers des individus ou des groupes : la paix, la justice, l'éducation de tous, les conditions d'une vie digne. Viser le bien commun, c'est gouverner pour tous et non pour soi ou pour son clan.
\end{definition}

\begin{popfigure}
\begin{center}
\begin{tikzpicture}[scale=1]
\draw[-{Stealth[length=3mm]},thick,popdark] (-6,0) -- (6.6,0);
\node[font=\small\bfseries,popdark] at (6.6,-0.55) {Degré d'intervention de l'État};
% Graduations
\draw[thick] (-5,0.15) -- (-5,-0.15);
\draw[thick] (0,0.15) -- (0,-0.15);
\draw[thick] (5,0.15) -- (5,-0.15);
% Pôles
\node[font=\small\bfseries,popblue,align=center] at (-5,-0.75) {État minimal\\(gendarme)};
\node[font=\small\bfseries,poporange,align=center] at (0,-0.75) {État providence\\(intervenant)};
\node[font=\small\bfseries,poppurple,align=center] at (5,-0.75) {État maximal\\(totalitaire)};
% Doctrines au-dessus
\node[font=\small,popblue,align=center] at (-5,0.85) {Libéralisme\\Locke, A. Smith};
\node[font=\small,poporange,align=center] at (0,0.85) {Socialisme,\\État social};
\node[font=\small,poppurple,align=center] at (5,0.85) {Totalitarismes\\(danger !)};
% Bien commun
\node[font=\small\itshape,popgreen,align=center] at (-2.5,1.75) {Aristote : la cité vise le bien commun,};
\node[font=\small\itshape,popgreen,align=center] at (-2.5,1.35) {ni trop peu ni trop d'État};
\end{tikzpicture}
\end{center}
\legende{Axe des doctrines : de l'État minimal (gendarme libéral) à l'État maximal (totalitaire), en passant par l'État providence. La notion aristotélicienne de bien commun invite à chercher la juste mesure.}
\end{popfigure}

\courssub{Les limites de l'État et la contestation anarchiste}

L'État, même utile, comporte des dangers qu'il faut nommer clairement :

\begin{itemize}
\item \textbf{l'arbitraire} : le pouvoir, s'il n'est pas contrôlé, agit selon le bon plaisir des gouvernants et non selon la loi ;
\item \textbf{la bureaucratie} : une administration trop lourde devient lente, coûteuse, parfois corrompue, et décourage les initiatives ;
\item \textbf{la violation des libertés} : censure, surveillance, répression — l'État protecteur peut devenir oppresseur, jusqu'au régime totalitaire qui contrôle toute la vie sociale.
\end{itemize}

D'où la nécessité du \cle{contrôle citoyen} : élections, presse libre, justice indépendante, séparation des pouvoirs. L'État est un serviteur puissant qu'il faut surveiller en permanence.

Certains penseurs vont plus loin et contestent l'État lui-même : c'est l'\cle{anarchisme} (Proudhon, Bakounine, Kropotkine). Pour les anarchistes, l'État n'est pas un remède mais le mal lui-même : toute autorité imposée est une domination, et les hommes pourraient s'organiser librement, par association volontaire et entraide, sans État. La question anarchiste est redoutable : \textbf{l'État est-il un mal nécessaire, ou un mal tout court ?} La plupart des philosophes répondent qu'il est un mal \emph{nécessaire} : on ne peut s'en passer (sans lui, la loi du plus fort), mais il faut le limiter et le contrôler pour qu'il reste au service des citoyens.

\begin{savaistu}[La Boétie : le pouvoir tient par notre consentement]
À seulement 16 ans, Étienne de La Boétie écrivit le \emph{Discours de la servitude volontaire} (1549), l'un des textes les plus étonnants de la philosophie politique. Il y pose une question simple : comment un seul homme peut-il commander à des milliers ? Sa réponse : le pouvoir ne tient que par le consentement — même passif — des gouvernés. « Soyez résolus de ne servir plus, et vous voilà libres. » Ce texte, lu et relu pendant des siècles, montre que la force de l'État repose en réalité sur l'obéissance que nous lui accordons : une idée qui fonde toutes les réflexions modernes sur la désobéissance civile et le contrôle citoyen du pouvoir.
\end{savaistu}

\courssub{Application au Cameroun et méthode de conclusion}

Au Cameroun, le débat sur le rôle de l'État est très concret. L'État assume les fonctions régaliennes (police, gendarmerie, armée, tribunaux) et les fonctions sociales : lycées publics et examens nationaux (BEPC, Probatoire, Baccalauréat), hôpitaux publics, grands travaux (routes, barrages, stades de la CAN). Mais un débat important anime la vie politique : celui de la \cle{décentralisation}. Faut-il transférer davantage de compétences (écoles, santé, routes locales) aux régions et aux communes, plus proches des citoyens ? Les partisans de la décentralisation y voient une administration plus efficace et plus démocratique ; ses critiques craignent une gestion locale inégale ou clientéliste. Ce débat illustre parfaitement la question philosophique de la section : \emph{jusqu'où}, et \emph{à quel niveau}, l'État doit-il intervenir ?

\begin{exoresolu}[Sujet de dissertation : « L'État doit-il tout faire pour les citoyens ? »]
\textbf{Énoncé.} Un élève affirme : « Puisque l'État nous protège, il doit aussi nous loger, nous nourrir et nous employer : l'État doit tout faire. » Discutez cette opinion en mobilisant les notions de la leçon.
\tcblower
\textbf{Solution détaillée (plan et arguments).}
\begin{enumerate}
\item \textbf{Analyse du sujet.} « Tout faire » signifie une intervention totale de l'État. La tension : entre la fonction protectrice légitime et le risque d'étouffement des libertés.
\item \textbf{Thèse — l'État a des fonctions indispensables.} Sans État, pas d'ordre ni de justice (Hobbes) ; l'État providence garantit l'égalité réelle : écoles publiques, hôpitaux, redistribution par l'impôt. Exemple : au Cameroun, le baccalauréat dans les lycées publics est accessible grâce à l'État.
\item \textbf{Antithèse — mais l'État qui fait tout devient dangereux.} Bureaucratie, arbitraire, violation des libertés ; l'État maximal glisse vers le totalitarisme. Le libéralisme rappelle que l'initiative individuelle et la liberté sont aussi des biens. L'anarchisme montre, par l'extrême, que toute puissance étatique appelle la méfiance.
\item \textbf{Synthèse.} L'État ne doit pas tout faire, mais faire \emph{l'essentiel} : sécurité, justice, services publics fondamentaux — dans le cadre d'un État de droit limité et contrôlé par les citoyens, au service du bien commun (Aristote).
\end{enumerate}
\end{exoresolu}

\begin{methode}[Conclure une dissertation sur l'État]
Une bonne conclusion sur ce thème suit quatre étapes :
\begin{enumerate}
\item \textbf{Rappeler la tension} qui structure le sujet : l'État est à la fois une puissance nécessaire (ordre, justice, protection) et un danger permanent (arbitraire, oppression).
\item \textbf{Reformuler la réponse} apportée au problème posé dans l'introduction, en une phrase claire.
\item \textbf{Justifier une position personnelle argumentée} : par exemple, défendre l'État de droit — un État assez fort pour protéger, mais limité par le droit et contrôlé par les citoyens.
\item \textbf{Ouvrir} sur une question voisine : quels moyens concrets les citoyens ont-ils de contrôler le pouvoir (élections, presse, justice) ?
\end{enumerate}
\end{methode}

\begin{exercice}[Atelier d'application : analyse de situation concrète]
\textbf{Situation.} La commune de votre arrondissement doit choisir entre deux projets d'investissement pour l'année prochaine, le budget ne permettant d'en financer qu'un seul :
\begin{itemize}
\item \textbf{Projet A :} Construction d'un nouveau commissariat de police dans un quartier en forte expansion où l'insécurité augmente.
\item \textbf{Projet B :} Construction d'un centre de santé intégré dans un quartier populaire mal desservi.
\end{itemize}
\textbf{Consignes.}
\begin{enumerate}
\item Identifiez, pour chaque projet, la \cle{fonction de l'État} concernée (régalienne, sociale, etc.).
\item En vous référant à la notion de \cle{bien commun} et aux doctrines étudiées (libéralisme, État providence), arguez en faveur du Projet A, puis du Projet B.
\item Proposez une \textbf{décision motivée} en tant que conseiller municipal, en précisant quel critère philosophique (sécurité première, égalité d'accès aux soins, bien commun) guide votre choix.
\end{enumerate}
\end{exercice}

\begin{corrige}[Exercice n°1 — Atelier d'application]
\begin{enumerate}
\item \textbf{Projet A (Commissariat) :} Fonction \emph{régalienne} (police, sécurité intérieure, protection des personnes et des biens). \textbf{Projet B (Centre de santé) :} Fonction \emph{sociale / État providence} (santé publique, accès aux soins, solidarité, réduction des inégalités).
\item \textbf{Argumentaire Projet A (inspiration libérale / hobbesienne) :} La sécurité est la condition \emph{sine qua non} de toute vie sociale (Hobbes). Sans ordre, le centre de santé ne pourrait pas fonctionner (vols, violences). L'État minimal a pour finalité première la protection. Investir dans la police, c'est garantir le cadre même de l'exercice des autres droits.
\item \textbf{Argumentaire Projet B (inspiration État providence / bien commun) :} La santé est un droit fondamental et une condition de la \emph{liberté réelle} (à quoi sert la sécurité si l'on meurt de soigner ?). L'État providence vise l'égalité réelle : ce quartier populaire est victime d'une inégalité territoriale. Le bien commun (Aristote) inclut la préservation de la vie et de la santé de la population la plus vulnérable.
\item \textbf{Décision motivée (exemple) :} Je choisis le Projet B. \emph{Justification :} Si l'insécurité est réelle, l'absence de soins tue silencieusement et massivement (paludisme, accouchements, malnutrition). Le bien commun commande de protéger la vie là où elle est le plus menacée par la négligence structurelle. De plus, un centre de santé crée du lien social et de l'emploi local. La sécurité pourra être renforcée par une coopération avec la gendarmerie existante en attendant le budget suivant.
\end{enumerate}
\end{corrige}

\begin{aretenir}[L'essentiel de la section]
\begin{itemize}
\item Les fonctions de l'État : régaliennes (police, armée, justice), de justice (droits, égalité devant la loi), économiques et sociales (éducation, santé, infrastructures, redistribution).
\item L'\cle{État de droit} : le pouvoir lui-même est soumis au droit ; le \cle{bien commun} : ce qui profite à toute la communauté.
\item Doctrines : libéralisme (État minimal), État providence (intervention pour l'égalité), Aristote (la cité vise la vie bonne).
\item L'État est un \emph{mal nécessaire} : indispensable mais dangereux — il doit être limité et contrôlé par les citoyens.
\end{itemize}
\end{aretenir}

\coursec{Atelier d'application : analyser des situations concrètes}

Après avoir examiné le \emph{rôle et la finalité de l'État} — sa mission de service public, sa responsabilité dans la justice, la sécurité et le développement —, nous devons maintenant passer de la théorie à la pratique. La philosophie n'est pas une simple culture générale : elle est une \cle{boîte à outils} pour penser le réel. Dans cette dernière section, nous allons apprendre à \textbf{diagnostiquer une situation politique concrète}, à \textbf{repérer les concepts en jeu} et à \textbf{construire une argumentation rigoureuse}, exactement comme l'exigent les épreuves du Probatoire et du Baccalauréat technique.

\courssub{Méthode d'analyse d'une situation politique : les 4 étapes}

Toute situation politique (un fait de société, un conflit d'autorité, un débat public) peut être disséquée selon une démarche en quatre temps. Cette méthode est votre \textbf{grille de lecture universelle} pour l'examen comme pour la vie citoyenne.

\begin{methode}[Analyser une situation politique en 4 étapes]
\begin{enumerate}
    \item \textbf{Les faits bruts (Description neutre)} : Que se passe-t-il ? Qui sont les acteurs ? Quelles sont leurs actions ou leurs paroles ? Pas d'interprétation, pas de jugement : seulement l'observation (comme un journaliste ou un ethnologue).
    \item \textbf{Les notions mobilisées (Qualification philosophique)} : Quels concepts du cours permettent de nommer ce qui se joue ? (Ex. : \cle{souveraineté}, \cle{légitimité}, \cle{contrat social}, \cle{séparation des pouvoirs}, \cle{devoir civique}, \cle{autorité traditionnelle} vs \cle{autorité étatique}).
    \item \textbf{Le problème philosophique (Problématisation)} : Quelle tension, quelle contradiction, quelle question fondamentale ce fait révèle-t-il ? Formuler sous forme de question ouverte opposant deux thèses plausibles. (Ex. : « L'obéissance à l'État est-elle due à la peur ou au consentement ? » ; « La tradition peut-elle coexister avec la loi républicaine ? »).
    \item \textbf{La thèse justifiée (Argumentation)} : Prendre position, étayer par des arguments logiques, des exemples concrets (dont le contexte camerounais) et des \textbf{références philosophiques précises} (Hobbes, Locke, Rousseau, Montesquieu, Kelsen, auteurs africains), puis répondre à l'objection adverse.
\end{enumerate}
\end{methode}

\begin{popfigure}
\begin{tikzpicture}[
    node distance=1.2cm and 2.5cm,
    box/.style={rectangle, rounded corners, draw=popblue, thick, fill=popblueL, text width=3.2cm, align=center, font=\small\bfseries, minimum height=1.8cm},
    arrow/.style={-Stealth, thick, popdark},
    labelbox/.style={rectangle, draw=poporange, thick, fill=poporangeL, text width=4.8cm, align=center, font=\small, minimum height=1.5cm}
]
\node[box, align=center] (etape1){1. FAITS\\Description neutre\\Acteurs, actions, paroles};
\node[box, right=of etape1, align=center] (etape2){2. NOTIONS\\Qualification\\Souveraineté, légitimité,\\contrat social, devoir...};
\node[box, right=of etape2, align=center] (etape3){3. PROBLÈME\\Problématisation\\Question opposant 2 thèses};
\node[box, right=of etape3, align=center] (etape4){4. THÈSE\\Argumentation\\Position + arguments +\\auteurs + objection};
\draw[arrow] (etape1.east) -- (etape2.west) node[midway, above, font=\tiny, popmuted, align=center] {Quels concepts\\éclairent les faits ?};
\draw[arrow] (etape2.east) -- (etape3.west) node[midway, above, font=\tiny, popmuted, align=center] {Quelle tension\\philosophique ?};
\draw[arrow] (etape3.east) -- (etape4.west) node[midway, above, font=\tiny, popmuted, align=center] {Quelle réponse\\argumentée ?};
% Boîtes d'explication sous chaque étape
\node[labelbox, below=0.8cm of etape1] (l1) {\textbf{Exemple :} « Un élève refuse que ses parents paient l'impôt car « l'État ne fait rien ». Acteur : l'élève (citoyen futur). Action : refus fiscal. Parole : critique de l'efficacité étatique.};
\node[labelbox, below=0.8cm of etape2] (l2) {\textbf{Notions :} Contrat social (Rousseau), devoir civique, contrepartie État-citoyen, consentement à l'impôt (Locke), légitimité de l'autorité.};
\node[labelbox, below=0.8cm of etape3] (l3) {\textbf{Problème :} L'obéissance à l'État est-elle conditionnelle à sa performance ? Le contrat social est-il un contrat commercial ?};
\node[labelbox, below=0.8cm of etape4] (l4) {\textbf{Thèse :} Non, le devoir fiscal est inconditionnel (solidarité nationale), mais le citoyen a le droit de contrôler l'usage de l'impôt (vote, presse, pétition). Réf. : Rousseau, Art. 14 Déclaration 1789.};
\end{tikzpicture}
\legende{Schéma de la méthode d'analyse en 4 étapes : du fait brut à la thèse philosophique justifiée.}
\end{popfigure}

\courssub{Situation 1 : Le refus de l'impôt — Contrat social et devoirs du citoyen}

\begin{exemplebox}[Situation concrète]
\textbf{Énoncé :} Lors d'une séance d'éducation à la citoyenneté, un élève de Terminale déclare : « \emph{Pourquoi mes parents paieraient des impôts ? L'État ne fait rien : pas d'eau courante, routes défoncées, hôpitaux sans médicaments. Tant qu'il ne remplit pas sa part du contrat, nous ne payons pas.} »
\end{exemplebox}

\begin{exoresolu}[Analyse guidée de la situation 1]
\textbf{Étape 1 — Faits :} Un élève (futur citoyen) conteste la légitimité de l'impôt au nom de l'inefficacité perçue des services publics. Il assimile la relation État-citoyen à un \emph{contrat commercial} (do ut des : je donne si tu donnes).

\textbf{Étape 2 — Notions :} \cle{Contrat social} (Hobbes, Locke, Rousseau), \cle{devoir civique}, \cle{consentement à l'impôt}, \cle{solidarité nationale}, \cle{légitimité de l'État}, \cle{droit de résistance} (lockeen) vs \cle{obéissance inconditionnelle} (hobbésienne).

\textbf{Étape 3 — Problème :} \textbf{L'obéissance fiscale est-elle conditionnelle à la performance de l'État ?} Autrement dit : le contrat social est-il un contrat synallagmatique (échange de prestations) ou un pacte de constitution de la société politique ?

\textbf{Étape 4 — Thèse justifiée :}
\begin{itemize}
    \item \textbf{Thèse :} Le devoir fiscal n'est \emph{pas} conditionnel à la satisfaction immédiate du citoyen. L'impôt est la \cle{contrepartie obligatoire} de l'appartenance à la communauté politique (solidarité), non le paiement d'un service à la carte.
    \item \textbf{Argument 1 (Rousseau, \emph{Du contrat social}, I, 6-7)} : Le pacte social crée le \cle{corps politique} ; chaque associé met sa personne et sa puissance sous la direction de la volonté générale. L'impôt n'est pas un prix, c'est la \emph{contribution à la chose commune}. « Afin que le pacte social ne soit pas un vain formulaire, il renferme tacitement cet engagement : que quiconque refusera d'obéir à la volonté générale y sera contraint par le corps tout entier. »
    \item \textbf{Argument 2 (Locke, \emph{Second traité du gouvernement}, §140)} : Le consentement à l'impôt se donne par l'intermédiaire des représentants (Assemblée nationale). Au Cameroun, l'Assemblée vote la loi de finances. Refuser l'impôt, c'est refuser la représentation nationale.
    \item \textbf{Argument 3 (Réalité camerounaise)} : L'État camerounais, malgré ses défaillances, assure des fonctions régaliennes (sécurité, administration, éducation, santé de base, infrastructures). La \cle{solidarité nationale} impose que les régions riches (Littoral, Centre) contribuent au développement des régions enclavées (Extrême-Nord, Est, Adamaoua).
    \item \textbf{Objection (anticipée) :} « Mais si l'État est corrompu ou défaillant, payer c'est cautionner ! » \textbf{Réponse :} La corruption ne supprime pas le devoir fiscal ; elle appelle le \cle{contrôle citoyen} (vote, presse libre, pétitions, Cour des comptes, Conseil constitutionnel — Art. 65 Constitution). La désobéissance civile (Thoreau, Gandhi, Mandela) est un acte \emph{public, non-violent et assumé}, non une fraude fiscale égoïste.
\end{itemize}
\end{exoresolu}

\begin{attention}[Piège fréquent au Bac : Confondre contrat social et contrat commercial]
Ne dites jamais : « L'État ne tient pas sa part du marché, donc je ne paie pas. » Le contrat social \textbf{n'est pas} un contrat de droit privé (Code civil). Il fonde la \textbf{légitimité de l'ordre politique}, pas un échange de services. La sanction du mauvais gouvernement, en démocratie, c'est \textbf{l'alternance par le vote}, pas la grève de l'impôt (sauf cas extrême de tyrannie — droit de résistance lockeen).
\end{attention}

\courssub{Situation 2 : Chef traditionnel vs Maire — Légitimité traditionnelle et légitimité étatique}

\begin{exemplebox}[Situation concrète]
\textbf{Énoncé :} Dans un village de la région de l'Ouest, le \emph{chef de troisième degré} (chef traditionnel) et le \emph{maire} de la commune d'arrondissement revendiquent tous deux l'autorité pour attribuer un terrain communal à un investisseur. Le chef invoque la « coutume et la tradition ancestrale » ; le maire invoque « la loi n° 2004/017 du 22 juillet 2004 d'orientation de la décentralisation » et le Code foncier.
\end{exemplebox}

\begin{methode}[Analyser un conflit de légitimité : Tableau comparatif]
Complétez ce tableau pour structurer votre raisonnement. C'est un exercice type de l'épreuve écrite (commentaire ou dissertation).
\end{methode}

\begin{popfigure}
\begin{center}
\begin{tabularx}{\linewidth}{|l|X|X|}
\hline
\rowcolor{popblueL}
\textbf{Critère de comparaison} & \textbf{Légitimité traditionnelle (Chef)} & \textbf{Légitimité étatique (Maire)} \\
\hline
\textbf{Fondement (Source)} & Coutume, lignage, ancêtres, investiture rituelle, reconnaissance communautaire & Constitution (Art. 1, 2, 55), Loi, Suffrage universel, Investiture administrative (Préfet) \\
\hline
\textbf{Mode de désignation} & Héréditaire / Élection par notables / Désignation par conseil de famille / Rite initiatique & Élection au suffrage universel direct (conseillers municipaux) puis élection du maire par le conseil municipal \\
\hline
\textbf{Domaine d'autorité (Compétences)} & Gestion des terres coutumières, règlement des conflits coutumiers, préservation rites/identité, mobilisation communautaire & Gestion domaine public/privé communal, urbanisme, état civil, développement local, police administrative (marchés, voirie), exécution budget \\
\hline
\textbf{Responsabilité / Contrôle} & Moral, symbolique, pression sociale, destitution par notables/famille (rare) & Juridique (tutelle administrative, Cour des comptes, Tribunal administratif), Politique (réélection, motion de censure), Pénale (détournement) \\
\hline
\textbf{Reconnaissance par l'État} & Loi n° 77/245 (chefs traditionnels) : auxiliaires de l'administration, rôle consultatif, pas de pouvoir réglementaire & Pleine légalité constitutionnelle, pouvoir réglementaire (arrêtés municipaux), pouvoir de police \\
\hline
\end{tabularx}
\end{center}
\legende{Tableau comparatif : légitimité traditionnelle vs légitimité étatique. À compléter et à apprendre par cœur pour le Bac.}
\end{popfigure}

\begin{exoresolu}[Analyse guidée de la situation 2]
\textbf{Problème :} \textbf{Comment articuler l'autorité coutumière et l'autorité républicaine dans un État de droit décentralisé ?}

\textbf{Thèse justifiée :} L'État de droit (Art. 1er Constitution : « Le Cameroun est une République... ») impose la \cle{primauté de la loi} sur la coutume. Le maire, élu au suffrage universel, détient la \cle{légalité} et la \cle{légitimité démocratique} pour gérer le domaine communal (Art. 27 Loi décentralisation). Le chef traditionnel conserve une \cle{légitimité symbolique et sociale} indispensable à l'acceptabilité locale (médiation, cohésion), mais son pouvoir foncier est \emph{subordonné} à la loi (Code foncier : domaine national = propriété de l'État). \textbf{Solution concrète :} La concertation obligatoire (Art. 30 Loi décentralisation : « Le maire consulte les autorités traditionnelles »). L'investisseur doit obtenir l'avis du chef (acceptabilité sociale) \textbf{ET} l'acte juridique du maire (sécurité juridique). C'est l'articulation \emph{coutume + loi} = \cle{gouvernance locale apaisée}.
\end{exoresolu}

\begin{savaistu}[Le Cameroun et le pluralisme juridique]
L'article 372 du Code pénal camerounais punit « quiconque, abusant de l'autorité que lui confère sa qualité de chef traditionnel... ». La Constitution de 1996 (Préambule) affirme l'attachement aux « valeurs traditionnelles » \textbf{compatibles avec les droits de l'homme et la démocratie} ». Le Cameroun est un \textbf{État de droit pluraliste} : la coutume est source de droit \emph{subsidiaire} (Art. 1er Code civil : « La coutume ne s'applique qu'en l'absence de disposition légale »). Le chef n'est pas un « mini-État » ; il est un \emph{partenaire institutionnel} de la commune.
\end{savaistu}

\courssub{Situation 3 : « Faut-il un homme fort ? » — Hobbes vs Démocratie}

\begin{exemplebox}[Situation concrète]
\textbf{Énoncé :} Sur les réseaux sociaux (Facebook, WhatsApp, X), un débat vif oppose deux camps : « \emph{Il nous faut un homme fort, un père de la nation qui décide vite et bien, sans perdre de temps en palabres parlementaires} » (Camp A) vs « \emph{Un homme fort, c'est un dictateur. Seule la démocratie, avec ses contre-pouvoirs, garantit nos libertés} » (Camp B).
\end{exemplebox}

\begin{exoresolu}[Analyse guidée de la situation 3 — Confrontation Hobbes / Démocratie moderne]
\textbf{Étape 1-2 : Faits \& Notions.} Le débat oppose \cle{efficacité décisionnelle} (Camp A) et \cle{garantie des libertés / séparation des pouvoirs} (Camp B). Notions : \cle{Souveraineté}, \cle{Léviathan} (Hobbes), \cle{Séparation des pouvoirs} (Montesquieu), \cle{Démocratie}, \cle{État de droit}, \cle{Populisme}, \cle{Autoritarisme}.

\textbf{Étape 3 — Problème :} \textbf{La concentration du pouvoir (homme fort) est-elle la condition de l'efficacité étatique ou la négation de la liberté politique ?}

\textbf{Étape 4 — Thèse justifiée (Dialectique) :}
\begin{itemize}
    \item \textbf{Thèse 1 (L'argument de l'efficacité — Hobbes, \emph{Léviathan}, 1651)} : Dans l'état de nature, c'est la « guerre de tous contre tous » (\emph{bellum omnium contra omnes}). La peur de la mort violente pousse les hommes à aliéner \emph{toute} leur liberté à un \cle{Souverain absolu} (Léviathan) qui assure la paix et la sécurité. \textbf{Argument :} L'homme fort incarne l'unité et la décision rapide. Au Cameroun, certains citent la stabilité relative comme preuve de la vertu d'un pouvoir fort.
    \item \textbf{Antithèse (Le risque de tyrannie — Montesquieu, \emph{De l'esprit des lois}, XI, 4)} : « Tout homme qui a du pouvoir est porté à en abuser ; il va jusqu'à ce qu'il trouve des limites. » \textbf{Séparation des pouvoirs} (exécutif, législatif, judiciaire) + \cle{contrôles et équilibres} (checks and balances) = garantie contre l'arbitraire. L'histoire camerounaise (années 90, crises anglophones) montre que l'absence de contre-pouvoirs réels engendre la frustration et le conflit.
    \item \textbf{Synthèse (Démocratie comme efficacité \emph{durable}) :} La démocratie n'est pas la « palabre » stérile ; c'est la \cle{procéduralisation du conflit} (Habermas). Un pouvoir fort \emph{contrôlé} (parlement, justice indépendante, presse, société civile) est plus efficace \emph{sur le long terme} car il intègre la diversité, prévient l'explosion sociale et assure la \cle{légitimité} (consentement). L'article 1er de la Constitution camerounaise : « La souveraineté nationale appartient au peuple » — pas à un homme.
    \item \textbf{Référence africaine :} \textbf{Achille Mbembe} (\emph{Critique de la raison nègre}, \emph{De la postcolonie}) : La « big-man » politics (politique du grand homme) est une pathologie du pouvoir postcolonial, pas une solution. \textbf{Felwine Sarr} (\emph{Afrotopia}) : Il faut inventer des « démocraties africaines » fondées sur la \cle{reddition de comptes} (accountability) et la \cle{participation}, pas sur le sauveur providentiel.
\end{itemize}
\end{exoresolu}

\begin{attention}[Erreur de rédaction : « Hobbes est pour la dictature »]
Hobbes justifie l'absolutisme \textbf{par la raison} (sortir de l'état de nature), pas par le goût du pouvoir. Il pose la \textbf{condition} : le Souverain doit assurer la \emph{sécurité} (salus populi suprema lex). S'il faillit, le pacte est rompu (droit de résistance implicite). Ne caricaturez jamais les auteurs.
\end{attention}

\courssub{Analyse guidée : Extraits de la Constitution camerounaise de 1996 (Articles 1 et 2)}

\begin{exemplebox}[Texte officiel — Constitution de la République du Cameroun (Loi n° 96-06 du 18 janvier 1996)]
\textbf{Article 1er :} « Le Cameroun est une République décentralisée. Elle est une et indivisible, laïque, démocratique et sociale. Elle assure l'égalité devant la loi de tous les citoyens sans distinction de sexe, de race, de religion, de langue, d'opinion politique, de croyance philosophique. Son organisation est présidentielle. Le peuple exerce sa souveraineté par l'intermédiaire de ses représentants élus ou par voie de référendum. Aucune section du peuple, aucune corporation, aucun parti, aucune tribu, aucun individu ne peut s'en attribuer l'exercice. »

\textbf{Article 2 :} « Le drapeau est le drapeau tricolore vert, rouge, jaune, chargé d'une étoile d'or au centre du voile rouge. L'hymne national est le « Chant de ralliement ». La devise de la République est : « Paix — Travail — Patrie ». Les langues officielles sont le français et l'anglais. La capitale de la République est Yaoundé. »
\end{exemplebox}

\begin{methode}[Comment analyser un article constitutionnel (Méthode commentaire de texte Bac)]
\begin{enumerate}
    \item \textbf{Contextualiser :} Date, auteur (Constituant), place dans la hiérarchie des normes (sommet de la pyramide kelsénienne).
    \item \textbf{Identifier les concepts-clés} (surligner) : République, décentralisée, une/indivisible, laïque, démocratique, sociale, souveraineté, représentation, référendum, indivisibilité de la souveraineté.
    \item \textbf{Expliquer chaque concept} en le reliant au cours (définition + auteur de référence).
    \item \textbf{Montrer l'articulation} : comment ces principes s'organisent-ils ? (Ex. : Décentralisation + Indivisibilité = tension gérée par la tutelle administrative).
    \item \textbf{Évaluer la portée} : Ce texte est-il une réalité effective ou un idéal normatif ? (Écart texte/réalité).
\end{enumerate}
\end{methode}

\begin{exoresolu}[Analyse des Articles 1 et 2]
\begin{itemize}
    \item \textbf{\cle{République} (Art. 1) :} \emph{Res publica} = « chose publique ». Le pouvoir n'est pas une propriété privée (monarchie) mais une charge publique, temporaire, responsable. \textbf{Réf.} : Cicéron, Montesquieu, Rousseau.
    \item \textbf{\cle{Décentralisée} :} Organisation territoriale avec des collectivités (communes, régions) dotées de personnalité juridique et d'autonomie financière (Art. 55). \textbf{Distinction :} Décentralisation (autonomie) $\neq$ Déconcentration (services déconcentrés de l'État : Préfet, Sous-préfet).
    \item \textbf{\cle{Une et indivisible} :} Principes d'unité nationale. Interdit la sécession, le fédéralisme (souvenir du fédéralisme 1961-1972), le tribalisme politique. Tension avec « décentralisée » : l'autonomie locale ne doit pas fracturer l'unité.
    \item \textbf{\cle{Laïque} :} Neutralité de l'État en matière religieuse. Séparation Églises/État. Liberté de conscience (Art. 1er al. 2). Pas de religion d'État.
    \item \textbf{\cle{Démocratique} :} \cle{Souveraineté nationale} = peuple (Art. 1 al. 4). Exercice : \textbf{représentation} (députés, conseillers) + \textbf{référendum} (démocratie semi-directe). \textbf{Principe clé :} « Nul ne peut s'en attribuer l'exercice » = \cle{interdiction de la confiscation du pouvoir} (coup d'État, présidence à vie, monopole partisan).
    \item \textbf{\cle{Sociale} :} L'État a une obligation de \cle{justice sociale} (santé, éducation, protection des plus faibles). Réf. : Préambule 1946, Droit du travail, Sécurité sociale (CNPS).
    \item \textbf{Symboles (Art. 2) :} Drapeau, hymne, devise, langues, capitale = \cle{marqueurs de l'identité nationale} et de la \cle{souveraineté} (Weber : monopole de la violence légitime + symboles).
\end{itemize}
\end{exoresolu}

\courssub{Méthode de la dissertation : Construire une problématique (Sujets type Bac)}

La dissertation philosophique n'est pas une récitation du cours. C'est une \textbf{réponse argumentée à une question précise}. La \cle{problématisation} est l'acte fondateur : elle transforme un sujet plat en une tension intellectuelle.

\begin{methode}[Problématiser un sujet en 3 temps (Modèle Bac)]
\textbf{Sujet type 1 : « L'État est-il nécessaire ? »}
\begin{enumerate}
    \item \textbf{Analyser les termes :} « État » (institution politique souveraine sur un territoire), « nécessaire » (qui ne peut pas ne pas être ; condition de possibilité ; indispensable).
    \item \textbf{Chercher l'opposition (Thèse vs Antithèse) :}
        \begin{itemize}
            \item \textbf{Thèse (OUI) :} L'État est \emph{nécessaire} pour sortir de l'état de nature (Hobbes), garantir les droits (Locke), réaliser la liberté civile (Rousseau), assurer la justice, la paix, les services publics. Sans État = chaos, insécurité, loi du plus fort.
            \item \textbf{Antithèse (NON / PAS FORCÉMENT) :} L'État peut être \emph{nuisible} (aliénation, oppression, bureaucratie, totalitarisme). La société civile peut s'autoréguler (anarchisme : Proudhon, Kropotkine ; société sans État : Clastres, anthropologie des sociétés « contre l'État »). L'État n'est pas \emph{ontologiquement} nécessaire, il est \emph{historiquement contingent}.
        \end{itemize}
    \item \textbf{Formuler la problématique (Question ouverte) :} \textbf{« L'État est-il la condition indispensable de la vie commune ou une menace pour la liberté qu'il prétend protéger ? »} ou \textbf{« La nécessité de l'État ne cache-t-elle pas le risque de sa toute-puissance ? »}
    \item \textbf{Annonce du plan (3 parties) :}
        \begin{enumerate}
            \item L'État comme condition de possibilité de l'ordre et de la liberté (Hobbes, Locke, Rousseau, Kelsen).
            \item Les dangers de l'État : pouvoir absolu, aliénation, bureaucratie (Marx, Weber, Arendt, anarchistes).
            \item L'État de droit démocratique : l'institutionnalisation de la nécessité \emph{et} de la limite (Montesquieu, Kelsen, Constitution 1996, Habermas).
        \end{enumerate}
\end{enumerate}

\textbf{Sujet type 2 : « L'État doit-il tout contrôler ? »}
\begin{enumerate}
    \item \textbf{Termes :} « Tout contrôler » = totalitarisme (contrôle total : économie, pensée, vie privée, mœurs) vs État régulateur (cadre général, services publics, justice).
    \item \textbf{Opposition :}
        \item Thèse : L'État doit tout contrôler pour le bien commun (sécurité, santé, éducation, économie planifiée, moralité) — \emph{État providence / État totalitaire / Léviathan}.
        \item Antithèse : L'État doit se limiter aux fonctions régaliennes (souveraineté, justice, police, défense) + services publics essentiels. La société (marché, associations, famille, religion) doit rester libre — \emph{État minimal (Nozick), État de droit, subsidiarité}.
    \item \textbf{Problématique :} \textbf{« La prétention de l'État à tout contrôler n'est-elle pas la négation de la société civile et de la liberté individuelle ? »}
    \item \textbf{Plan :}
        \begin{enumerate}
            \item Les arguments du contrôle total : efficacité, égalité, sécurité (Platon, Hobbes, Rousseau \emph{législateur}, totalitarismes XXe).
            \item La nécessité de la limite : séparation pouvoir/société, droits de l'homme, économie de marché, vie privée (Locke, Constant, Tocqueville, Hayek, Nozick).
            \item L'État \emph{subsidiaire} et \emph{partenaire} : contrôle des règles du jeu (justice, régulation) + services publics, mais autonomie de la société (décentralisation, société civile, Art. 1er \& 55 Constitution).
        \end{enumerate}
\end{enumerate}
\end{methode}

\begin{attention}[Erreur fatale au Bac : Réciter le cours sans répondre au sujet]
\begin{itemize}
    \item \textbf{Mauvais :} « L'État est nécessaire. Hobbes dit que sans État c'est la guerre. Locke dit qu'il protège la propriété. Rousseau dit qu'il exprime la volonté générale. Montesquieu dit qu'il faut séparer les pouvoirs. Conclusion : l'État est nécessaire. » $\rightarrow$ \textbf{Hors sujet / Dissertation « catalogue ». Note éliminatoire.}
    \item \textbf{Bon :} « \emph{Problématique :} L'État est-il la condition indispensable de la vie commune ou une menace pour la liberté qu'il prétend protéger ? \textbf{Thèse 1} : Oui, il est nécessaire car sans lui c'est l'insécurité (Hobbes : Léviathan). \textbf{Thèse 2} : Mais cette nécessité engendre le risque de toute-puissance (Weber : monopole de la violence ; Arendt : totalitarisme). \textbf{Synthèse} : Seule la démocratie constitutionnelle (séparation des pouvoirs, droits fondamentaux, Art. 1er Constitution camerounaise) permet de concilier nécessité de l'État et liberté des citoyens. »
\end{itemize}
\end{attention}

\courssub{Méthode de justification d'une conclusion : La structure « T.A.R.E.R »}

Au terme de votre dissertation ou de votre commentaire, la conclusion n'est pas un résumé. C'est votre \textbf{prise de position finale}, solidement ancrée.

\begin{methode}[Justifier sa conclusion : La méthode T.A.R.E.R]
\begin{enumerate}
    \item \textbf{T — Thèse (Position claire)} : Répondez \emph{oui}, \emph{non}, ou \emph{oui mais...} / \emph{ni tout à fait l'un ni tout à fait l'autre}. Une phrase nette. \emph{Ex. : « L'État est nécessaire, mais seulement s'il est un État de droit démocratique. »}
    \item \textbf{A — Argument principal (Le plus fort)} : Résumez en une phrase l'argument décisif de votre 3e partie (synthèse). \emph{Ex. : « Seule la séparation des pouvoirs et le respect des droits fondamentaux transforment la contrainte étatique en liberté politique. »}
    \item \textbf{R — Référence philosophique précise} : Citez \textbf{un auteur} et \textbf{une œuvre/concept} qui valide votre thèse. \emph{Ex. : « Comme l'écrit Kelsen dans \emph{Théorie pure du droit}, la norme fondamentale (Grundnorm) confère légitimité à l'ordre juridique seulement si elle est respectée par les organes qu'elle crée. »}
    \item \textbf{E — Exemple concret (Cameroun / Actualité)} : Ancrez dans le réel. \emph{Ex. : « Au Cameroun, l'article 1er de la Constitution pose le principe, mais l'effectivité du contrôle juridictionnel (Conseil constitutionnel, Tribunal administratif) reste le test de la réalité de l'État de droit. »}
    \item \textbf{R — Réponse à l'objection adverse (Ouverture)} : Montrez que vous avez vu l'argument contraire et pourquoi il ne renverse pas votre thèse. \emph{Ex. : « Certes, l'État de droit peut sembler lent ou inefficace face aux urgences (sécurité, développement), mais l'histoire montre que l'arbitraire « efficace » mène toujours à l'effondrement (ex. années de braise, crises anglophones). La liberté n'est pas un luxe, c'est la condition de la stabilité durable. »}
\end{enumerate}
\end{methode}

\begin{exemplebox}[Sujet type Bac — Application de la méthode T.A.R.E.R]
\textbf{Sujet :} « \emph{La force suffit-elle à fonder l'État ?} »

\textbf{Conclusion attendue (Modèle T.A.R.E.R) :}
\begin{itemize}
    \item \textbf{T (Thèse) :} \textbf{Non}, la force \emph{seule} ne suffit pas à fonder l'État ; elle ne produit que la soumission, non l'obéissance légitime.
    \item \textbf{A (Argument) :} La force sans légitimité engendre la révolte latente ou ouverte ; l'État durable repose sur le \emph{consentement} (contrat social) et la \emph{légalité rationnelle} (Weber).
    \item \textbf{R (Référence) :} \textbf{Rousseau}, \emph{Du contrat social}, I, 3 : « Le plus fort n'est jamais assez fort pour être toujours le maître, s'il ne transforme sa force en droit et l'obéissance en devoir. »
    \item \textbf{E (Exemple) :} Au Cameroun, la répression pure (force) face aux revendications syndicales ou régionales (2016-2017) a exacerbé le conflit ; le \emph{Grand Dialogue National} (2019) a cherché à restaurer la \emph{légitimité par la parole et la réforme} (décentralisation, statut spécial).
    \item \textbf{R (Réponse objection) :} \emph{Objection :} « Mais sans force (police, armée), l'État s'effondre ! » \textbf{Réponse :} La force est \emph{nécessaire} (monopole de la violence légitime — Weber) pour \emph{faire respecter} la loi, mais elle n'est pas le \emph{fondement} de l'État. Le fondement, c'est la \cle{légitimité} (croyance en la validité de l'ordre — Weber) issue du suffrage et du respect des droits.
\end{itemize}
\end{exemplebox}

\begin{savaistu}[Le Conseil constitutionnel camerounais : Gardien de la régularité électorale]
L'article 65 de la Constitution de 1996 dispose : « Le Conseil constitutionnel veille à la régularité de l'élection du Président de la République et des élections législatives. Il proclame les résultats. » C'est un exemple concret de \textbf{contrôle institutionnel du pouvoir} : l'organe juridictionnel (judiciaire) contrôle l'organe politique (exécutif/législatif) sur la procédure même qui fonde la \cle{légitimité démocratique}. Depuis 2018, ses décisions (contentieux électoraux) sont publiées et font jurisprudence. C'est la concrétisation camerounaise de la \cle{séparation des pouvoirs} (Montesquieu) et du \cle{constitutionnalisme} (Kelsen) : le pouvoir constituant (peuple) $\rightarrow$ Constitution $\rightarrow$ organes constitués (Président, Parlement, Conseil constitutionnel) qui se contrôlent mutuellement.
\end{savaistu}

\begin{aretenir}[Synthèse de l'atelier — Ce qu'il faut maîtriser pour l'examen]
\begin{itemize}
    \item La \textbf{méthode en 4 étapes} (Faits $\rightarrow$ Notions $\rightarrow$ Problème $\rightarrow$ Thèse justifiée) pour toute situation politique.
    \item Les \textbf{3 situations types} : Impôt/Contrat social ; Chef/Maire (Légitimités) ; Homme fort/Démocratie (Hobbes vs Montesquieu).
    \item L'\textbf{analyse constitutionnelle} (Art. 1 \& 2) : République, Décentralisation, Indivisibilité, Laïcité, Démocratie, Sociale, Souveraineté, Représentation, Symboles.
    \item La \textbf{problématisation} : Transformer le sujet en question opposant 2 thèses + Annoncer un plan en 3 parties dialectiques.
    \item La \textbf{justification de conclusion (T.A.R.E.R)} : Thèse + Argument + Référence + Exemple + Réponse à l'objection.
    \item \textbf{Ne jamais réciter le cours.} Toujours \textbf{répondre à la question posée} en mobilisant le cours comme \emph{boîte à outils}.
\end{itemize}
\end{aretenir}

\begin{exercice}[Entraînement Bac — À faire chez vous]
\begin{enumerate}
    \item Appliquez la méthode des 4 étapes à cette situation : « Une association de parents d'élèves refuse la nomination d'un nouveau proviseur nommé par le Ministère, au profit d'un enseignant local « qui connaît les problèmes de l'établissement ». » (Identifiez : acteurs, notions [légitimité administrative vs légitimité locale/charismatique], problème, thèse).
    \item Problématisez et annoncez le plan pour le sujet : « \emph{L'État a-t-il pour seule fonction d'assurer la sécurité ?} »
    \item Rédigez une conclusion T.A.R.E.R pour le sujet : « \emph{La loi est-elle l'expression de la volonté générale ?} » (Utilisez Rousseau, Montesquieu, Kelsen, Constitution camerounaise Art. 1er).
\end{enumerate}
\end{exercice}

\begin{corrige}[Exercice 2 — Pistes de correction]
\textbf{Sujet :} « L'État a-t-il pour seule fonction d'assurer la sécurité ? »
\textbf{Analyse termes :} « Seule fonction » = exclusivité. « Sécurité » = protection vie/biens (police, armée, justice).
\textbf{Opposition :}
\begin{itemize}
    \item \textbf{Thèse (Hobbes, État minimal / Nozick) :} Oui, la sécurité (sortir de l'état de nature) est la \emph{raison d'être} première. Tout le reste (économie, culture, éducation) relève de la société civile.
    \item \textbf{Antithèse (État providence, État social, Rousseau, Marx, Constitution camerounaise Art. 1er « sociale ») :} Non, l'État moderne a des fonctions \emph{positives} : éducation, santé, redistribution, développement, protection environnement, cohésion sociale. La sécurité seule = « État gendarme » (night-watchman state), insuffisant pour la justice sociale.
\end{itemize}
\textbf{Problématique :} \textbf{« La fonction sécuritaire épuise-t-elle la finalité de l'État ou n'en est-elle que la condition nécessaire, insuffisante sans les fonctions sociales et politiques ? »}
\textbf{Plan :}
\begin{enumerate}
    \item La sécurité comme fondement et condition \emph{sine qua non} de l'État (Hobbes, Weber, fonctions régaliennes).
    \item L'élargissement historique et normatif des fonctions : État social, État providence, État développeur (Bismarck, Beveridge, Keynes, Art. 1er Const. Cameroun « sociale », ODD ONU).
    \item L'État \emph{complet} : sécurité \textbf{ET} justice sociale \textbf{ET} liberté politique (État de droit social et démocratique) — synthèse par le constitutionnalisme et la démocratie participative.
\end{enumerate}
\end{corrige}

\newpage

\coursec{Activité d'intégration}

\coursec{Leçon 07 : L'État et le pouvoir}

\courssub{Objectifs de l'activité}
\begin{itemize}
    \item Mobiliser les notions de la leçon (\cle{définition de l'État}, \cle{fondement de l'autorité}, \cle{formes de gouvernement}, \cle{rôle et finalité}) pour analyser une situation concrète ancrée dans le contexte camerounais.
    \item Exercer le \cle{savoir-faire} « Appliquer les notions de l'unité à des situations concrètes de la vie courante » et « Rédiger et justifier une démarche, une réponse ou une conclusion ».
    \item S'entraîner à la \cle{dissertation philosophique} courte : thèse, arguments structurés, exemples précis, références d'auteurs, traitement de l'objection.
    \item Développer l'esprit critique face aux discours anti-institutionnels en distinguant \cle{dysfonctionnement} et \cle{inutilité ontologique} de l'État.
\end{itemize}

\courssub{Situation : Dossier documentaire « Mon quartier a-t-il besoin de l'État ? »}
\textbf{Contexte :} Vous habitez le quartier \emph{Nkolbisson}, une commune d'arrondissement de la ville de Yaoundé (Cameroun). Une assemblée générale de quartier est convoquée pour discuter de la pertinence de la présence de l'État au quotidien. Un habitant, \emph{Papa Paul}, tient un discours virulent contestant l'utilité de l'État. Le conseil de quartier vous mandate, en tant que lycéens en Terminale, pour analyser la situation et produire une synthèse argumentée à partir du dossier ci-dessous.

\medskip
\noindent\textbf{Document 1 : Extrait de la Constitution de la République du Cameroun (Loi n°96-06 du 18 janvier 1996)}
\begin{quote}
\textbf{Article 1er :} Le Cameroun est une République décentralisée. Il est un et indivisible, laïque, démocratique et dédié au service social.\\
\textbf{Article 2 :} (1) La souveraineté nationale appartient au peuple qui l'exerce soit par référendum, soit par l'intermédiaire de ses représentants élus. (2) Aucune fraction du peuple, aucun corps, aucun individu ne peut s'en attribuer l'exercice. (3) Le suffrage est universel, égal et secret. [...] (5) L'État assure la protection des minorités et la liberté d'opinion dans le respect de la loi.
\end{quote}

\medskip
\noindent\textbf{Document 2 : Extrait du \emph{Contrat social} de Jean-Jacques Rousseau (Livre I, Chapitre 6)}
\begin{quote}
« Trouver une forme d'association qui défende et protège de toute la force commune la personne et les biens de chaque associé, et par laquelle chacun, s'unissant à tous, n'obéisse pourtant qu'à lui-même et reste aussi libre qu'auparavant. » \\
« Cette formule nous donne, pour le contrat social, des termes qui, bien entendus, se réduisent à celui-ci : \textbf{chacun de nous met en commun sa personne et toute sa puissance sous la suprême direction de la volonté générale} ; et nous recevons encore chaque membre comme partie indivisible du tout. »
\end{quote}

\medskip
\noindent\textbf{Document 3 : Données chiffrées des services publics de la Commune de Nkolbisson (Année budgétaire 2023-2024)}
\begin{center}
\begin{tabularx}{\linewidth}{|l|X|X|}\hline
\rowcolor{popblueL}
\textbf{Domaine} & \textbf{Infrastructures / Effectifs publics} & \textbf{Financement (en millions FCFA)} \\ \hline
\textbf{Éducation} & 5 écoles publiques maternelles et primaires (3 200 élèves) ; 1 lycée public (1 800 élèves) & Subventions État (personnel, manuels) : 180 \\ \hline
\textbf{Santé} & 2 Centres de Santé Intégrés (CSI) publics ; 1 hôpital de district (référence) & Subventions État (médicaments, personnel) : 95 \\ \hline
\textbf{Sécurité} & 1 Commissariat de police (45 agents) ; 1 Brigade de gendarmerie (proximité) & Budget fonctionnement (État) : 60 \\ \hline
\textbf{Voirie / Eau / Électricité} & 12 km de routes bitumées (entretien) ; Réseau eau potable (CAMWATER) ; Réseau électrique (ENEO) & Budget communal investissement : 45 \\ \hline
\textbf{Total Budget Communal} & \textbf{Recettes : 380 M FCFA} (Dotations État : 280 M | Fiscalité locale : 100 M) & \textbf{Dépenses : 375 M FCFA} \\ \hline
\end{tabularx}
\end{center}
\legende{Source : Données fictives réalistes inspirées des comptes administratifs types des Communes d'arrondissement de Yaoundé (MINEPAT / MINAT).}

\medskip
\noindent\textbf{Document 4 : Témoignage de Papa Paul, 52 ans, menuisier, résident à Nkolbisson depuis 20 ans}
\begin{quote}
« Écoutez-moi bien, mes enfants. Moi, Papa Paul, je me lève à 5h, j'ouvre mon atelier, je travaille dur. Chaque fin de mois, l'État vient me prendre ma patente, ma taxe sur le chiffre d'affaires, la TVA sur mes bois. Et qu'est-ce que je vois en retour ? 
La route devant mon atelier est un trou depuis trois ans ; ma femme a accouché par terre au CSI parce qu'il n'y avait pas de drap, pas de gant, pas de médecin la nuit. Mes gamins ? Je les ai mis au privé parce que l'école publique, c'est la grève tout le temps, les profs pas payés, les toits qui fuient. 
La sécurité ? Quand les \emph{« microbes »} (jeunes délinquants) ont volé ma scie circulaire l'an dernier, le commissariat m'a dit : « Allez chercher vous-même, on n'a pas d'essence pour la moto. » J'ai payé 15 000 FCFA aux agents pour qu'ils daignent bouger. 
Franchement, l'État, c'est quoi ? Une machine à pomper l'argent du pauvre pour engraisser les gros à Yaoundé. Nous, ici, on s'organise : on a cotisé pour réparer le transformateur électrique, on a payé des vigiles pour la nuit, on nettoie nos caniveaux nous-mêmes. \textbf{Si l'État disparaissait demain, on s'en sortirait mieux !} C'est un \textbf{mal nécessaire} au mieux, un parasite en réalité. »
\end{quote}

\courssub{Consignes}
\begin{enumerate}
    \item \textbf{Analyse documentaire (Groupe) :} 
    \begin{enumerate}
        \item Dans le Document 1, identifiez les \textbf{quatre éléments constitutifs de l'État} (population, territoire, souveraineté, organisation politique) et la \textbf{forme de gouvernement} indiquée (République, décentralisée, laïque, démocratique).
        \item Dans le Document 3, relevez les \textbf{quatre fonctions régaliennes ou sociales} exercées par l'État à Nkolbisson (éducation, santé, sécurité, aménagement du territoire) et calculez le \textbf{poids financier de l'État} dans le budget communal (part des dotations de l'État dans les recettes totales).
    \end{enumerate}
    \item \textbf{Classification doctrinale (Groupe) :} 
    \begin{enumerate}
        \item Relisez le Document 4 (Papa Paul). Classez ses arguments principaux dans trois colonnes correspondant aux doctrines étudiées : 
        \begin{itemize}
            \item \textbf{Libéralisme classique} (État minimal, protection vie/propriété, méfiance pouvoir) ;
            \item \textbf{Anarchisme} (Rejet de l'État, auto-organisation, société sans contrainte) ;
            \item \textbf{Contractualisme} (Rousseau / Locke : rupture du contrat, État ne remplit pas sa fin protectrice).
        \end{itemize}
        \item Justifiez chaque classement par une citation précise du témoignage.
    \end{enumerate}
    \item \textbf{Construction d'un tableau comparatif (Groupe) :} 
    \begin{enumerate}
        \item Construisez un tableau \textbf{« Avec État / Sans État »} à partir des données du dossier (Docs 1, 3, 4). Cinq lignes minimum : Sécurité, Éducation/Santé, Infrastructures, Libertés/Droits, Cohésion nationale.
        \item Indiquez pour chaque ligne : \textit{Réalité actuelle (Doc 3)}, \textit{Argument « Sans État » (Doc 4)}, \textit{Risque réel « Sans État » (raisonnement philosophique)}.
    \end{enumerate}
    \item \textbf{Production écrite individuelle (Évaluation) :} 
    \begin{enumerate}
        \item Rédigez une conclusion argumentée de \textbf{15 à 20 lignes} maximum répondant à la question : 
        \begin{center}
        \textbf{« L'État est-il un mal nécessaire ou un instrument du bien commun ? »}
        \end{center}
        \item \textbf{Contraintes obligatoires :} 
        \begin{itemize}
            \item Thèse claire en ouverture.
            \item Au moins \textbf{deux arguments philosophiques distincts} (ex: protection droits, services publics, volonté générale, monopole violence légitime).
            \item Citation explicite d'\textbf{au moins deux auteurs} étudiés en classe (ex: Hobbes, Locke, Rousseau, Montesquieu, Kelsen, Rawls).
            \item Utilisation d'\textbf{au moins un exemple camerounais précis} (Doc 1, 3 ou connaissance personnelle : décentralisation, CSI, école publique, Constitution).
            \item Traitement d'une \textbf{objection} (reprendre l'argument fort de Papa Paul) et réponse.
        \end{itemize}
    \end{enumerate}
    \item \textbf{Restitution orale et correction collective :} Présentation des tableaux, lecture de 3 conclusions types, évaluation par les pairs via la \textbf{grille de justification} (Thèse / Argument / Exemple / Référence / Objection-Réponse).
\end{enumerate}

\courssub{Ressources à mobiliser}
\begin{definition}[Rappel : Les 4 éléments constitutifs de l'État (Droit constitutionnel / Science politique)]
\begin{itemize}
    \item \textbf{Population :} Ensemble d'individus soumis à la même autorité (peuple).
    \item \textbf{Territoire :} Espace géographique délimité (frontières) où s'exerce la souveraineté.
    \item \textbf{Souveraineté :} Pouvoir suprême, exclusif, inaliénable (interne : monopole violence légitime ; externe : indépendance).
    \item \textbf{Organisation politique (Gouvernement) :} Ensemble d'institutions qui font la loi, l'appliquent, la jugent (Séparation des pouvoirs, Administration).
\end{itemize}
\end{definition}

\begin{definition}[Fondements de l'autorité (Doctrines)]
\begin{itemize}
    \item \textbf{Contractualisme (Hobbes, Locke, Rousseau) :} L'État naît d'un contrat. \textbf{Hobbes} (\emph{Léviathan}) : État = rempart contre la « guerre de tous contre tous » (sécurité). \textbf{Locke} : État = protection droits naturels (vie, liberté, propriété). \textbf{Rousseau} (\emph{Contrat social}) : État = expression de la \cle{volonté générale} pour la liberté civile.
    \item \cle{Légalité / Légitimité (Weber) :} Domination légale-rationnelle (lois), traditionnelle (coutumes), charismatique (chef).
    \item \cle{Séparation des pouvoirs (Montesquieu, \emph{De l'esprit des lois}) :} Législatif, Exécutif, Judiciaire pour éviter la tyrannie.
\end{itemize}
\end{definition}

\begin{definition}[Formes de gouvernement (Aristote / Montesquieu / Droit moderne)]
\begin{itemize}
    \item \textbf{Monarchie / République} (Chef de l'État : héréditaire ou élu).
    \item \textbf{Démocratie / Dictature} (Origine du pouvoir : peuple ou force).
    \item \textbf{Unitaire / Fédéral / Décentralisé} (Répartition territoire). \rightarrow \textbf{Cameroun = République unitaire décentralisée.}
    \item \textbf{Parlementaire / Présidentiel / Semi-présidentiel} (Rapports Exécutif/Législatif). \rightarrow \textbf{Cameroun = Régime présidentiel.}
\end{itemize}
\end{definition}

\begin{definition}[Rôle et finalité de l'État (Fonctions)]
\begin{itemize}
    \item \textbf{Fonctions régaliennes (indispensables) :} Sécurité intérieure (police), défense nationale (armée), justice, monnaie, diplomatie.
    \item \textbf{Fonctions sociales / de service public :} Éducation, santé, infrastructure, protection sociale, environnement, culture. \cle{Principe d'égalité d'accès}.
    \item \textbf{Finalité philosophique :} \cle{Bien commun} (Aristote, Thomas d'Aquin, Rousseau) vs \cle{Intérêt particulier}. L'État n'est pas une fin en soi, mais un \cle{moyen} au service de la vie bonne (\emph{eu zên}).
\end{itemize}
\end{definition}

\begin{methode}[Rédiger une conclusion argumentée (Dissertation courte)]
\begin{enumerate}
    \item \textbf{Accroche / Thèse (1 phrase) :} Réponse directe à la question (ex: « L'État, loin d'être un mal nécessaire, est l'instrument indispensable du bien commun... »).
    \item \textbf{Argument 1 + Auteur + Exemple :} « D'abord, [idée]. Comme l'écrit [Auteur], "[citation/idée]". Au Cameroun, [exemple Doc 3 ou réel] l'illustre. »
    \item \textbf{Argument 2 + Auteur + Exemple :} « Ensuite, [idée 2]. [Auteur 2] montre que... La Constitution de 1996 (Art. 1) garantit... »
    \item \textbf{Objection (Papa Paul) + Réponse :} « Certes, comme le dit Papa Paul, [dysfonctionnement réel]. Mais [distinction dysfonctionnement/essence] : l'État imparfait se réforme, il ne se supprime pas (Rousseau : "la volonté générale est toujours droite"). »
    \item \textbf{Ouverture (1 phrase) :} Perspective (ex: « Le défi camerounais n'est pas la suppression de l'État, mais son efficacité par la reddition des comptes et la participation citoyenne. »).
\end{enumerate}
\end{methode}

\courssub{Démarche et corrigé commenté}
\begin{exoresolu}[Corrigé guidé de l'activité d'intégration]
\noindent\textbf{Énoncé de l'activité :} Voir les Consignes 1 à 5 ci-dessus (Analyse docs, Classification doctrines, Tableau comparatif, Rédaction conclusion, Restitution).

\tcblower

\noindent\textbf{1. Analyse documentaire (Corrigé Consigne 1)}
\begin{itemize}
    \item \textbf{Doc 1 (Constitution) :}
    \begin{itemize}
        \item \textit{Population :} « Le peuple » (Art. 2), « minorités » (Art. 2).
        \item \textit{Territoire :} Implicite dans « République », « Cameroun » (un et indivisible, Art. 1).
        \item \textit{Souveraineté :} « Appartient au peuple » (Art. 2), exercée par référendum ou représentants. \cle{Monopole de la légitimité}.
        \item \textit{Organisation :} « République décentralisée », « laïque », « démocratique » (Art. 1). Séparation des pouvoirs (implicitement, Titre II-IV Constitution).
    \end{itemize}
    \item \textbf{Doc 3 (Données chiffrées) :}
    \begin{itemize}
        \item \textit{Fonctions identifiées :} 
        \begin{enumerate}
            \item \textbf{Éducation} (Écoles, Lycée) — Service public social.
            \item \textbf{Santé} (CSI, Hôpital) — Service public social / Protection vie.
            \item \textbf{Sécurité} (Commissariat, Gendarmerie) — Fonction régalienne (monopole violence légitime).
            \item \textbf{Aménagement / Voirie / Réseaux} — Service public économique / Cadre de vie.
        \end{enumerate}
        \item \textit{Calcul poids financier :} Dotations État = 280 M FCFA. Recettes totales = 380 M FCFA. 
        \[ \frac{280}{380} \times 100 \approx \textbf{73,7 \%} \]
        \textbf{Interprétation :} La commune dépend très largement de l'État central pour exister. Sans dotations, effondrement des services.
    \end{itemize}
\end{itemize}

\noindent\textbf{2. Classification doctrinale du témoignage de Papa Paul (Corrigé Consigne 2)}

\begin{center}
\begin{tabularx}{\linewidth}{|l|X|X|}\hline
\rowcolor{popblueL}
\textbf{Doctrine} & \textbf{Arguments de Papa Paul (Doc 4)} & \textbf{Justification / Nuance} \\ \hline
\textbf{Libéralisme classique} & « L'État vient me prendre ma patente, ma taxe... » ; « Police : pas d'essence, j'ai payé 15 000 FCFA pour qu'ils bougent. » & Critique de la \cle{prédation fiscale} et de l'\cle{inefficacité} de l'État minimal (Locke/Nozick) : l'État ne remplit pas sa fonction basique (protection propriété/sécurité) malgré l'impôt. \\ \hline
\textbf{Anarchisme} & « Nous, on s'organise : cotisé pour transformateur, payé vigiles, nettoyé caniveaux. » ; « Si l'État disparaissait, on s'en sortirait mieux. » & Apologie de l'\cle{auto-organisation} (Proudhon, Bakounine) : la société civile peut remplacer l'État par la libre association, l'entraide, sans contrainte politique. \\ \hline
\textbf{Contractualisme (Rousseau/Locke)} & « École publique : grève, profs pas payés, toits qui fuient. » ; « CSI : accouchement par terre, pas de drap. » & \cle{Rupture du contrat} : L'État ne protège pas les biens (éducation, santé = vie/liberté). Pour Locke, droit de résistance ; pour Rousseau, la \cle{volonté générale} est trahie par l'administration (volonté particulière des élites). \\ \hline
\end{tabularx}
\end{center}

\noindent\textbf{3. Tableau comparatif « Avec État / Sans État » (Corrigé Consigne 3)}

\begin{center}
\begin{tabularx}{\linewidth}{|l|X|X|X|}\hline
\rowcolor{popblueL}
\textbf{Dimension} & \textbf{Réalité AVEC État (Docs 1, 3)} & \textbf{Argument « SANS État » (Doc 4 / Papa Paul)} & \textbf{Risque réel « SANS État » (Analyse philosophique)} \\ \hline
\textbf{Sécurité / Justice} & Commissariat, Gendarmerie, Justice (monopole violence légitime). Budget 60 M. & Auto-défense : vigiles payés, « on cherche soi-même ». & \textbf{Guerre de tous contre tous (Hobbes)} : Justice privée = loi du plus fort, vendettas, pas de procédure équitable, insécurité majeure pour les plus faibles. \\ \hline
\textbf{Éducation / Santé} & 5 écoles, 1 lycée (3200+1800 élèves), 2 CSI, 1 Hôpital. Subventions 275 M. & École privée pour les siens ; CSI défaillant = argument pour partir. & \textbf{Inégalités radicales} : Seuls les riches accèdent à l'instruction/soins. Perte de \cle{capital humain} national. Violation Droit à l'éducation (Constitution, PIDESC). \\ \hline
\textbf{Infrastructures / Services} & Routes (12 km), Eau (CAMWATER), Électricité (ENEO). Budget investissement 45 M. & Cotisation transformateur, nettoyage caniveaux (auto-organisation locale). & \textbf{Effet de seuil / Externalités} : Réseaux eau/électricité impossibles à l'échelle locale. Routes : coordination inter-quartiers impossible. Désertification territoires non rentables. \\ \hline
\textbf{Libertés / Droits / Citoyenneté} & Constitution Art. 1 \& 2 : Suffrage universel, laïcité, protection minorités, service public. & « L'État pompe l'argent pour engraisser les gros. » (Sentiment d'aliénation). & \textbf{Perte de la citoyenneté} : Plus de cadre juridique commun, plus de droits opposables (ex: grève, plainte), arbitraire des puissants locaux (chefs traditionnels, seigneurs de guerre). \\ \hline
\textbf{Cohésion nationale / Unité} & « Un et indivisible » (Art. 1). Redistribution (Dotations 73,7\% budget). Solidarité nationale. & « Nous, ici, on gère. » (Repli identitaire / communalisme). & \textbf{Balkanisation / Sécessionnisme} : Fragmentation en micro-entités hostiles. Fin de la redistribution Nord/Sud, Urbain/Rural. Conflits pour ressources (eau, terres). \\ \hline
\end{tabularx}
\end{center}

\noindent\textbf{4. Proposition de conclusion argumentée type (Corrigé Consigne 4 — \textit{Exemple élève})}

\begin{quote}
\textbf{Thèse :} L'État n'est pas un « mal nécessaire » au sens d'un fardeau fatal, mais l'\textbf{instrument indispensable du bien commun}, condition de possibilité de la liberté effective et de la justice sociale, pourvu qu'il soit démocratique et responsable.

\textbf{Argument 1 (Sécurité / Droit) :} D'abord, comme l'affirme \textbf{Hobbes} dans \emph{Le Léviathan}, sans État c'est la « guerre de tous contre tous » où la vie est « misérable, brutale et courte ». Le monopole de la violence légitime (Weber) permet seul de substituer le droit à la force. Au Cameroun, malgré les dysfonctionnements relevés par Papa Paul (corruption, manque d'essence), le commissariat et la gendarmerie de Nkolbisson (Doc 3) offrent un cadre légal de recours impossible dans l'auto-défense anarchique.

\textbf{Argument 2 (Services publics / Égalité) :} Ensuite, \textbf{Rousseau} définit l'État comme l'association qui « défend et protège la personne et les biens de chaque associé » par la \cle{volonté générale}. L'État-providence (éducation, santé Docs 1 \& 3) réalise l'égalité des chances : les 5 000 élèves publics de Nkolbisson (Doc 3) n'auraient aucune alternative sans l'école publique gratuite. La Constitution de 1996 (Art. 1 : « dédié au service social ») inscrit cette finalité dans le marbre.

\textbf{Objection (Papa Paul) \& Réponse :} Certes, Papa Paul a raison de dénoncer la \cle{prédation fiscale} et la \cle{mauvaise gestion} (routes, CSI, corruption). Mais confondre \textbf{dysfonctionnement} (État mal géré, capté par des intérêts particuliers) et \textbf{inutilité ontologique} est une erreur logique (sophisme). La réponse n'est pas la suppression (anarchisme) mais la \cle{réforme démocratique} : reddition des comptes, décentralisation effective (Art. 1 Constitution), participation citoyenne (budgets participatifs).

\textbf{Ouverture :} Le défi du Cameroun n'est pas de se passer d'État, mais de construire un \textbf{État de droit efficace} où l'impôt de Papa Paul finance visiblement le bitume, le drap du CSI et le salaire du prof, réalisant ainsi la \cle{justice sociale} promise par le contrat social.
\end{quote}

\noindent\textbf{5. Grille d'évaluation de la restitution (Corrigé Consigne 5)}
\begin{center}
\begin{tabularx}{\linewidth}{|l|X|X|}\hline
\rowcolor{popblueL}
\textbf{Critère} & \textbf{Indicateur de réussite (Niveau Bac/Probatoire)} & \textbf{Points} \\ \hline
\textbf{Thèse} & Claire, nuancée (pas binaire), répond directement au sujet. & /2 \\ \hline
\textbf{Argumentation} & 2 arguments philosophiques distincts, logiquement enchaînés. & /4 \\ \hline
\textbf{Références auteurs} & 2 auteurs minimum cités correctement (Hobbes, Locke, Rousseau, Montesquieu, Rawls, Kelsen...). & /3 \\ \hline
\textbf{Exemple camerounais} & Précis, issu du dossier (Doc 1, 3) ou connaissance réelle (Décentralisation, PNDP, CAN 2021...), bien exploité. & /3 \\ \hline
\textbf{Traitement objection} & Reprise explicite de l'argument adverse (Papa Paul) + Réponse philosophique argumentée (distinction essence/accident). & /3 \\ \hline
\textbf{Qualité langue / Structure} & 15-20 lignes, connecteurs logiques, vocabulaire précis (souveraineté, régalien, bien commun, volonté générale). & /2 \\ \hline
\textbf{TOTAL} & & \textbf{/17} \\ \hline
\end{tabularx}
\end{center}

\end{exoresolu}

\courssub{Bilan de l'activité}
\begin{aretenir}[Ce que cette activité a permis de mettre en évidence]
\begin{itemize}
    \item \textbf{Ancrage réel :} Les élèves constatent que l'État n'est pas une abstraction : il finance 74\% du budget de leur quartier (Doc 3), gère leur école, leur santé, leur sécurité. Le « mal nécessaire » de Papa Paul est contredit par les chiffres.
    \item \textbf{Distinction conceptuelle opératoire :} Différencier \cle{État} (institution permanente, cadre juridique) et \cle{Gouvernement / Administration} (équipe dirigeante, gestionnaires). Papa Paul critique l'administration (corruption, incompétence), pas l'État comme structure. C'est la clé pour ne pas basculer dans l'anarchisme naïf.
    \item \textbf{Articulation théorie/pratique :} Le contractualisme (Rousseau, Locke, Hobbes) n'est pas une « vieille philosophie » : il donne les outils pour juger l'action publique (Le contrat est-il respecté ? La volonté générale est-elle servie ?).
    \item \textbf{Méthodologie Bac :} L'exercice impose la structure thèse/arguments/auteurs/exemple/objection/réponse/ouverture. C'est le « squelette » de la dissertation de 4h. L'entraînement sur 15-20 lignes automatise le réflexe.
    \item \textbf{Citoyenneté active :} L'activité transforme la plainte (« l'État ne fait rien ») en exigence démocratique (« l'État doit rendre des comptes, je participe au contrôle »). Passage du statut de \cle{sujet} passif à \cle{citoyen} souverain (Art. 2 Constitution).
\end{itemize}
\end{aretenir}

\begin{attention}[Pièges à éviter pour le Bac / Probatoire]
\begin{itemize}
    \item \textbf{Confondre « État » et « Gouvernement » :} Dire « L'État est corrompu » est un abus de langage. C'est tel gouvernement, telle administration. L'État (institution) reste le rempart.
    \item \textbf{Citer des auteurs « au feeling » :} « Comme dit Montesquieu... » sans lien avec l'argument (séparation des pouvoirs ? climat ?). \textbf{Sanction :} Hors-sujet partiel.
    \item \textbf{Oublier l'exemple camerounais :} Sujet technique = ancrage local obligatoire. « En France... » = 0 point à l'exemple.
    \item \textbf{Traiter l'objection par l'ignorance :} « Papa Paul a tort. » \rightarrow \textbf{Échec.} Il faut dire : « Papa Paul pointe un réel (dysfonctionnement), mais sa conclusion (suppression) est fausse car... »
    \item \textbf{Dépasser 20 lignes :} Concision exigée. Apprenez à densifier : « Hobbes (\emph{Léviathan}) montre que sans souverain, c'est la guerre de tous contre tous » (1 ligne au lieu de 3).
\end{itemize}
\end{attention}

\newpage

\coursec{Méthodes et exercices résolus}

\coursec{L'État et le pouvoir}

\courssub{Qu'est-ce que l'État ? Définition et éléments constitutifs}

\begin{definition}[L'État (approche sociologique et juridique)]
L'État est une \textbf{personne morale de droit public} qui possède quatre éléments constitutifs indissociables (théorie classique, reprise par la Convention de Montevideo 1933 et Max Weber) :
\begin{enumerate}
    \item \textbf{Une population :} Un groupe humain stable, installé durablement sur un territoire (critère démographique).
    \item \textbf{Un territoire :} Un espace géographique délimité par des frontières (terrestres, maritimes, aériennes), base spatiale de l'exercice de la souveraineté (critère territorial).
    \item \textbf{Un gouvernement (organisation politique) :} Un ensemble d'institutions (exécutif, législatif, judiciaire, administration, armée, police) qui exerce l'autorité sur la population et le territoire (critère organisationnel).
    \item \textbf{La souveraineté :} Le pouvoir suprême à l'intérieur (monopole de la contrainte, pas de rival) et l'indépendance à l'extérieur (pas de tutelle, égalité entre États). C'est l'élément \textbf{distinctif} de l'État (critère juridique/politique).
\end{enumerate}
\end{definition}

\begin{aretenir}[Distinctions fondamentales]
\begin{itemize}
    \item \textbf{État vs Nation :} L'État est une structure politico-juridique ; la Nation est une communauté culturelle, historique, linguistique (fait sociologique). \emph{État-nation} = coïncidence idéale ; \emph{État multinational} (ex: Cameroun, Belgique) ou \emph{Nation sans État} (ex: Kurdes, Palestiniens).
    \item \textbf{État vs Gouvernement :} L'État est \textbf{permanent} (institutions, territoire, population) ; le Gouvernement est \textbf{temporaire} (équipe dirigeante, majorité politique). « L'État, c'est moi » (Louis XIV) est une confusion absolutiste.
    \item \textbf{État vs Société civile :} L'État détient le monopole de la violence légitime ; la société civile = ensemble des associations, entreprises, syndicats, familles, autonomes vis-à-vis de l'État.
\end{itemize}
\end{aretenir}

\begin{exemplebox}[Les éléments de l'État au Cameroun]
\begin{itemize}
    \item \textbf{Population :} ~28 millions d'habitants (estimation 2024), diversité ethnique (~250 groupes), bilingue officiel (français/anglais).
    \item \textbf{Territoire :} 475 442 km$^2$, frontières héritées de la colonisation (allemande, française, britannique), contentieux résolus (Bakassi, Lac Tchad).
    \item \textbf{Gouvernement :} Président de la République (Chef de l'État), Premier Ministre (Chef du Gouvernement), Parlement (Assemblée nationale, Sénat), Administration décentralisée/déconcentrée, Forces de défense.
    \item \textbf{Souveraineté :} Indépendance (1er janv 1960 / 1er oct 1961), membre ONU, UA, CEMAC, Commonwealth, OIF. Souveraineté interne affirmée (lutte contre Boko Haram, crise NOSO), mais défis (pression extérieure, dette, insécurité localisée).
\end{itemize}
\end{exemplebox}

\courssub{Le fondement de l'État : d'où vient son autorité ?}

\begin{definition}[Légitimité, Légalité, Domination (Max Weber)]
\begin{itemize}
    \item \textbf{Légalité :} Conformité aux règles juridiques en vigueur (constitution, lois, procédures). \emph{Ex:} Élection respectant le code électoral.
    \item \textbf{Légitimité :} \textbf{Croyance subjective} des dominés en la validité de l'ordre de domination, en le « droit de commander » du dirigeant. C'est le fondement \emph{moral} et \emph{sociologique} de l'obéissance.
    \item \textbf{Domination :} « Probabilité de trouver obéissance à un commandement donné » (Weber). L'État revendique le \textbf{monopole de la violence physique légitime} sur un territoire.
\end{itemize}
\end{definition}

\begin{propriete}[Les trois types purs de domination légitime (Max Weber, \emph{Économie et Société})]
\begin{center}
\begin{tabularx}{\linewidth}{|l|X|X|X|}\hline
\textbf{Type} & \textbf{Fondement de la validité} & \textbf{Base de l'obéissance} & \textbf{Exemple type} \\ \hline
\textbf{Traditionnelle} & Coutume, tradition immémoriale, héritage (droit du sang, investiture rituelle) & Fidélité personnelle au « seigneur » / au chef dépositaire de la tradition & Monarchies absolues, Chefferies coutumières (Cameroun : Chefs de 1er, 2e, 3e degré) \\ \hline
\textbf{Charismatique} & Qualités exceptionnelles du chef (héroïsme, prophétisme, vertu, grâce), reconnaissance personnelle & Dévouement personnel au « héros » / au prophète, preuve constante (miracles, victoires) & Leaders révolutionnaires (Gandhi, Mandela, Sankara), Prophètes, Chefs de guerre victorieux \\ \hline
\textbf{Rationnelle-légale} & Règles juridiques rationnelles, procédures codifiées, compétence technique & Obéissance à la \emph{fonction} (impersonnelle), à la loi, non à la personne & Démocraties modernes, Bureaucratie, État de droit (Cameroun : Président élu, Parlement, Administration) \\ \hline
\end{tabularx}
\end{center}
\cle{Routinisation} : La domination charismatique est instable ; pour durer, elle doit se \emph{routiniser} en traditionnelle (héritage dynastique) ou en légale (institutionalisation, élections).
\end{propriete}

\begin{attention}[Légalité $\neq$ Légitimité]
\begin{itemize}
    \item Un putschiste qui fait voter une constitution par référendum contrôlé a la \textbf{légalité} (texte voté) mais pas la \textbf{légitimité} (croyance libre du peuple).
    \item Un chef coutumier respecté par sa communauté a la \textbf{légitimité traditionnelle} mais pas forcément la \textbf{légalité} formelle (s'il n'est pas reconnu par l'administration).
    \item \textbf{Crise de légitimité :} Quand l'écart entre légalité formelle et adhésion réelle se creuse (abstention massive, contestation rue, répression), l'État devient fragile.
\end{itemize}
\end{attention}

\begin{savaistu}[Le contrat social : fondement philosophique de la légitimité moderne]
\begin{itemize}
    \item \textbf{Hobbes (\emph{Léviathan}, 1651) :} État de nature = « guerre de tous contre tous ». Contrat de \textbf{soumission} à un Souverain absolu pour la \textbf{Sécurité}. Légitimité = Utilité (paix civile).
    \item \textbf{Locke (\emph{Gouvernement civil}, 1690) :} État de nature = insécurité des biens. Contrat de \textbf{confiance} pour protéger \textbf{Vie, Liberté, Propriété}. Légitimité = Consentement + Respect droits naturels. Droit de résistance si tyrannie.
    \item \textbf{Rousseau (\emph{Contrat social}, 1762) :} « Trouver une forme d'association qui défende la personne et les biens... et où chacun n'obéisse qu'à lui-même ». Contrat d'\textbf{association} : aliénation totale à la \textbf{Volonté Générale}. Légitimité = \textbf{Souveraineté populaire}. « Obéir à la loi qu'on s'est prescrite, c'est la liberté ».
\end{itemize}
\end{savaistu}

\courssub{Les formes de gouvernements}

\begin{definition}[Forme de l'État vs Forme de gouvernement]
\begin{itemize}
    \item \textbf{Forme de l'État (Organisation territoriale) :} \emph{État unitaire} (souveraineté unique, centre fort : France, Cameroun, Chine) vs \emph{État fédéral} (partage de souveraineté entre État fédéral et États fédérés : USA, Allemagne, Nigeria, Afrique du Sud). \emph{Note :} Le Cameroun est un \textbf{État unitaire décentralisé} (Loi n°2019/024), pas un État fédéral.
    \item \textbf{Forme de gouvernement (Relations Exécutif / Législatif) :} Définit comment le pouvoir exécutif est désigné, dont il émane, et comment il est contrôlé par le parlement.
\end{itemize}
\end{definition}

\begin{propriete}[Typologie classique des régimes démocratiques (Maurice Duverger)]
\begin{center}
\begin{tabularx}{\linewidth}{|l|X|X|X|}\hline
\textbf{Critère} & \textbf{Régime Présidentiel} & \textbf{Régime Parlementaire} & \textbf{Régime Semi-présidentiel} \\ \hline
\textbf{Chef de l'État / Chef du Gouvernement} & \textbf{Même personne} (Président) & \textbf{Séparés} (Monarque/Président symbolique + Premier Ministre politique) & \textbf{Séparés} (Président élu + Premier Ministre) \\ \hline
\textbf{Origine de l'Exécutif} & Président élu au \textbf{suffrage universel direct} & PM issu de la \textbf{majorité parlementaire} (investiture) & Président élu au \textbf{suffrage universel} + PM issu de la majorité parlementaire \\ \hline
\textbf{Responsabilité du Gouvernement devant le Parlement} & \textbf{NON} (Séparation rigide, mandats fixes) & \textbf{OUI} (Vote de censure, motion de défiance, question de confiance) & \textbf{OUI} (Pour le PM et son gouvernement) \\ \hline
\textbf{Dissolution du Parlement} & \textbf{NON} (Mandats fixes) & \textbf{OUI} (Par Chef de l'État, souvent sur avis PM) & \textbf{OUI} (Par le Président) \\ \hline
\textbf{Séparation des pouvoirs} & \textbf{Rigide} (Checks and balances) & \textbf{Souple} (Fusion/Confiance) & \textbf{Hybride} (Dualité exécutive, responsabilité partagée) \\ \hline
\textbf{Exemple type} & \textbf{USA} (Référence pure) & \textbf{Royaume-Uni, Allemagne, Japon, Espagne} & \textbf{France (V$^e$ République), Cameroun, Sénégal, Roumanie} \\ \hline
\end{tabularx}
\end{center}
\end{propriete}

\begin{exemplebox}[Le régime semi-présidentiel camerounais (Constitution 1996, révisée 2008)]
\begin{itemize}
    \item \textbf{Président de la République :} Élu au suffrage universel direct pour 7 ans (renouvelable une fois). Chef de l'État, Chef des Armées, nomme le PM et les ministres, préside le Conseil des ministres, droit de dissolution de l'Assemblée, pouvoir réglementaire propre, grâce.
    \item \textbf{Premier Ministre :} Chef du Gouvernement, nommé par le Président. Responsable devant l'Assemblée nationale (vote de censure possible, Art. 34). Conduit la politique de la nation, assure l'exécution des lois.
    \item \textbf{Parlement :} Bicaméral (Assemblée nationale + Sénat). Vote la loi, contrôle l'action du Gouvernement (questions, commissions d'enquête, censure).
    \item \textbf{Particularité :} « Présidentialisme majoritaire » : Le Président domine l'exécutif (nomination PM, ministres, dissolution) ; le PM est un collaborateur direct. Cohabitation possible mais rare (majorité présidentielle stable).
\end{itemize}
\end{exemplebox}

\begin{attention}[Dictature / Régime autoritaire : Hors classification constitutionnelle démocratique]
Ce n'est pas une « forme de gouvernement » au sens juridique, mais une \textbf{pathologie du pouvoir} :
\begin{itemize}
    \item Origine illégitime (coup d'État, fraude massive).
    \item Confusion des pouvoirs (Exécutif contrôle Législatif et Judiciaire).
    \item Parti unique ou hégémonique, élections non libres (liste unique, bourrage d'urnes).
    \item Non-responsabilité du dirigeant, immunité, répression opposition.
    \item Absence d'État de droit (justice aux ordres, détentions arbitraires).
\end{itemize}
\end{attention}

\courssub{Rôle et finalité de l'État}

\begin{definition}[Les fonctions de l'État]
\begin{itemize}
    \item \textbf{Fonctions régaliennes (de souveraineté) :} \emph{Indispensables à l'existence de l'État.} Justice, Police/Sécurité intérieure, Armée/Défense nationale, Diplomatie, Monnaie, Législation. \textbf{Fondement : Monopole de la violence physique légitime (Weber).}
    \item \textbf{Fonctions sociales / Services publics :} Éducation, Santé, Transports, Énergie, Eau, Logement, Protection sociale, Culture, Sport. \textbf{Principes :} Continuité, Égalité (accès non discriminatoire), Mutabilité (adaptation), Adaptabilité (nouvelles technologies).
    \item \textbf{Fonctions économiques :} Régulation (concurrence, prix), Planification/Stratégie (PNDP, Vision 2035), Propriété publique (secteur parapublic), Redistribution fiscale (impôt progressif, transferts sociaux).
\end{itemize}
\end{definition}

\begin{aretenir}[La finalité de l'État : Le Bien Commun]
\begin{itemize}
    \item \textbf{Aristote / Thomas d'Aquin / Vatican II (\emph{Gaudium et Spes} §26) :} « L'ensemble des conditions de la vie sociale qui permettent aux groupes et à chacun de leurs membres de parvenir à leur perfection plus pleinement et plus aisément. »
    \item \textbf{Rawls (\emph{Théorie de la justice}) :} Justice comme équité. Deux principes : 1) Libertés égales maximales. 2) Inégalités sociales utiles aux plus défavorisés (principe de différence) + égalité des chances équitable.
    \item L'État n'est pas une fin en soi, mais un \textbf{moyen} au service de la personne humaine.
\end{itemize}
\end{aretenir}

\begin{methode}[Débat philosophique : Quelle étendue pour l'État ?]
\begin{enumerate}
    \item \textbf{État minimal / État de garde (Libéralisme classique, Nozick) :} Seules fonctions régaliennes (Police, Justice, Armée). Toute redistribution = vol / travail forcé. Société civile gère le reste (marché, charité).
    \item \textbf{État-providence / État social (Social-démocratie, Rawls, Droit social) :} Intervention forte pour corriger les inégalités de naissance/marché (Éducation, Santé, Retraite, Allocations). Solidarité nationale obligatoire.
    \item \textbf{État totalitaire (Arendt) :} Contrôle total de la société (économie, culture, vie privée, pensée) par le parti unique et la police politique. Négation de la personne.
    \item \textbf{Principle de Subsidiarité (Doctrine sociale de l'Église, Droit UE) :} L'État ne doit faire que ce que les échelons inférieurs (communes, associations, familles, individus) ne peuvent pas faire efficacement. Proximité et efficacité.
\end{enumerate}
\end{methode}

\courssub{Atelier d'application : analyser des situations concrètes}

\begin{methode}[Identifier les éléments constitutifs de l'État dans une situation concrète]
\textbf{Objectif :} Savoir repérer les quatre éléments constitutifs de l'État (population, territoire, gouvernement, souveraineté) dans un énoncé ou un article de presse.
\begin{enumerate}
    \item \textbf{Lire attentivement} le texte ou la description du pays imaginaire/réel.
    \item \textbf{Repérer la population :} Y a-t-il mention d'un peuple, d'habitants, d'une nation ? (Critère : permanence).
    \item \textbf{Repérer le territoire :} Y a-t-il des frontières, une superficie, un sol, des ressources naturelles ? (Critère : base spatiale).
    \item \textbf{Repérer le gouvernement :} Y a-t-il des institutions, une administration, une armée, une police, des lois votées ? (Critère : organisation politique).
    \item \textbf{Repérer la souveraineté :} Le pouvoir est-il suprême à l'intérieur (pas de rival) et indépendant à l'extérieur (pas de tutelle) ? (Critère : fondement juridique).
    \item \textbf{Conclure :} Si les quatre éléments sont présents $\rightarrow$ c'est un État. Si l'un manque $\rightarrow$ ce n'est pas un État (ex: nation sans État, gouvernement en exil).
\end{enumerate}
\cle{Règle} : L'État est une \textbf{personne morale de droit public} ; il naît de la réunion stable de ces quatre éléments.
\end{methode}

\begin{exoresolu}[Analyse de la « République de Mboa »]
\textbf{Énoncé :} La « République de Mboa » est une entité politique située en Afrique centrale. Elle compte 12 millions d'habitants (population). Son territoire de 475 000 km$^2$ est délimité par des frontières reconnues internationalement (territoire). Un Président élu au suffrage universel dirige le pays, assisté d'un Premier ministre et d'un Parlement bicaméral qui vote les lois. Une administration recouvre les impôts et une armée assure la défense (gouvernement). Cependant, une rébellion armée contrôle 30\% du territoire au Nord depuis deux ans, y percevant des taxes et rendant la justice. Par ailleurs, le pays dépend à 80\% de l'aide budgétaire d'une puissance étrangère qui impose ses choix économiques.
\textbf{Question :} La République de Mboa remplit-elle les critères de l'État ? Justifiez chaque élément.
\tcblower
\textbf{Solution détaillée :}
\begin{enumerate}
    \item \textbf{Population (OUI) :} Le texte mentionne explicitement « 12 millions d'habitants ». Il y a un peuple installé durablement. Critère rempli.
    \item \textbf{Territoire (OUI, mais fragilisé) :} Le territoire existe (475 000 km$^2$, frontières reconnues). Cependant, le contrôle effectif n'est pas total (30\% échappent à l'autorité centrale). En droit international, l'État existe tant que le territoire est \emph{défini}, même si le contrôle n'est pas parfait, mais cela affaiblit la souveraineté interne.
    \item \textbf{Gouvernement (OUI) :} Institutions complètes : Président, Premier ministre, Parlement, administration fiscale, armée. L'organisation politique est effective sur la majeure partie du territoire. Critère rempli.
    \item \textbf{Souveraineté (NON / PARTIELLE) :} C'est le point de rupture.
    \begin{itemize}
        \item \textbf{Interne :} La rébellion exerce un pouvoir concurrent (perception d'impôts, justice) sur une partie du territoire. Le monopole de la violence légitime (Weber) n'est pas assuré par l'État central. Souveraineté interne compromise.
        \item \textbf{Externe :} La dépendance budgétaire à 80\% et l'ingérence d'une puissance étrangère sur les choix économiques remettent en cause l'indépendance politique. L'État n'est pas pleinement souverain à l'extérieur.
    \end{itemize}
\end{enumerate}
\textbf{Conclusion :} Formellement, Mboa possède les attributs de l'État (reconnaissance internationale, institutions). Matériellement, c'est un \textbf{État failli} ou \textbf{État fragile} : la souveraineté, fondement même de l'État, est gravement entamée tant à l'intérieur (rivalité de pouvoir) qu'à l'extérieur (tutelle économique). Il ne remplit pas \emph{pleinement} la définition weberienne de l'État.
\end{exoresolu}

\begin{methode}[Qualifier le fondement de la légitimité du pouvoir (Weber)]
\textbf{Objectif :} Classer une justification du pouvoir parmi les trois types purs de domination légitime selon Max Weber.
\begin{enumerate}
    \item \textbf{Identifier le discours justificatif :} Sur quoi le dirigeant ou le régime s'appuie-t-il pour dire « j'ai le droit de commander » ?
    \item \textbf{Tester la \textbf{domination traditionnelle} :} Le pouvoir repose-t-il sur la coutume, l'héritage, « c'est ainsi depuis toujours » ? (Ex: Monarchie absolue, chefferies coutumières).
    \item \textbf{Tester la \textbf{domination charismatique} :} Le pouvoir repose-t-il sur la personne exceptionnelle du chef, ses qualités héroïques, prophétiques, sa « grâce » personnelle ? (Ex: Leader révolutionnaire, prophète, chef de guerre victorieux). \emph{Instable, routinisation nécessaire}.
    \item \textbf{Tester la \textbf{domination rationnelle-légale} :} Le pouvoir repose-t-il sur des règles juridiques écrites, une constitution, des procédures électorales, la compétence technique ? L'obéissance va à la \emph{fonction}, pas à la personne. (Ex: Démocraties modernes, bureaucratie).
    \item \textbf{Noter les formes mixtes :} Dans la réalité, les régimes combinent souvent ces types (ex: Président élu [légal] qui cultive un culte de la personnalité [charismatique] et s'appuie sur des notables traditionnels [traditionnel]).
\end{enumerate}
\cle{Rappel} : La légitimité = croyance en la validité de l'ordre de domination (croyance subjective des dominés). La légalité = conformité aux lois en vigueur. Un pouvoir peut être légal sans être légitime (putschiste qui fait voter une constitution) et vice-versa.
\end{methode}

\begin{exoresolu}[Typologie des justifications du pouvoir au Cameroun et ailleurs]
\textbf{Énoncé :} Pour chacune des situations suivantes, identifiez le(s) type(s) de domination légitime (Weber) mobilisé(s) et justifiez votre réponse.
\begin{enumerate}
    \item \textit{Situation A :} Un chef de village est désigné selon la coutume matrilinéaire par le conseil des notables. Il porte les insignes royaux transmis par son oncle maternel. Les villageois lui obéissent « parce que la tradition l'exige ».
    \item \textit{Situation B :} Un général prend le pouvoir par un coup d'État. Il promet des élections dans deux ans. En attendant, il gouverne par ordonnances. Il déclare : « Je suis le seul capable de sauver la nation du chaos, le peuple me fait confiance. »
    \item \textit{Situation C :} Un Président est élu au suffrage universel direct pour 7 ans, renouvelable une fois, selon la Constitution de 1996 (révisée). Il nomme un Premier ministre issu de la majorité parlementaire. Les lois sont votées par l'Assemblée nationale.
    \item \textit{Situation D :} Un mouvement de libération nationale arrive au pouvoir après une guerre. Son leader historique devient Président à vie par acclamation populaire. La Constitution est suspendue, remplacée par des « directives du Guide ».
\end{enumerate}
\tcblower
\textbf{Solution détaillée :}
\begin{enumerate}
    \item \textbf{Situation A : Domination TRADITIONNELLE.}
    \textbf{Justification :} La source de l'autorité est la coutume (matrilinéarité), la transmission héréditaire (oncle $\rightarrow$ neveu), la permanence des rites (insignes). L'obéissance est due à la personne en tant que dépositaire de la tradition.
    \item \textbf{Situation B : Domination CHARISMATIQUE (en phase de routinisation vers le légal).}
    \textbf{Justification :} Le général s'appuie sur ses qualités personnelles exceptionnelles (« seul capable », « sauver la nation »), sur un lien affectif direct avec le peuple (« le peuple me fait confiance »), non sur des règles préétablies (il a violé la constitution par le coup d'État). La promesse d'élections marque une tentative de \emph{routinisation} vers le type légal pour pérenniser le pouvoir.
    \item \textbf{Situation C : Domination RATIONNELLE-LÉGALE (type pur).}
    \textbf{Justification :} L'accès au pouvoir suit une procédure juridique codifiée (Constitution, suffrage universel, limite de mandats). L'organisation administrative repose sur la séparation des pouvoirs (Président, PM, Parlement) et la compétence technique (bureaucratie). L'obéissance va à la \emph{fonction} (Président de la République) définie par la loi, non à l'homme.
    \item \textbf{Situation D : Domination CHARISMATIQUE routinisée en TRADITIONNELLE / AUTORITAIRE.}
    \textbf{Justification :} Origine charismatique (leader historique, guerre de libération, « Guide »). La suspension de la Constitution et le pouvoir à vie par acclamation marquent le refus de la rationalité légale. La transformation en « tradition » nouvelle (directives du Guide comme nouvelle coutume sacralisée) est une \emph{routinisation traditionaliste} du charisme.
\end{enumerate}
\end{exoresolu}

\begin{methode}[Classifier une forme de gouvernement (Critères juridiques)]
\textbf{Objectif :} Distinguer Régime parlementaire, Régime présidentiel, Régime semi-présidentiel (et dictature) à partir des relations Exécutif/Législatif.
\begin{enumerate}
    \item \textbf{Identifier le Chef de l'État et le Chef du Gouvernement :} Sont-ce la même personne ou deux personnes distinctes ?
    \item \textbf{Analyser l'origine de l'Exécutif :}
    \begin{itemize}
        \item Élu directement par le peuple $\rightarrow$ Présidentiel / Semi-présidentiel (pour le Président).
        \item Issu du Parlement (majorité) $\rightarrow$ Parlementaire (Premier ministre).
    \end{itemize}
    \item \textbf{Analyser la responsabilité politique :}
    \begin{itemize}
        \item Le Gouvernement est-il responsable devant le Parlement ? (Vote de censure, motion de défiance, question de confiance) $\rightarrow$ \textbf{OUI} = Parlementaire ou Semi-présidentiel. \textbf{NON} = Présidentiel.
        \item Le Parlement peut-il être dissous par l'Exécutif ? $\rightarrow$ \textbf{OUI} = Parlementaire ou Semi-présidentiel. \textbf{NON} = Présidentiel (mandats fixes).
    \end{itemize}
    \item \textbf{Synthétiser :}
    \begin{center}
    \begin{tabular}{|l|c|c|c|}\hline
    \textbf{Critère} & \textbf{Présidentiel} & \textbf{Parlementaire} & \textbf{Semi-présidentiel} \\ \hline
    Chef État / Chef Gov & \textbf{Même personne} & \textbf{Séparés} & \textbf{Séparés} \\ \hline
    Origine Exécutif & Peuple (Président) & Parlement (PM) & Peuple (Président) + Parlement (PM) \\ \hline
    Resp. Govt devant Parl. & \textbf{Non} & \textbf{Oui} & \textbf{Oui} (PM) \\ \hline
    Dissolution Parlement & \textbf{Non} & \textbf{Oui} (par Chef État) & \textbf{Oui} (par Président) \\ \hline
    Exemple type & USA & Royaume-Uni, Allemagne & France (V$^e$ République), Cameroun \\ \hline
    \end{tabular}
    \end{center}
    \item \textbf{Attention :} La \textbf{dictature} / régime autoritaire se définit par l'absence de séparation des pouvoirs réelle, l'absence d'élections libres, la confusion des pouvoirs au profit d'un seul. Ce n'est pas une « forme de gouvernement » au sens constitutionnel démocratique.
\end{enumerate}
\end{methode}

\begin{exoresolu}[Classification de régimes politiques]
\textbf{Énoncé :} Classer les régimes politiques décrits ci-dessous en choisissant parmi : \emph{Régime présidentiel, Régime parlementaire, Régime semi-présidentiel, Dictature}. Justifiez en citant les critères décisifs.
\begin{enumerate}
    \item \textit{Régime X :} Le Président de la République est élu au suffrage universel direct pour 5 ans. Il nomme le Premier ministre. Le Gouvernement est responsable devant l'Assemblée nationale (vote de censure possible). Le Président peut dissoudre l'Assemblée. Le Président préside le Conseil des ministres.
    \item \textit{Régime Y :} Le Roi est Chef de l'État héréditaire. Le Premier ministre est le chef du parti majoritaire au Parlement. Le Gouvernement doit obtenir la confiance du Parlement. Le Roi ne peut pas dissoudre le Parlement sans contreseing du PM. Le Roi « règne mais ne gouverne pas ».
    \item \textit{Régime Z :} Le Président est élu au suffrage universel pour 4 ans. Il est à la fois Chef de l'État et Chef du Gouvernement. Il nomme les ministres (secrétaires d'État) qui ne sont pas députés. Le Congrès (Parlement) vote les lois mais ne peut pas renverser le Gouvernement. Le Président ne peut pas dissoudre le Congrès. Mandats fixes.
    \item \textit{Régime W :} Le « Guide Suprême » arrive au pouvoir par un coup d'État. Il cumule les fonctions de Chef de l'État, Chef du Gouvernement, Chef de l'armée et Chef du parti unique. L'Assemblée est élue au scrutin uninominal majoritaire à un tour avec liste unique du parti au pouvoir. La Constitution accorde l'immunité totale au Guide. Pas de séparation des pouvoirs.
\end{enumerate}
\tcblower
\textbf{Solution détaillée :}
\begin{enumerate}
    \item \textbf{Régime X : SEMI-PRÉSIDENTIEL (Type français / camerounais).}
    \textbf{Critères décisifs :} Double tête de l'exécutif (Président élu au suffrage universel + PM responsable devant Parlement). \textbf{Double responsabilité :} PM responsable devant Parlement (critère parlementaire) MAIS Président élu par le peuple et disposant de la dissolution (critère présidentiel). C'est l'archétype du semi-présidentiel (Duverger).
    \item \textbf{Régime Y : PARLEMENTAIRE (Type monarchie parlementaire / Allemagne).}
    \textbf{Critères décisifs :} Séparation stricte Chef de l'État (Roi, symbolique, non élu) / Chef du Gouvernement (PM, politique, issu du Parlement). Responsabilité du Gouvernement \textbf{uniquement} devant le Parlement. Dissolution possible mais contrôlée (contreseing). Le pouvoir effectif est entre les mains de la majorité parlementaire.
    \item \textbf{Régime Z : PRÉSIDENTIEL (Type américain).}
    \textbf{Critères décisifs :} \textbf{Fusion} Chef de l'État / Chef du Gouvernement (Président). Élection directe du Président. \textbf{Séparation rigide des pouvoirs :} Pas de responsabilité politique du Gouvernement devant le Congrès (pas de motion de censure). Pas de dissolution du Congrès par le Président. Mandats fixes (4 ans). Checks and balances (freins et contrepoids) mais pas de responsabilité politique.
    \item \textbf{Régime W : DICTATURE / RÉGIME AUTORITAIRE.}
    \textbf{Critères décisifs :} Origine illégitime (coup d'État). \textbf{Confusion totale des pouvoirs} (exécutif, législatif via parti unique, judiciaire via immunité, militaire). Absence d'élections libres et pluralistes (liste unique). Non-responsabilité du dirigeant. Ne rentre dans aucune classification constitutionnelle démocratique.
\end{enumerate}
\end{exoresolu}

\begin{methode}[Analyser le rôle et la finalité de l'État (Service public, Monopole de la violence, Bien commun)]
\textbf{Objectif :} Évaluer une action publique au regard des fonctions régaliennes et sociales de l'État.
\begin{enumerate}
    \item \textbf{Identifier la fonction mobilisée :}
    \begin{itemize}
        \item \textbf{Fonction régalienne (de souveraineté) :} Justice, Police, Armée, Diplomatie, Monnaie. \emph{Monopole de la violence physique légitime} (Weber). Indispensable à l'existence de l'État.
        \item \textbf{Fonction sociale / de service public :} Éducation, Santé, Transports, Énergie, Eau, Protection sociale. \emph{Continuité, Égalité, Mutabilité, Adaptabilité} (principes du service public).
        \item \textbf{Fonction économique :} Régulation, Planification, Propriété publique, Redistribution fiscale.
    \end{itemize}
    \item \textbf{Vérifier les principes du service public (si applicable) :}
    \begin{itemize}
        \item \textbf{Continuité :} Service assuré sans interruption.
        \item \textbf{Égalité :} Accès pour tous sans discrimination (tarifs, accueil).
        \item \textbf{Mutabilité :} Adaptation aux besoins nouveaux (numérisation, etc.).
    \end{itemize}
    \item \textbf{Évaluer la finalité :} L'action vise-t-elle l'\textbf{intérêt général} (bien commun) ou un \textbf{intérêt particulier} (clientélisme, corruption, clan) ?
    \item \textbf{Confronter aux limites :} État minimal (libéralisme/Nozick) vs État providence (social-démocratie/Rawls) vs État totalitaire (contrôle total). L'État doit-il \emph{tout} faire ? (Subsidiarité).
\end{enumerate}
\cle{Concept clé} : La \textbf{finalité de l'État} est la réalisation du \textbf{Bien Commun} (Aristote, Thomas d'Aquin, Vatican II : « ensemble des conditions de la vie sociale qui permettent aux groupes et à chacun de leurs membres de parvenir à leur perfection plus pleinement et plus aisément »).
\end{methode}

\begin{exoresolu}[Évaluation d'une politique publique : « Opération École Pour Tous »]
\textbf{Énoncé :} L'État du « Pays M » lance l'opération « École Pour Tous » : construction de 500 écoles primaires en zone rurale, recrutement de 2000 instituteurs contractuels payés par l'État, gratuité des manuels scolaires. Financement : augmentation de 2\% de la TVA (taxe sur la valeur ajoutée) et emprunt extérieur. Un an après, un rapport d'ONG révèle : 10\% des écoles non achevées (détournement de fonds), instituteurs non payés depuis 6 mois (grèves), manuels introuvables dans 30\% des écoles (corruption à la douane).
\textbf{Questions :}
\begin{enumerate}
    \item Quelle fonction de l'État est visée ? Quel principe de service public est principalement mis en avant ?
    \item Analysez les dysfonctionnements à la lumière de la distinction \emph{Légalité / Légitimité} et du \emph{Monopole de la violence légitime}.
    \item L'État a-t-il rempli sa finalité (Bien Commun) ?
\end{enumerate}
\tcblower
\textbf{Solution détaillée :}
\begin{enumerate}
    \item \textbf{Fonction :} Fonction \textbf{sociale / de service public} (Éducation). C'est une mission de service public à caractère administratif (non marchande).
    \textbf{Principe mis en avant :} L'\textbf{Égalité} (accès gratuit pour tous, zones rurales défavorisées) et la \textbf{Continuité} (ouverture d'écoles là où il n'y en avait pas).
    \item \textbf{Analyse des dysfonctionnements :}
    \begin{itemize}
        \item \textbf{Légalité vs Légitimité :} La loi (légalité) a été votée (budget, recrutement). Mais la \textbf{réalité administrative} (non-paiement, détournements) détruit la \textbf{légitimité} de l'action étatique. Les citoyens ne croient plus à la parole de l'État. La « croyance en la validité de l'ordre » (Weber) s'effondre.
        \item \textbf{Monopole de la violence légitime :} L'État a le monopole de la contrainte (police, justice) pour punir les coupables (détournement, corruption). Si l'État \emph{ne poursuit pas} les responsables (impunité), il renonce à son monopole de la violence légitime au profit de réseaux parallèles (mafieux, claniques). La grève des instituteurs (violence symbolique/légitime ?) est une réaction à la défaillance de l'État employeur.
    \end{itemize}
    \item \textbf{Finalité (Bien Commun) :} \textbf{Échec partiel / relatif.}
    \begin{itemize}
        \item \textbf{Positif :} 90\% des écoles construites, 2000 postes créés, intention redistributive (TVA, dette) vers les plus pauvres. Volonté affichée de Bien Commun.
        \item \textbf{Négatif :} La corruption et la mauvaise gestion (détournement, impayés) transforment le service public en \textbf{prédation privée}. L'argent public (impôt des citoyens) sert des intérêts particuliers. Le Bien Commun est \emph{trahi} par les agents de l'État. L'État devient un « État prédateur » (Bayart) plutôt qu'un État serviteur.
    \end{itemize}
    \textbf{Conclusion :} L'État a la \emph{forme} juridique de l'action pour le Bien Commun, mais en perd la \emph{substance} morale et effective par défaut de \textbf{gouvernance} (reddition des comptes, État de droit).
\end{enumerate}
\end{exoresolu}

\begin{methode}[Structurer une dissertation philosophique sur l'État (Bac Tech)]
\textbf{Objectif :} Produire une introduction, un plan dialectique et une conclusion conformes aux attendus du Probatoire/Baccalauréat technique.
\begin{enumerate}
    \item \textbf{Analyser le sujet :} Définir les termes clés (ex: « État », « Mal nécessaire », « Légitimité », « Liberté »). Identifier la tension (ex: Contrainte vs Liberté, Ordre vs Justice).
    \item \textbf{Problématiser :} Transformer le sujet en question philosophique ouverte. \emph{Ex: « L'État, en garantissant l'ordre, ne risque-t-il pas d'aliéner la liberté qu'il est censé protéger ? »}
    \item \textbf{Construire le plan dialectique (Thèse / Antithèse / Synthèse) :}
    \begin{itemize}
        \item \textbf{Thèse (L'État est nécessaire / légitime) :} Hobbes (Léviathan : sortir de l'état de nature « guerre de tous contre tous »), Locke (protection vie/biens/liberté), Kelsen (ordre juridique). Arguments : Sécurité, Justice, Services publics, Cohésion nationale.
        \item \textbf{Antithèse (L'État est un danger / mal) :} Libertaire/Anarchistes (Bakounine, Proudhon : État = oppression, monopole violence), Marx/Engels (État = instrument de domination de classe), Libéraux radicaux (Nozick : État minimal seulement). Arguments : Bureaucratie, totalitarisme, atteinte aux libertés, inégalités structurelles.
        \item \textbf{Synthèse (Dépassement : L'État de Droit / Démocratie) :} Kant (État de droit : loi = raison universelle), Rousseau (Légitimité = Volonté générale / Contrat social), Rawls (Justice comme équité). L'État n'est pas un mal \emph{en soi}, mais un \emph{risque permanent}. La solution : \textbf{Séparation des pouvoirs}, \textbf{Droit}, \textbf{Contrôle citoyen}, \textbf{Subsidiarité}. L'État doit être \textbf{serviteur}, pas maître.
    \end{itemize}
    \item \textbf{Rédiger l'introduction (entonoir) :} Accroche $\rightarrow$ Définitions $\rightarrow$ Problématique $\rightarrow$ Annonce du plan.
    \item \textbf{Rédiger la conclusion :} Réponse synthétique à la problématique $\rightarrow$ Ouverture (ex: Mondialisation, État numérique, Gouvernance mondiale).
\end{enumerate}
\cle{Conseil Bac} : Citez \textbf{noms d'auteurs} + \textbf{concepts précis} (État de nature, Contrat social, Volonté générale, Monopole violence, État de droit). Exemples camerounais bienvenus (Décentralisation, CHAN 2021, Gestion Covid).
\end{methode}

\begin{exoresolu}[Dissertation : « L'État est-il un mal nécessaire ? »]
\textbf{Énoncé :} Traitez le sujet : « L'État est-il un mal nécessaire ? » (Introduction complète + Plan détaillé avec arguments/auteurs + Conclusion).
\tcblower
\textbf{Solution détaillée :}

\textbf{INTRODUCTION}
\textbf{[Accroche]} « L'État est le plus froid de tous les monstres froids » (Nietzsche, \emph{Ainsi parlait Zarathoustra}). Cette métaphore saisit l'ambivalence fondamentale de l'institution étatique : protectrice et menaçante.
\textbf{[Définitions]} L'État (au sens moderne, Weber) : « Communauté humaine qui, dans les limites d'un territoire donné, revendique avec succès le monopole de la violence physique légitime ». \emph{Mal} : ce qui porte atteinte au bien, à la liberté, au bonheur. \emph{Nécessaire} : ce qui ne peut pas ne pas être, condition indispensable.
\textbf{[Problématique]} Si l'État est indispensable pour sortir du chaos (Hobbes), sa puissance coercitive ne fait-elle pas de lui une menace permanente pour la liberté des citoyens ? En d'autres termes, \textbf{la contrainte étatique est-elle le prix inévitable de la sécurité, ou peut-on concevoir un ordre politique sans domination ?}
\textbf{[Annonce du plan]} Nous montrerons d'abord que l'État est une \textbf{nécessité vitale} pour garantir l'ordre et les droits (Thèse). Ensuite, que sa nature même en fait un \textbf{mal structurel}, porteur de risques totalitaires (Antithèse). Enfin, que la philosophie politique moderne tente de le \textbf{civiliser par le Droit et la Démocratie} pour le rendre supportable, non plus « mal nécessaire » mais « bien commun » (Synthèse).

\textbf{DÉVELOPPEMENT}

\textbf{I. THÈSE : L'État, condition nécessaire de la vie commune (Le rempart contre le chaos).}
\begin{itemize}
    \item \textbf{Hobbes (\emph{Léviathan}) :} État de nature = « guerre de tous contre tous » (\emph{bellum omnium contra omnes}). La vie est « solitaire, misérable, nasty, brutish, and short ». Le \textbf{Contrat social} : aliéner sa liberté naturelle à un Souverain absolu (Léviathan) pour obtenir la \textbf{Sécurité}. L'État est \textbf{nécessaire} car sans lui, pas de vie possible.
    \item \textbf{Locke :} État de nature précaire (pas de juge impartial). Contrat pour protéger \textbf{Vie, Liberté, Propriété}. État = « Juge commun » + « Force publique ». Nécessaire pour l'impartialité de la justice.
    \item \textbf{Fonctions positives (Ex. Cameroun) :} Unité nationale (bilingue, diversité), Infrastructures (barrage de Lom Pangar, routes), Santé (Couverture Santé Universelle en cours), Éducation. Sans État : fragmentation, insécurité (Boko Haram, sécessionnisme NOSO).
\end{itemize}

\textbf{II. ANTITHÈSE : L'État, mal structurel (La machine à broyer l'individu).}
\begin{itemize}
    \item \textbf{Anarchistes (Bakounine, Proudhon) :} L'État = \textbf{« Légitimation de la violence »}. Centralisation, bureaucratie, armée, police = instruments de domination d'une minorité sur la majorité. « Là où est l'État, là n'est pas la liberté ». L'État \emph{est} le mal ; la société peut s'auto-organiser (fédéralisme libertaire, entraide).
    \item \textbf{Marx / Engels (\emph{L'Origine de la famille...}) :} L'État = « Comité pour gérer les affaires communes de la bourgeoisie ». Instrument de \textbf{domination de classe} ». Il disparaîtra (« s'éteindra ») dans le communisme (administration des choses, non gouvernement des hommes).
    \item \textbf{Libéralisme radical (Nozick, \emph{Anarchie, État, Utopie}) :} Seul l'\textbf{État minimal} (police, justice, défense) est justifié. Tout État-providence (redistribution) = travail forcé / vol (atteinte au droit de propriété de soi). L'État qui dépasse le minimum est un mal.
    \item \textbf{Risque totalitaire (Arendt, \emph{Les Origines du totalitarisme}) :} Banalité du mal, bureaucratie sans pensée, terreur. L'État moderne peut devenir « monstre froid » absolu.
\end{itemize}

\textbf{III. SYNTHÈSE : L'État de Droit démocratique, dépassion de l'alternative (Serviteur, pas Maître).}
\begin{itemize}
    \item \textbf{Kant / Kelsen :} L'État n'est pas une personne morale mystique, mais un \textbf{ordre juridique normatif}. La loi (légalité) lie \emph{tous}, y compris les gouvernants. \textbf{État de Droit (Rechtsstaat)} : séparation des pouvoirs, hiérarchie des normes, justiciabilité.
    \item \textbf{Rousseau (\emph{Du Contrat Social}) :} La légitimité ne vient pas de la force mais de la \textbf{Volonté Générale}. « Obéir à la loi qu'on s'est prescrite, c'est la liberté ». L'État \emph{légitime} (République) n'est pas un mal, il \emph{est} la liberté civile.
    \item \textbf{Rawls (\emph{Théorie de la justice}) :} Principes de justice (Libertés égales + Inégalités utiles aux plus défavorisés). L'État \textbf{juste} corrige les hasards sociaux (éducation, santé). Finalité = \textbf{Bien Commun}.
    \item \textbf{Mécanismes de contrôle :} Séparation des pouvoirs (Montesquieu), Contrôle juridictionnel (Conseil Constitutionnel), Société civile, Presse libre, Décentralisation (proximité). \emph{Subsidiarité} : l'État ne fait que ce que le niveau inférieur ne peut pas faire.
\end{itemize}

\textbf{CONCLUSION}
\textbf{[Réponse]} L'État n'est pas un « mal » \emph{par essence}, mais une \textbf{institution risquée}. Il devient « mal nécessaire » seulement lorsqu'il est \textbf{absolu, arbitraire, non contrôlé} (État policier, prédateur, totalitaire).
\textbf{[Synthèse]} La philosophie politique moderne (Kant, Rousseau, Rawls, Droit international) montre que la \textbf{démocratie constitutionnelle} (État de Droit + Souveraineté populaire + Droits de l'Homme) transforme la contrainte en \textbf{liberté partagée}.
\textbf{[Ouverture]} Aujourd'hui, la \textbf{mondialisation} (firmes transnationales, GAFAM, changement climatique) et le \textbf{numérique} (surveillance algorithmique, identité numérique) déplacent les frontières de l'État. La question devient : \textbf{Comment garantir la souveraineté populaire et les droits fondamentaux face à des puissances non étatiques qui échappent au contrôle démocratique ?} L'État-nation doit se réinventer (État régulateur, coopération supranationale) pour ne pas redevenir ce « monstre froid » ou disparaître.
\end{exoresolu}

\begin{methode}[Rédiger une réponse argumentée à une étude de cas (Cas pratique / Commentaire de texte)]
\textbf{Objectif :} Répondre à une question de type « Appliquer les notions à une situation concrète » (SAVOIR-FAIRE 1 \& 2).
\begin{enumerate}
    \item \textbf{Lire le document / la situation} : Identifier les acteurs, le conflit, les faits juridiques/politiques.
    \item \textbf{Mobiliser les notions du cours} (Boîte à outils) : Souveraineté, Légitimité, Légalité, Service public, Bien commun, Forme de gouvernement, Séparation des pouvoirs, État de droit.
    \item \textbf{Structurer la réponse (Méthode ARE : Affirmer - Raisonner - Exemplifier/Étayer) :}
    \begin{itemize}
        \item \textbf{Affirmer :} Réponse directe à la question (Oui/Non, C'est un cas de...).
        \item \textbf{Raisonner :} Expliciter \textbf{POURQUOI} en utilisant la définition précise de la notion. Citer l'auteur si pertinent (Weber, Hobbes, Kelsen, Montesquieu).
        \item \textbf{Étayer :} Lier la notion aux \textbf{faits précis} du document/de la situation (chiffres, citations, articles de loi).
    \end{itemize}
    \item \textbf{Conclure :} Une phrase qui résout le problème posé, ouvre sur la portée générale (ex: « Cela illustre la fragilité de l'État de droit quand... »).
    \item \textbf{Soigner la rédaction :} Phrases complètes, vocabulaire précis (pas de familier), connecteurs logiques (donc, cependant, par conséquent, en effet).
\end{enumerate}
\cle{Grille d'évaluation type Bac :} 1) Justesse des notions (4 pts) 2) Pertinence de l'application au cas (4 pts) 3) Rigueur du raisonnement (4 pts) 4) Qualité de la langue/structure (3 pts) 5) Conclusion/ouverture (2 pts) + Bonus originalité/exemple local (1-2 pts).
\end{methode}

\begin{exoresolu}[Cas pratique : Conflit foncier et autorité de l'État]
\textbf{Énoncé :} Dans le village de \emph{Nkolbisson}, un chef traditionnel (Chef de 3$^e$ degré) vend un terrain communal à un promoteur immobilier privé sans consulter le conseil villageois ni la mairie. Les jeunes du village bloquent le chantier, brûlent les engins. Le Sous-Préfet (représentant de l'État) fait intervenir la gendarmerie qui arrête 5 jeunes. Le Chef traditionnel crie à l'ingérence de l'État dans la coutume. Le promoteur porte plainte pour « trouble à l'ordre public » et « destruction de biens ».
\textbf{Questions :}
\begin{enumerate}
    \item Qualifiez le type de domination légitime (Weber) dont relève l'autorité du Chef traditionnel. Justifiez.
    \item Le Sous-Préfet a-t-il légitimement utilisé le « monopole de la violence physique légitime » ? Argumentez.
    \item Ce conflit révèle-t-il une défaillance de l'État de droit ? Expliquez en mobilisant la notion de \textbf{légalité} et de \textbf{hiérarchie des normes}.
\end{enumerate}
\tcblower
\textbf{Solution détaillée :}

\textbf{1. Domination du Chef traditionnel : TRADITIONNELLE.}
\textbf{Raisonnement :} Selon Weber, la domination traditionnelle repose sur la \textbf{croyance en la sainteté des traditions immémoriales} et la légitimité de ceux qui y sont appelés par la coutume (héritage, investiture coutumière).
\textbf{Application :} Le Chef est « Chef de 3$^e$ degré » (reconnu par l'administration), désigné selon la coutume locale. Son autorité est acceptée « parce que cela a toujours été ainsi ». \textbf{Nuance :} Sa décision de vente \emph{unilatérale} sans consulter le conseil villageois trahit la coutume (abus de pouvoir traditionnel), ce qui peut éroder sa légitimité traditionnelle aux yeux des jeunes.

\textbf{2. Intervention du Sous-Préfet / Gendarmerie : OUI, usage légitime (formellement).}
\textbf{Raisonnement :} Weber définit l'État par le \textbf{monopole de la violence physique légitime} sur un territoire. Seul l'État (via ses agents : police, gendarmerie, armée) a le droit d'user de la contrainte physique pour faire respecter la loi.
\textbf{Application :}
\begin{itemize}
    \item \textbf{Légalité formelle :} Le blocage de chantier et la destruction de biens (engins) constituent des \textbf{troubles à l'ordre public} (atteinte à la propriété privée, entrave à la liberté du travail). Le Code pénal et la loi sur la décentralisation donnent au Sous-Préfet (représentant de l'État dans l'arrondissement) le devoir de maintenir l'ordre public.
    \item \textbf{Proportionnalité (condition de légitimité matérielle) :} L'arrestation de 5 jeunes doit être \textbf{proportionnée} aux faits (flagrance, nécessité). Si la gendarmerie a agi dans le cadre de la loi (procès-verbal, droits de la défense respectés), l'usage de la force est \textbf{légitime}.
    \item \textbf{Point de vigilance :} Si l'arrestation est arbitraire (pas de flagrance, torture, détention au-delà de la garde à vue sans procureur), l'État bascule dans l'illégalité/illégitimité (violence d'État, pas violence \emph{légitime}).
\end{itemize}

\textbf{3. Défaillance de l'État de droit ? OUI, à plusieurs niveaux.}
\textbf{Raisonnement :} L'État de droit (Kelsen, Constitution) implique : \textbf{Hiérarchie des normes} (Constitution $>$ Loi $>$ Règlement $>$ Coutume), \textbf{Séparation des pouvoirs}, \textbf{Protection des droits fondamentaux}, \textbf{Imputabilité de l'administration}.
\textbf{Application :}
\begin{itemize}
    \item \textbf{Conflit de normes (Hiérarchie) :} La \textbf{Constitution} (Préambule, Art. 1er) et la \textbf{Loi n°74/1 sur l'organisation des chefferies} placent le Chef traditionnel \textbf{sous l'autorité de l'administration} (Sous-Préfet). La coutume ne peut primer sur la loi républicaine (propriété foncière, domaine public/privé de l'État, code foncier). Le Chef a agi \emph{ultra vires} (au-delà de ses pouvoirs) en vendant un terrain communal (domaine public ?) sans procédure légale (décret, enquête publique).
    \item \textbf{Défaillance de la Légalité administrative :} L'administration (Mairie, Services du cadastre, Sous-Préfecture) a-t-elle \textbf{contrôlé} la régularité de la vente \emph{avant} le conflit ? L'absence de contrôle préalable (délivrance de titre foncier, avis de non-objection) est une défaillance de l'État de droit (insécurité juridique).
    \item \textbf{Accès à la Justice :} Les jeunes ont recours à la \textbf{violence privée} (autodéfense, émeute) car ils estiment (à tort ou à raison) que la justice étatique est inaccessible, lente, ou corrompue (promoteur « intouchable »). C'est l'indice d'un \textbf{déficit de légitimité de l'institution judiciaire}.
    \item \textbf{Bien Commun vs Intérêt Privé :} La vente profite à un promoteur privé et au Chef (intérêt particulier) au détriment de la communauté (terrain communal = bien commun). L'État (Sous-Préfet) protège l'acquéreur (légalité formelle du titre ?) au lieu de protéger le bien commun (annulation de la vente illicite).
\end{itemize}
\textbf{Conclusion :} Ce cas illustre la \textbf{tension entre légalité formelle} (État protège l'ordre public et la propriété privée issue d'un acte peut-être irrégulier) \textbf{et légitimité sociale} (population défend le bien commun coutumier). Il révèle un \textbf{État de droit incomplet} : la loi existe mais n'est pas appliquée préventivement (contrôle de légalité), la justice n'est pas perçue comme accessible/équitable, et l'administration peine à articuler droit moderne et droit coutumier (pluralisme juridique non géré).
\end{exoresolu}

\begin{methode}[Comparer des conceptions philosophiques de l'État (Tableau synoptique)]
\textbf{Objectif :} Maîtriser les repères auteurs/thèses pour le commentaire de texte ou la dissertation.
\begin{enumerate}
    \item \textbf{Construire un tableau mental (ou sur brouillon) avec les colonnes :} Auteur / Conception de l'État / Fondement (Origine) / Finalité / Limites / Citation clé.
    \item \textbf{Remplir pour les 4 figures incontournables du programme :}
    \begin{center}
    \begin{tabularx}{\linewidth}{|l|X|X|X|X|}\hline
    \textbf{Auteur} & \textbf{Conception} & \textbf{Fondement} & \textbf{Finalité} & \textbf{Citation / Concept clé} \\ \hline
    \textbf{Hobbes} & \textbf{État absolu (Léviathan)} & Contrat de soumission (Peur mort violente) & \textbf{Sécurité / Paix civile} (Ordre) & « L'État est ce mortel dieu » ; Monopole violence \\ \hline
    \textbf{Locke} & \textbf{État libéral (Gouvernement de confiance)} & Contrat de confiance (Droits naturels) & \textbf{Protection Vie, Liberté, Propriété} & Droit de résistance si tyrannie ; Séparation pouvoirs \\ \hline
    \textbf{Rousseau} & \textbf{État souverain (République)} & Contrat d'association (Volonté générale) & \textbf{Liberté civile / Bien commun} & « Obéir à la loi qu'on s'est prescrite, c'est la liberté » ; Souveraineté inaliénable \\ \hline
    \textbf{Marx/Engels} & \textbf{État instrument de classe / À abolir} & Structure économique (Infrastructure) & \textbf{Domination de classe (Bourgeoisie)} & « Comité affaires bourgeoisie » ; Dépérissement (Communisme) \\ \hline
    \end{tabularx}
    \end{center}
    \item \textbf{Ajouter les modernes (Bac Tech) :} \textbf{Weber} (Définition sociologique : Monopole violence + Territoire), \textbf{Kelsen} (Normativisme : État = Ordre juridique), \textbf{Rawls} (État juste : Justice comme équité, Voile d'ignorance).
    \item \textbf{Savoir opposer :} Contractualistes (Hobbes/Locke/Rousseau) vs Critiques (Marx/Anarchistes/Libéraux radicaux). Naturel vs Historique. Fin vs Moyen.
\end{enumerate}
\end{methode}

\begin{exoresolu}[Commentaire guidé : Extrait de « Du Contrat Social » (Rousseau), Livre I, Chap. 6]
\textbf{Énoncé :} « Trouver une forme d'association qui défende et protège de toute la force commune la personne et les biens de chaque associé, et par laquelle chacun, s'unissant à tous, n'obéisse pourtant qu'à lui-même et reste aussi libre qu'auparavant. Telle est la difficulté fondamentale dont le contrat social donne la solution. »
\textbf{Questions :}
\begin{enumerate}
    \item Expliquez la « difficulté fondamentale » identifiée par Rousseau.
    \item En quoi le « contrat social » résout-il cette difficulté ? (Notions : Aliénation totale, Volonté générale, Souveraineté).
    \item Cette conception permet-elle d'éviter le « mal nécessaire » (Hobbes) ou la « domination de classe » (Marx) ? Argumentez.
\end{enumerate}
\tcblower
\textbf{Solution détaillée :}

\textbf{1. La difficulté fondamentale : Le paradoxe de l'association politique.}
Rousseau pose l'équation impossible en apparence : \textbf{Comment s'associer pour se protéger (force commune) sans s'asservir (perdre sa liberté naturelle) ?}
\begin{itemize}
    \item \textbf{Force commune :} Nécessaire car l'état de nature devient précaire (apparition propriété, inégalités, conflits). L'individu seul ne peut plus faire face.
    \item \textbf{Obéir qu'à soi-même :} Condition de la liberté (autonomie = se donner sa propre loi / hétéronomie = subir loi d'autrui).
    \item \textbf{Reste aussi libre :} La liberté naturelle (indépendance physique) doit se muer en liberté civile (obéissance à la loi rationnelle), sans perte nette.
\end{itemize}
C'est le problème de la \textbf{légitimité du pouvoir} : comment un pouvoir \emph{extérieur} (l'État) peut-il être \emph{intérieur} à ma volonté ?

\textbf{2. La solution du Contrat Social : L'aliénation totale à la communauté.}
\begin{itemize}
    \item \textbf{Aliénation totale :} « Chacun de nous met en commun sa personne et toute sa puissance sous la suprême direction de la volonté générale. » Je donne \emph{tout} (biens, liberté, vie) à \emph{tous} (la communauté), pas à un maître (Roi).
    \item \textbf{Réciprocité / Égalité :} Comme tous font de même, je récupère en droit ce que je donne. Je suis \textbf{auteur} de la loi (Souverain) et \textbf{sujet} de la loi.
    \item \textbf{Volonté Générale $\neq$ Volonté de tous :} La Volonté Générale vise l'\textbf{intérêt commun} (Bien Commun), pas la somme des intérêts privés. Elle est \textbf{toujours droite} (tend au bien commun), mais le peuple peut se tromper sur les moyens (besoin de Législateur).
    \item \textbf{Souveraineté inaliénable / indivisible :} Le Souverain, c'est le Peuple réuni. Il ne peut être représenté (critique du parlementarisme anglais). La loi émane du Souverain ; le Gouvernement (Prince) n'est qu'exécuteur.
\end{itemize}
$\rightarrow$ Je n'obéis qu'à \emph{moi-même} (en tant que membre du Souverain) $\rightarrow$ \textbf{Liberté = Obéissance à la loi qu'on s'est prescrite.}

\textbf{3. Dépassement de Hobbes et Marx ?}
\begin{itemize}
    \item \textbf{Vs Hobbes (Léviathan) :} \textbf{OUI, radicalement.}
    \begin{itemize}
        \item Hobbes : Aliénation \emph{à un tiers} (Souverain absolu, personne physique ou assemblée). Je perds ma liberté pour la sécurité. Souverain \emph{au-dessus} des lois.
        \item Rousseau : Aliénation \emph{au tout} (Communauté). Je gagne la liberté morale/civile. Souverain \emph{est} la loi. Pas de « mal nécessaire », mais \textbf{naissance de la moralité politique}.
    \end{itemize}
    \item \textbf{Vs Marx (Instrument de classe) :} \textbf{PARTIELLEMENT / EN THÉORIE.}
    \begin{itemize}
        \item Rousseau \emph{idéalise} : Il suppose des citoyens vertueux, égaux en droits, transparents à la Volonté Générale. Il ignore les \textbf{conditions matérielles} (économie, classes sociales) que Marx analyse.
        \item Marx répondrait : Dans une société divisée en classes, la « Volonté Générale » masque la volonté de la classe dominante (bourgeoisie). L'État rousseauiste (République bourgeoise) \emph{est} l'instrument de la bourgeoisie (Code civil, propriété sacralisée).
        \item \textbf{Nuance Rousseau :} Il dénonce la propriété privée comme source des maux (\emph{Discours sur l'inégalité}) et veut une économie encadrée (censure, religion civile). Mais il ne pense pas la \textbf{structure de classe} comme moteur de l'histoire.
    \end{itemize}
\end{itemize}
\textbf{Conclusion :} Rousseau offre la \textbf{théorie normative la plus haute de la légitimité démocratique} (Souveraineté populaire, État de droit). Mais sans transformation des \textbf{rapports sociaux réels} (économie, éducation, inégalités), le « Contrat » risque de rester une fiction juridique masquant une domination réelle (critique marxiste) ou de s'effondrer dans le despotisme si le Peuple ne se réunit pas (critique libérale / réaliste).
\end{exoresolu}

\begin{attention}[Les 5 erreurs qui coûtent des points au Bac (Probatoire / Baccalauréat Technique)]
\begin{enumerate}
    \item \textbf{Confondre \emph{Légalité} et \emph{Légitimité}.}
    \emph{Erreur :} « Le Président est légitime car il a été élu. » $\rightarrow$ \textbf{Correction :} L'élection confère la \textbf{légalité} (conformité à la procédure). La \textbf{légitimité} est la \emph{croyance} des gouvernés en le droit de commander (Weber). Un putschiste peut faire voter une loi (légalité) sans être légitime. Un chef coutumier est légitime (tradition) sans base légale formelle.
    \item \textbf{Réduire l'État au Gouvernement / au Président.}
    \emph{Erreur :} « L'État a décidé d'augmenter les prix. » $\rightarrow$ \textbf{Correction :} L'État est une \textbf{personne morale permanente} (institutions, administration, territoire, population). Le Gouvernement/President est un \textbf{organe} temporaire de l'État. Distinguer \emph{L'État} (structure) et \emph{Les gouvernants} (personnes).
    \item \textbf{Oublier le \emph{Monopole de la violence légitime} (Weber) comme critère \emph{définitif} de l'État.}
    \emph{Erreur :} Définir l'État seulement par « population + territoire + gouvernement ». $\rightarrow$ \textbf{Correction :} Une entreprise multinationale a population, territoire (usines), gouvernement (PDG). Ce qui fait l'État, c'est la \textbf{revendication réussie du monopole de la violence physique légitime} sur un territoire (police, justice, armée, prison). C'est LA définition de référence au programme.
    \item \textbf{Mélanger \emph{Forme de l'État} (Unitaire/Fédéral) et \emph{Forme de Gouvernement} (Présidentiel/Parlementaire).}
    \emph{Erreur :} « La France est un État fédéral car il y a un Président et un Premier ministre. » $\rightarrow$ \textbf{Correction :} \textbf{Forme de l'État} = Organisation territoriale du pouvoir (Unitaire : centre fort / Fédéral : partage souveraineté États fédérés). \textbf{Forme de Gouvernement} = Relations Exécutif/Législatif (Présidentiel, Parlementaire, Semi-présidentiel). La France = État Unitaire + Régime Semi-présidentiel. Le Cameroun = État Unitaire décentralisé + Régime Semi-présidentiel.
    \item \textbf{Rédiger une dissertation sans \emph{Problématique} explicite ni \emph{Plan annoncé}.}
    \emph{Erreur :} Commencer direct : « Hobbes dit que... Ensuite Locke dit... » $\rightarrow$ \textbf{Correction :} L'introduction \textbf{doit} contenir : 1) Accroche, 2) Définitions des termes du sujet, 3) \textbf{Problématique (sous forme de question)}, 4) \textbf{Annonce du plan (Thèse / Antithèse / Synthèse)}. Sans cela : note plafonnée (souvent $< 8/20$) car « hors sujet » ou « absence de démarche philosophique ».
\end{enumerate}
\end{attention}

\newpage

\coursec{Exercices}

\courssub{Je vérifie mes connaissances}

\begin{exercice}[QCM : notions fondamentales]
\textbf{Choisissez la bonne réponse parmi les propositions.}
\begin{enumerate}
    \item L'ensemble des institutions qui détiennent le monopole de la violence physique légitime sur un territoire donné définit, selon Max Weber :
    \begin{itemize}
        \item[A.] Le Gouvernement
        \item[B.] La Nation
        \item[C.] \cle{L'État}
        \item[D.] La Société civile
    \end{itemize}
    \item Parmi les éléments constitutifs de l'État, celui qui désigne le pouvoir suprême, absolu et inaliénable est :
    \begin{itemize}
        \item[A.] Le territoire
        \item[B.] La population
        \item[C.] \cle{La souveraineté}
        \item[D.] L'administration
    \end{itemize}
    \item Pour \cle{John Locke}, le fondement de l'autorité politique réside dans :
    \begin{itemize}
        \item[A.] La force brute
        \item[B.] \cle{Le consentement des gouvernés (contrat social)}
        \item[C.] Le droit divin
        \item[D.] La tradition historique
    \end{itemize}
    \item La classification des régimes selon \cle{Aristote} repose sur le critère du :
    \begin{itemize}
        \item[A.] Mode de succession au pouvoir
        \item[B.] Nombre de gouvernants et finalité du pouvoir (intérêt général vs intérêt particulier)
        \item[C.] Degré de centralisation administrative
        \item[D.] Existence d'une constitution écrite
    \end{itemize}
\end{enumerate}
\end{exercice}

\begin{exercice}[Vrai ou Faux ? Justifiez brièvement]
\textbf{Indiquez si l'affirmation est vraie ou fausse, puis justifiez votre réponse en une phrase en mobilisant une notion du cours.}
\begin{enumerate}
    \item « L'État et le Gouvernement sont des synonymes : le Président de la République est l'État. »
    \item « Un régime où le pouvoir est exercé par un seul pour l'intérêt de tous est une tyrannie selon Aristote. »
    \item « La \cle{séparation des pouvoirs} (Montesquieu) implique que les trois pouvoirs (législatif, exécutif, judiciaire) soient totalement étanches et n'aient aucune relation entre eux. »
    \item « La finalité ultime de l'État, selon la tradition contractualiste (Hobbes, Locke, Rousseau), est la sécurité et la liberté des citoyens. »
\end{enumerate}
\end{exercice}

\begin{exercice}[Définitions attendues au Bac]
\textbf{Donnez une définition précise et philosophique (2 à 3 lignes) de chacun des concepts suivants :}
\begin{enumerate}
    \item \cle{L'État} (au sens politique et juridique).
    \item \cle{La légitimité} (distinction fait/droit, Weber).
    \item \cle{L'État de droit}.
    \item \cle{La volonté générale} (Rousseau).
\end{enumerate}
\end{exercice}

\begin{exercice}[Correspondance : Auteurs et Thèses]
\textbf{Reliez chaque auteur à sa thèse principale sur le fondement ou la forme du pouvoir.}
\begin{center}
\begin{tabularx}{\linewidth}{|l|X|l|}\hline
\textbf{Auteur} & \textbf{Thèse principale} & \textbf{N°} \\ \hline
Thomas Hobbes & 1. Le pouvoir naît d'un contrat pour sortir de l'état de nature guerrier ; souveraineté absolue. & \\ \hline
John Locke & 2. Séparation des pouvoirs (législatif, exécutif, judiciaire) pour éviter la tyrannie. & \\ \hline
Jean-Jacques Rousseau & 3. Le pouvoir légitime repose sur le consentement ; droit de résistance si violation des droits naturels. & \\ \hline
Montesquieu & 4. Souveraineté du peuple exercée par la volonté générale ; liberté = obéissance à la loi qu'on s'est prescrite. & \\ \hline
Aristote & 5. Classification selon le nombre (un, quelques-uns, tous) et la finalité (bien commun vs intérêt privé). & \\ \hline
\end{tabularx}
\end{center}
\end{exercice}

\courssub{Je m'entraîne}

\begin{exercice}[Analyse de citation : Fondement de l'obéissance]
\textbf{Lisez la citation suivante et répondez aux questions :}
\begin{quote}
« Le plus fort n'est jamais assez fort pour être toujours le maître, s'il ne transforme sa force en droit et l'obéissance en devoir. » \\
\emph{Jean-Jacques Rousseau, \textit{Du Contrat social}, Livre I, Chap. 3.}
\end{quote}
\begin{enumerate}
    \item Identifiez la thèse défendue par Rousseau dans cet extrait.
    \item Expliquez la distinction opérée entre « force » et « droit », « obéissance » et « devoir ».
    \item En quoi cette citation critique-t-elle la thèse de la « force » comme fondement unique de l'État (thèse du fait accompli) ?
    \item Donnez un exemple historique ou contemporain (Cameroun ou ailleurs) où un pouvoir issu de la force a dû chercher une \cle{légitimité} juridique (élections, constitution).
\end{enumerate}
\end{exercice}

\begin{exercice}[Étude de cas : La Constitution camerounaise de 1996]
\textbf{L'article 2 de la Constitution de la République du Cameroun (Loi n°96-06 du 18 janvier 1996) dispose :}
\begin{quote}
« La souveraineté nationale appartient au peuple qui l'exerce soit par voie de référendum, soit par l'intermédiaire de ses représentants élus. [...] L'État est unitaire et décentralisé. [...] Le suffrage est universel, égal et secret. [...] Les partis politiques concourent à l'expression du suffrage. »
\end{quote}
\begin{enumerate}
    \item Identifiez dans ce texte les \textbf{trois principes fondamentaux} de l'organisation politique de l'État camerounais.
    \item Montrez en quoi ce texte ancre le régime camerounais dans la tradition \cle{contractualiste} (Rousseau, Locke).
    \item Expliquez la portée de la mention « L'État est unitaire et décentralisé » sur la forme de l'État (distinction État unitaire / État fédéral).
    \item Quel rôle la Constitution assigne-t-elle aux \cle{partis politiques} ? Reliez cela à la notion de \cle{démocratie pluraliste}.
\end{enumerate}
\end{exercice}

\begin{exercice}[Piège classique : « L'État, c'est moi »]
\textbf{Lors d'un meeting électoral, un candidat déclare : « Si je suis élu, je serai l'État. Je déciderai seul de ce qui est bon pour le pays. »}
\textbf{En vous appuyant sur la distinction \cle{État / Gouvernement} et la théorie de la \cle{séparation des pouvoirs} (Montesquieu), discutez cette affirmation en structurant votre réponse en trois arguments numérotés.}
\end{exercice}

\begin{exercice}[Classification des régimes : Application comparative]
\textbf{Complétez le tableau suivant en classant les régimes proposés selon la typologie d'Aristote (critère : nombre + finalité) et selon la typologie moderne (Présidentiel / Parlementaire / Mixte). Justifiez chaque case en une phrase.}
\begin{center}
\begin{tabularx}{\linewidth}{|l|X|X|X|}\hline
\textbf{Régime réel} & \textbf{Aristote : Forme (Nombre/Finalité)} & \textbf{Type moderne (Exécutif/Législatif)} & \textbf{Justification brève} \\ \hline
Royaume-Uni (Monarchie parlementaire) & & & \\ \hline
États-Unis (République présidentielle) & & & \\ \hline
Cameroun (République à régime présidentiel) & & & \\ \hline
France V\textsuperscript{e} République (Régime semi-présidentiel) & & & \\ \hline
\end{tabularx}
\end{center}
\end{exercice}

\courssub{Je me prépare au Bac}

\begin{exercice}[Sujet de Dissertation Type Bac Tech]
\textbf{Sujet : « L'État est-il nécessaire à la liberté des citoyens ? »}
\textbf{Consignes :} Vous traiterez ce sujet en suivant la méthode de la dissertation philosophique (Introduction : accroche, définition des termes, problématique, annonce du plan ; Développement en 2 ou 3 parties argumentées avec exemples ; Conclusion : bilan, ouverture).
\begin{enumerate}
    \item \textbf{Problématisez} : Montrez la tension entre l'État comme garant de la liberté (loi, protection) et l'État comme menace potentielle (contrainte, surveillance, totalitarisme).
    \item \textbf{Argumentation attendue} (pistes) :
    \begin{itemize}
        \item Thèse : L'État comme condition de la liberté (Hobbes : sécurité ; Locke : protection des droits naturels ; Rousseau : liberté civile = obéissance à la loi qu'on s'est prescrite ; \cle{État de droit}).
        \item Antithèse : L'État comme limite ou danger pour la liberté (constrainte légale, risque de tyrannie, surveillance de masse, raison d'État).
        \item Synthèse : La liberté politique ne s'exerce que \emph{dans} et \emph{par} l'État de droit démocratique (séparation des pouvoirs, constitution, contrôle juridictionnel, participation citoyenne). Distinction liberté des Anciens / Modernes (Benjamin Constant).
    \end{itemize}
    \item \textbf{Exemples camerounais obligatoires} : Rôle de la Constitution de 1996, loi sur les libertés, fonctionnement du Conseil constitutionnel, décentralisation (loi 2019/024).
\end{enumerate}
\end{exercice}

\begin{exercice}[Explication de Texte : Rousseau, \textit{Du Contrat Social}]
\textbf{Texte :}
\begin{quote}
« [...] Pour que le pacte social ne soit pas un vain formulaire, il renferme tacitement cet engagement, qui seul peut donner de la force aux autres : que quiconque refusera d'obéir à la volonté générale y sera contraint par tout le corps ; ce qui ne signifie autre chose sinon qu'on le forcera d'être libre. Car c'est la condition qui, donnant chaque citoyen à la patrie, le garantit de toute dépendance personnelle ; c'est l'artifice qui a rendu légitimes les engagements civils, qui, sans être mutuels, nuls et sans effet, deviennent, par le pacte social, volontaires, légitimes et sacrés. » \\
\emph{Jean-Jacques Rousseau, \textit{Du Contrat social}, Livre I, Chap. 7.}
\end{quote}
\textbf{Questions :}
\begin{enumerate}
    \item \textbf{Dégagez la thèse} principale du texte en une phrase.
    \item \textbf{Identifiez et expliquez} les \textbf{trois arguments} majeurs qui la soutiennent (le pacte, la contrainte/liberté, la légitimité).
    \item \textbf{Expliquez la formule paradoxale} : « on le forcera d'être libre ». En quoi révèle-t-elle la spécificité de la \cle{volonté générale} par rapport à la volonté particulière ?
    \item \textbf{Portée et actualité} : En quoi cette conception de la loi comme expression de la volonté générale fonde-t-elle la \cle{légitimité} de l'État de droit moderne ? Illustrer par l'exemple de l'obligation vaccinale ou du code de la route au Cameroun.
\end{enumerate}
\end{exercice}

\begin{exercice}[Question de Réflexion Guidée / Analyse de Situation Complexe]
\textbf{Contexte :} Un débat oppose deux citoyens camerounais à propos de la gestion de la chose publique.
\begin{itemize}
    \item \textbf{M. A} affirme : « Le Président de la République a été élu avec une large majorité. Il incarne donc la volonté du peuple. Il n'a pas de comptes à rendre au Parlement ni à la Justice entre deux élections. La force de son mandat lui donne tous les pouvoirs. »
    \item \textbf{Mme B} répond : « La Constitution de 1996 institue une République \emph{démocratique} et \emph{unitaire}. Le Président est le chef de l'Exécutif, mais le pouvoir législatif appartient au Parlement et le pouvoir judiciaire aux tribunaux. Nul n'est au-dessus de la loi. »
\end{itemize}
\textbf{Travail demandé :}
\begin{enumerate}
    \item \textbf{Qualifiez} la conception du pouvoir défendue par M. A (référence auteur/théorie : souveraineté absolue, plébiscitaire, confusion État/Gouvernement).
    \item \textbf{Analysez} la réponse de Mme B en identifiant les \textbf{trois principes constitutionnels} qu'elle mobilise (Souveraineté populaire, Séparation des pouvoirs, Hiérarchie des normes / Légalité).
    \item \textbf{Discutez} : « La force du mandat électoral suffit-elle à fonder l'exercice illimité du pouvoir ? » Structurez votre réponse en confrontant \cle{Légitimité d'origine} (élection) et \cle{Légitimité d'exercice} (respect de la Constitution, contrôle des comptes, reddition des comptes, alternance).
    \item \textbf{Concluez} en définissant ce qu'est un \cle{Gouvernement responsable} dans un régime présidentiel (responsabilité politique devant le Parlement ? non, mais responsabilité constitutionnelle et morale).
\end{enumerate}
\end{exercice}

\newpage

\coursec{Corrigés des exercices}

\coursec{Exercices d'application et corrigés détaillés — L'État et le pouvoir}

\begin{corrige}[Exercice 1 — QCM : notions fondamentales]
\textbf{Question 1 : réponse C — L'État.}\\
C'est la définition célèbre de Max Weber (\emph{Le Savant et le politique}, 1919) : l'État est « une communauté humaine qui, dans les limites d'un territoire déterminé, revendique avec succès pour son propre compte le monopole de la violence physique légitime ». Le Gouvernement (A) n'est que l'organe de direction de l'État ; la Nation (B) est une communauté culturelle et historique ; la société civile (D) désigne l'ensemble des citoyens et associations hors sphère étatique.

\textbf{Question 2 : réponse C — La souveraineté.}\\
Selon Jean Bodin (\emph{Les Six Livres de la République}, 1576), la souveraineté est « la puissance absolue et perpétuelle d'une République ». Le territoire et la population sont des éléments matériels de l'État ; l'administration est un simple instrument d'exécution.

\textbf{Question 3 : réponse B — Le consentement des gouvernés.}\\
Pour Locke (\emph{Second Traité du gouvernement civil}, 1690), les hommes quittent l'état de nature par un contrat pour protéger leurs droits naturels (vie, liberté, propriété). L'autorité politique ne vaut donc que par le consentement des gouvernés, qui conservent un droit de résistance en cas de violation.

\textbf{Question 4 : réponse B — Nombre de gouvernants et finalité du pouvoir.}\\
Aristote (\emph{Les Politiques}) croise deux critères : \emph{qui gouverne ?} (un seul, quelques-uns, la multitude) et \emph{pour qui ?} (l'intérêt général : monarchie, aristocratie, république ; l'intérêt particulier : tyrannie, oligarchie, démagogie).
\end{corrige}

\begin{attention}[Points de vigilance — Exercice 1]
Ne pas confondre \cle{État} et \cle{Gouvernement} : c'est la confusion la plus fréquente au Bac. Retenir aussi que la souveraineté est à la fois \emph{absolue} (aucun pouvoir supérieur à l'intérieur) et \emph{inaliénable} (elle ne peut être cédée).
\end{attention}

\begin{corrige}[Exercice 2 — Vrai ou Faux ?]
\textbf{1. FAUX.} L'État est une institution permanente (territoire, population, souveraineté), tandis que le Gouvernement est un organe temporaire et révocable qui exerce le pouvoir au nom de l'État. Le Président incarne le Gouvernement, non l'État : les gouvernants passent, l'État demeure.

\textbf{2. FAUX.} Selon Aristote, le pouvoir d'un seul exercé pour l'intérêt de tous est une \cle{monarchie} (forme droite). La \cle{tyrannie} est sa forme dégénérée : le pouvoir d'un seul exercé pour son intérêt particulier.

\textbf{3. FAUX.} Pour Montesquieu (\emph{De l'Esprit des lois}, 1748), la séparation des pouvoirs exige leur distribution, mais aussi des mécanismes de contrôle réciproque (\emph{checks and balances}) : « le pouvoir arrête le pouvoir ». Une étanchéité totale paralyserait l'État ; c'est la collaboration contrôlée qui garantit la liberté.

\textbf{4. VRAI.} Pour Hobbes (sortir de la guerre de tous contre tous), Locke (protéger les droits naturels) et Rousseau (garantir la liberté civile par la volonté générale), le pacte social a toujours pour finalité la sécurité et la liberté des associés.
\end{corrige}

\begin{attention}[Points de vigilance — Exercice 2]
Dans un Vrai/Faux, la justification doit mobiliser explicitement une notion du cours (ici : État/Gouvernement, typologie d'Aristote, séparation des pouvoirs, contractualisme). Une réponse sans justification ne vaut aucun point.
\end{attention}

\begin{corrige}[Exercice 3 — Définitions attendues au Bac]
\textbf{1. L'État.} L'État est l'ensemble des institutions politiques et juridiques qui exercent, sur un territoire délimité et sur la population qui l'habite, une autorité souveraine caractérisée par le monopole de la violence physique légitime (Weber). Ses éléments constitutifs sont : un territoire, une population et une souveraineté.

\textbf{2. La légitimité.} La légitimité est le caractère d'un pouvoir conforme au droit et reconnu comme juste par les gouvernés. Elle se distingue du simple fait de la force : un pouvoir peut être \emph{en fait} (imposé par la violence) sans être \emph{en droit}. Weber distingue trois légitimités : traditionnelle, charismatique et rationnelle-légale (fondée sur la loi).

\textbf{3. L'État de droit.} L'État de droit est un État dont le pouvoir est lui-même soumis au droit : il existe une hiérarchie des normes dominée par la Constitution, les pouvoirs sont séparés et contrôlés, et les citoyens disposent de recours juridiques contre l'arbitraire de l'administration.

\textbf{4. La volonté générale.} Selon Rousseau (\emph{Du Contrat social}), la volonté générale est la volonté du corps politique visant l'intérêt commun ; elle n'est pas la somme des volontés particulières (volonté de tous), mais ce qui demeure lorsque sont écartés les intérêts privés. Elle est la source unique de la loi légitime.
\end{corrige}

\begin{attention}[Points de vigilance — Exercice 3]
Une définition « attendue au Bac » comporte : le genre proche (institution, caractère, principe), la différence spécifique (ce qui distingue le concept) et, si possible, l'auteur de référence. Éviter les définitions par l'exemple seul.
\end{attention}

\begin{corrige}[Exercice 4 — Correspondance : Auteurs et Thèses]
\begin{center}
\begin{tabularx}{\linewidth}{|l|X|}\hline
\textbf{Auteur} & \textbf{Thèse correspondante} \\ \hline
Thomas Hobbes & Thèse 1 : contrat pour sortir de l'état de nature guerrier ; souveraineté absolue (le \emph{Léviathan}). \\ \hline
John Locke & Thèse 3 : consentement des gouvernés et droit de résistance en cas de violation des droits naturels. \\ \hline
Jean-Jacques Rousseau & Thèse 4 : souveraineté populaire exercée par la volonté générale ; la liberté est l'obéissance à la loi qu'on s'est prescrite. \\ \hline
Montesquieu & Thèse 2 : séparation des pouvoirs législatif, exécutif et judiciaire pour prévenir la tyrannie. \\ \hline
Aristote & Thèse 5 : classification des régimes selon le nombre des gouvernants et la finalité du pouvoir. \\ \hline
\end{tabularx}
\end{center}

\textbf{Justifications :} Hobbes fonde l'État sur la peur de la mort violente dans l'état de nature (« l'homme est un loup pour l'homme ») ; Locke, plus optimiste, limite le pouvoir par les droits naturels ; Rousseau radicalise la souveraineté populaire ; Montesquieu pense les institutions comme équilibre ; Aristote inaugure la science politique comparative.
\end{corrige}

\begin{attention}[Point de vigilance — Exercice 4]
Ne pas attribuer à Hobbes la séparation des pouvoirs : il défend au contraire une souveraineté \emph{absolue et indivisible}. C'est Locke puis Montesquieu qui théorisent la limitation du pouvoir.
\end{attention}

\begin{corrige}[Exercice 5 — Analyse de citation : Fondement de l'obéissance]
\textbf{1. Thèse.} Rousseau affirme que la force ne peut jamais fonder durablement le pouvoir : seul un pouvoir transformé en droit, et une obéissance transformée en devoir, rendent l'autorité politique stable et légitime.

\textbf{2. Les distinctions.} La \emph{force} est un fait physique, brut et précaire ; le \emph{droit} est une norme reconnue qui fonde l'autorité en légitimité. De même, l'\emph{obéissance} à la force est une soumission contrainte, révocable dès que la force faiblit ; le \emph{devoir} est une obligation morale et juridique que le citoyen reconnaît librement. On obéit à la force par nécessité, on accomplit un devoir par adhésion.

\textbf{3. Critique du « fait accompli ».} Si la force fondait le droit, alors toute révolte victorieuse serait légitime : le droit changerait à chaque rapport de forces, ce qui détruirait toute idée de droit. La force produit un pouvoir de fait, jamais une autorité de droit ; elle engendre la révolte, non l'obligation.

\textbf{4. Exemple.} Au Cameroun, les régimes issus de l'indépendance ont organisé des élections et adopté des constitutions (1960, 1972, 1996) pour convertir un pouvoir de fait en autorité de droit. Ailleurs, les régimes nés de coups d'État (par exemple en Afrique de l'Ouest) organisent presque toujours ensuite des « transitions » électorales : preuve que la force seule ne suffit pas et cherche à se transformer en droit.
\end{corrige}

\begin{attention}[Points de vigilance — Exercice 5]
Bien articuler le couple fait/droit. Ne pas écrire que Rousseau refuse toute force : il dit qu'elle doit être \emph{transformée} en droit. La citation doit être citée et expliquée, non paraphrasée vaguement.
\end{attention}

\begin{corrige}[Exercice 6 — Étude de cas : La Constitution camerounaise de 1996]
\textbf{1. Les trois principes fondamentaux.}
\begin{itemize}
    \item La \cle{souveraineté nationale} (ou populaire) : le peuple est la source du pouvoir (art. 2).
    \item L'exercice \cle{démocratique} du pouvoir : référendum (démocratie directe) et représentants élus au suffrage universel, égal et secret (démocratie représentative) (art. 3).
    \item La forme de l'État : \cle{unitaire et décentralisé} (art. 1).
\end{itemize}

\textbf{2. Ancrage contractualiste.} Dire que « la souveraineté appartient au peuple » reprend la thèse de Rousseau : le peuple seul est souverain et la loi exprime la volonté générale. Le recours au suffrage et à la représentation traduit l'idée lockéenne du consentement des gouvernés : les gouvernants ne tiennent leur autorité que d'une délégation du peuple, exprimée périodiquement par les élections.

\textbf{3. « Unitaire et décentralisé ».} Un État \emph{unitaire} concentre la souveraineté au niveau national : il n'existe qu'une seule Constitution et un seul pouvoir central, contrairement à l'État \emph{fédéral} (États-Unis, Nigeria) où les États fédérés disposent de compétences propres garanties. La \emph{décentralisation} signifie que des collectivités territoriales (régions, communes — loi n°2019/024 du 24 décembre 2019) gèrent librement les affaires locales, mais par délégation de l'État central, sans partage de la souveraineté.

\textbf{4. Rôle des partis politiques.} Les partis « concourent à l'expression du suffrage » : ils structurent l'offre électorale et permettent la compétition pour le pouvoir. Cela fonde une \cle{démocratie pluraliste} : la diversité des opinions et des candidatures est reconnue comme condition de la liberté politique, par opposition au parti unique (rappel historique : le Cameroun est passé du parti unique au multipartisme en 1990).
\end{corrige}

\begin{attention}[Points de vigilance — Exercice 6]
Bien distinguer \emph{décentralisation} (délégation dans un État unitaire) et \emph{fédéralisme} (partage constitutionnel de la souveraineté). Citer les articles 1, 2, 3 et la loi de 2019 valorise la copie.
\end{attention}

\begin{corrige}[Exercice 7 — Piège classique : « L'État, c'est moi »]
\textbf{Argument 1 — Confusion État / Gouvernement.} Le candidat confond l'État, institution permanente et impersonnelle (territoire, population, souveraineté), avec le Gouvernement, organe temporaire chargé de diriger. Même élu, un président n'est qu'un \emph{dépositaire} du pouvoir : il l'exerce au nom du peuple et pour une durée limitée. Dire « je serai l'État » relève de la monarchie absolue de droit divin (« L'État, c'est moi », formule attribuée à Louis XIV), incompatible avec une République.

\textbf{Argument 2 — Négation de la séparation des pouvoirs.} « Je déciderai seul » concentre dans une seule main les pouvoirs législatif (faire les lois), exécutif (les appliquer) et potentiellement judiciaire (juger). Or Montesquieu a montré que « tout homme qui a du pouvoir est porté à en abuser » et que seule la distribution des pouvoirs garantit la liberté : « le pouvoir arrête le pouvoir ». La concentration des pouvoirs est la définition même du \cle{despotisme}.

\textbf{Argument 3 — Négation de la souveraineté populaire et de l'État de droit.} Dans une République démocratique (article 2 de la Constitution de 1996), la souveraineté appartient au peuple, non au gouvernant. L'élection est une \emph{délégation limitée} par la Constitution : le président élu reste soumis à la loi, au contrôle juridictionnel et à la reddition des comptes. Un pouvoir qui décide « seul de ce qui est bon » substitue la volonté particulière d'un homme à la volonté générale (Rousseau) : il devient illégitime, même s'il est élu.
\end{corrige}

\begin{attention}[Points de vigilance — Exercice 7]
Trois arguments numérotés étaient exigés : respecter la consigne formelle. Chaque argument doit mobiliser une notion précise (État/Gouvernement ; séparation des pouvoirs ; souveraineté populaire/État de droit) et un auteur.
\end{attention}

\begin{corrige}[Exercice 8 — Classification des régimes : Application comparative]
\begin{center}
\begin{tabularx}{\linewidth}{|l|X|X|X|}\hline
\textbf{Régime réel} & \textbf{Aristote} & \textbf{Type moderne} & \textbf{Justification brève} \\ \hline
Royaume-Uni & République au sens aristotélicien du gouvernement de tous pour le bien commun (monarchie constitutionnelle où le roi ne gouverne pas) & Parlementaire & Le Premier ministre est responsable devant le Parlement ; le monarque règne sans gouverner. \\ \hline
États-Unis & République (gouvernement de tous, intérêt général) & Présidentiel & Séparation stricte : le Président élu ne dépend pas du Congrès et ne peut en être renversé. \\ \hline
Cameroun & République (souveraineté du peuple, art. 2 de la Constitution de 1996) & Présidentiel (régime à forte dominance présidentielle) & Le Président, élu au suffrage universel, nomme le Gouvernement qui n'est pas responsable devant le Parlement. \\ \hline
France V\textsuperscript{e} République & République (gouvernement de tous pour l'intérêt général) & Mixte (semi-présidentiel) & Double légitimité : Président élu au suffrage universel et Gouvernement responsable devant l'Assemblée nationale. \\ \hline
\end{tabularx}
\end{center}

\textbf{Méthode :} pour Aristote, on demande \emph{qui gouverne} (un, quelques-uns, tous) et \emph{pour qui} (intérêt général ou particulier) ; pour la typologie moderne, on examine les \emph{rapports entre l'exécutif et le législatif} : séparation stricte (présidentiel), collaboration avec responsabilité du Gouvernement devant le Parlement (parlementaire), ou combinaison des deux (mixte).
\end{corrige}

\begin{attention}[Points de vigilance — Exercice 8]
Ne pas classer le Royaume-Uni comme « monarchie » au sens aristotélicien du pouvoir effectif : le roi n'y exerce pas le pouvoir réel. Ne pas confondre « régime présidentiel » (catégorie juridique) et « régime autoritaire » (catégorie politique).
\end{attention}

\begin{corrige}[Exercice 9 — Sujet de Dissertation : « L'État est-il nécessaire à la liberté des citoyens ? »]
\textbf{Corrigé détaillé (plan et contenu attendus).}

\textbf{Introduction.} \emph{Accroche} : « On ne peut pas vivre avec les gens, on ne peut pas vivre sans eux » (formule traditionnelle) : la vie en société impose des règles. \emph{Définitions} : l'État (institution souveraine détenant le monopole de la violence légitime, Weber) ; la liberté (pouvoir de faire ce que les lois permettent, Montesquieu). \emph{Problème} : l'État impose des lois qui contraignent ; mais sans loi, la liberté du plus fort écrase celle du plus faible. \emph{Problématique} : l'État est-il l'ennemi de la liberté, sa condition, ou les deux selon sa forme ? \emph{Annonce du plan}.

\textbf{Première partie — L'État comme condition de la liberté.} Sans État, l'état de nature est insécurité : pour Hobbes, « l'homme est un loup pour l'homme » et la vie y est « solitaire, misérable, dangereuse, animale et brève » ; la liberté sans loi n'est que la liberté du plus fort. Le contrat social institue un pouvoir qui garantit la sécurité (Hobbes), protège les droits naturels — vie, liberté, propriété (Locke) — et fonde la liberté civile : « obéir à la loi qu'on s'est prescrite, c'est être libre » (Rousseau). \emph{Exemple camerounais} : la Constitution de 1996 garantit les libertés fondamentales ; le code de la route protège la vie de tous en contraignant chacun.

\textbf{Deuxième partie — L'État comme menace pour la liberté.} Le pouvoir tend à s'étendre : « tout homme qui a du pouvoir est porté à en abuser » (Montesquieu). La loi est une contrainte ; la raison d'État peut justifier l'arbitraire ; les régimes totalitaires montrent l'État devenu destructeur des libertés (surveillance, parti unique, terreur). L'anarchisme (Proudhon) voit dans l'État l'ennemi par essence de la liberté. \emph{Exemple} : les dérives des états d'urgence prolongés ou de la censure montrent que la puissance publique peut étouffer la liberté qu'elle devait protéger.

\textbf{Troisième partie — Synthèse : la liberté par l'État de droit démocratique.} Tout dépend de la \emph{forme} de l'État. Dans l'État de droit, le pouvoir est lui-même soumis à la loi : hiérarchie des normes, séparation des pouvoirs (Montesquieu), contrôle constitutionnel, pluralisme politique et participation citoyenne. La liberté n'est pas l'absence d'État mais le produit d'un État limité et contrôlé : distinction de Benjamin Constant entre la liberté des Anciens (participation collective) et celle des Modernes (indépendance individuelle garantie par le droit). \emph{Exemples camerounais} : le Conseil constitutionnel, la décentralisation (loi n°2019/024), le pluralisme partisan depuis 1990.

\textbf{Conclusion.} L'État est nécessaire à la liberté à condition d'être un État de droit démocratique : la liberté ne s'exerce que \emph{dans} et \emph{par} la loi commune. \emph{Ouverture} : la vigilance citoyenne est-elle le dernier rempart de la liberté ?

\textbf{Barème indicatif (sur 20) :} Introduction (accroche, définitions, problématique, plan) : 4 pts ; Première partie (thèse argumentée + exemple) : 4 pts ; Deuxième partie (antithèse argumentée + exemple) : 4 pts ; Synthèse (dépassement, État de droit) : 4 pts ; Conclusion et ouverture : 2 pts ; Langue, cohérence, exemples camerounais : 2 pts.
\end{corrige}

\begin{attention}[Points de vigilance — Exercice 9]
La problématique doit naître d'une \emph{tension réelle} (contrainte vs protection), non d'une simple question de cours. Chaque partie doit contenir au moins un auteur et un exemple. Les exemples camerounais sont explicitement exigés par le sujet : leur absence est pénalisée.
\end{attention}

\begin{corrige}[Exercice 10 — Explication de Texte : Rousseau, \emph{Du Contrat social}]
\textbf{1. Thèse.} La légitimité du pacte social repose sur la soumission de chacun à la volonté générale : être contraint d'obéir à cette volonté, c'est être contraint à sa propre liberté, car la loi commune seule arrache le citoyen à toute dépendance personnelle.

\textbf{2. Les trois arguments.}
\begin{itemize}
    \item \emph{Le pacte} : le contrat social contient tacitement l'engagement d'obéir à la volonté générale ; sans cet engagement, le pacte serait un « vain formulaire », sans force obligatoire.
    \item \emph{La contrainte libératrice} : le corps social peut forcer le récalcitrant, car cette contrainte le « garantit de toute dépendance personnelle » : n'être soumis qu'à la loi, c'est n'être l'esclave de personne.
    \item \emph{La légitimité} : le pacte transforme des engagements civils, nuls s'ils sont unilatéraux, en obligations « volontaires, légitimes et sacrées » : la contrainte de l'État est juste parce qu'elle découle du consentement de chacun.
\end{itemize}

\textbf{3. « On le forcera d'être libre ».} La formule est paradoxale car elle unit la force et la liberté. Elle s'explique par la distinction entre la \cle{volonté particulière} (mes intérêts privés, mes passions immédiates) et la \cle{volonté générale} (l'intérêt commun, qui est aussi mon intérêt en tant que citoyen). Celui qui refuse la loi suit sa volonté privée contre sa volonté de citoyen ; le contraindre à la loi, c'est donc le ramener à sa liberté véritable, qui est autonomie : obéir à la loi qu'on s'est prescrite.

\textbf{4. Portée et actualité.} Cette conception fonde l'État de droit moderne : une loi est légitime non parce qu'elle est imposée, mais parce qu'elle exprime la volonté générale via le suffrage et la représentation ; l'obéissance au droit est alors un devoir de citoyen, non une soumission. \emph{Illustrations} : le code de la route au Cameroun contraint chaque conducteur (ceinture, limitations), mais il protège la vie de tous, y compris la sienne ; l'obligation vaccinale limite un choix individuel au nom de la santé publique, intérêt commun. Dans les deux cas, la contrainte légale réalise la liberté collective.

\textbf{Barème indicatif (sur 20) :} Thèse dégagée : 3 pts ; Analyse des trois arguments : 6 pts ; Explication du paradoxe (volonté générale / particulière) : 5 pts ; Portée et exemples actuels : 4 pts ; Clarté et exactitude des références : 2 pts.
\end{corrige}

\begin{attention}[Points de vigilance — Exercice 10]
Ne pas lire « on le forcera d'être libre » comme une apologie de la dictature : la contrainte n'est légitime que si la loi émane de la volonté générale, donc du peuple souverain. Toujours distinguer volonté générale et volonté de tous (somme des volontés particulières).
\end{attention}

\begin{corrige}[Exercice 11 — Analyse de Situation Complexe]
\textbf{1. La conception de M. A.} M. A défend une conception \emph{absolutiste et plébiscitaire} du pouvoir : l'élection serait un blanc-seing conférant tous les pouvoirs jusqu'au scrutin suivant. Il commet la confusion \cle{État / Gouvernement} (le président « incarne » le peuple) et rejoint la thèse hobbesienne de la souveraineté absolue, voire la logique de la raison d'État. C'est une conception de la légitimité réduite à la force du mandat, proche du « pouvoir personnel » dénoncé par la tradition libérale.

\textbf{2. Les trois principes mobilisés par Mme B.}
\begin{itemize}
    \item La \cle{souveraineté populaire} : le pouvoir appartient au peuple, qui ne fait que le déléguer (art. 2 de la Constitution de 1996) ; le président est un mandataire, non un propriétaire du pouvoir.
    \item La \cle{séparation des pouvoirs} (Montesquieu) : l'exécutif revient au Président, le législatif au Parlement, le judiciaire aux tribunaux ; aucun ne peut s'approprier les compétences des autres.
    \item La \cle{hiérarchie des normes} et le principe de légalité : « nul n'est au-dessus de la loi » ; la Constitution s'impose à tous, y compris au chef de l'État — c'est la définition de l'\cle{État de droit}.
\end{itemize}

\textbf{3. Discussion : la force du mandat suffit-elle ?} Il faut distinguer la \cle{légitimité d'origine} et la \cle{légitimité d'exercice}. L'élection confère la première : le gouvernant tient son titre du suffrage universel, ce qui le distingue du tyran. Mais la légitimité d'origine n'est pas un chèque en blanc : la légitimité d'exercice exige le respect continu de la Constitution, le contrôle juridictionnel, la reddition des comptes et la possibilité de l'alternance. Un pouvoir élu qui viole la loi détruit sa propre légitimité : comme le montre Rousseau, la force (même électorale) ne fait pas le droit ; et Locke reconnaît aux citoyens un droit de résistance contre le pouvoir qui trahit sa mission. La majorité électorale fonde le droit de gouverner, jamais celui de gouverner sans limites.

\textbf{4. Conclusion : qu'est-ce qu'un Gouvernement responsable ?} Dans un régime présidentiel comme celui du Cameroun, le Gouvernement n'est pas politiquement responsable devant le Parlement (pas de motion de censure renversant le Président). Mais il demeure \emph{constitutionnellement responsable} : soumis à la Constitution, au contrôle de la légalité de ses actes par les juridictions, à la Haute Cour de justice pour les infractions commises dans l'exercice des fonctions, et \emph{moralement responsable} devant le peuple souverain, qui juge périodiquement par le suffrage. La responsabilité est ainsi la contrepartie nécessaire de tout pouvoir légitime.

\textbf{Barème indicatif (sur 20) :} Qualification de la position de M. A : 4 pts ; Identification des trois principes chez Mme B : 6 pts ; Discussion structurée (légitimité d'origine / d'exercice) : 6 pts ; Conclusion sur le Gouvernement responsable : 3 pts ; Présentation : 1 pt.
\end{corrige}

\begin{attention}[Points de vigilance — Exercice 11]
Ne pas écrire que le Gouvernement camerounais est « responsable devant le Parlement » au sens parlementaire : c'est faux dans un régime présidentiel. Bien distinguer responsabilité politique (parlementarisme), responsabilité constitutionnelle et juridictionnelle, et responsabilité devant l'électeur. La distinction légitimité d'origine / légitimité d'exercice doit être explicitement nommée et définie.
\end{attention}

\newpage

\coursec{Fiche bilan}

\begin{popfigure}
\centering
\begin{tikzpicture}[
    mindmap,
    concept color=popblue,
    level 1 concept/.append style={level distance=4.5cm, sibling angle=360/7},
    level 2 concept/.append style={level distance=3cm},
    every node/.style={font=\small, text=white},
    grow cyclic,
    text width=2.8cm,
    align=center
]
\node[concept, text width=3.2cm] {L'État\\et le Pouvoir}
    child[concept color=poporange] { node[concept, align=center]{1. Définition\\\& Éléments\\Population, Territoire,\\Gouvernement, Souveraineté} }
    child[concept color=popgreen] { node[concept, align=center]{2. Fondements\\\& Légitimité\\Force, Droit, Contrat\\social, Reconnaissance} }
    child[concept color=poppurple] { node[concept, align=center]{3. Formes de\\Gouvernement\\Monarchie, République,\\Démocratie, Dictature,\\Séparation des pouvoirs} }
    child[concept color=poppink] { node[concept, align=center]{4. Rôle\\\& Finalité\\Ordre, Justice, Bien\\commun, Services,\\Développement} }
    child[concept color=popgold] { node[concept, align=center]{5. Souveraineté\\Interne / Externe\\Moniste / Pluraliste\\Limites (Droit int.)} }
    child[concept color=popblue!70!black] { node[concept, align=center]{6. Citoyenneté\\\& État de droit\\Droits / Devoirs,\\Constitution, Loi,\\Contrôle juridictionnel} }
    child[concept color=popmuted] { node[concept, align=center]{7. Problèmes\\actuels\\État failli, Corruption,\\Gouvernance, Décentralisation,\\Mondialisation} };
\end{tikzpicture}
\legende{Carte mentale récapitulative : les 7 piliers de la leçon sur l'État et le pouvoir.}
\end{popfigure}

\begin{bilan}

\courssub{Définitions clés à connaître par cœur}

\begin{itemize}
    \item \textbf{L'État} : Institution politique souveraine, dotée d'un territoire, d'une population, d'un gouvernement et de la capacité de faire la loi (Monopole de la violence physique légitime \emph{— Weber}).
    \item \textbf{La Souveraineté} : Pouvoir suprême, absolu, inaliénable, imprescriptible, indivisible (Bodin). Distinction : \emph{interne} (autorité sur le territoire) / \emph{externe} (indépendance internationale).
    \item \textbf{Le Gouvernement} : Ensemble des organes qui exercent le pouvoir exécutif (administration, police, armée) ; distinct de l'État (permanent) car il est temporaire et révocable.
    \item \textbf{L'État de droit} : État soumis à la loi, où la séparation des pouvoirs est effective, les droits fondamentaux garantis et la justice indépendante.
    \item \textbf{La Légitimité} : Croyance en la justesse du pouvoir (Weber) : \emph{traditionnelle} (coutume), \emph{charismatique} (chef providentiel), \emph{légale-rationnelle} (loi, élections).
    \item \textbf{Le Contrat social} : Fondement théorique (Hobbes, Locke, Rousseau) : les hommes quittent l'état de nature pour créer l'État par convention, source de la légitimité.
\end{itemize}

\courssub{Auteurs \& Thèses de référence (Tableau de synthèse)}

\begin{center}
\begin{tabularx}{\linewidth}{|l|X|X|}\hline
\rowcolor{popblueL}
\textbf{Auteur} & \textbf{Thèse centrale sur l'État} & \textbf{Concept clé} \\ \hline
\textbf{Hobbes} (\emph{Léviathan}) & État absolu nécessaire pour fuir la guerre de tous contre tous (\emph{homo homini lupus}). & Contrat d'aliénation totale ; Souveraineté indivisible. \\ \hline
\textbf{Locke} (\emph{Gouvernement civil}) & État limité : protège vie, liberté, propriété. Droit de résistance si tyrannie. & Contrat de confiance ; Séparation pouvoirs (exécutif/législatif) ; Propriété naturelle. \\ \hline
\textbf{Rousseau} (\emph{Contrat social}) & Souveraineté = Peuple (Volonté générale). Gouvernement = simple mandataire. & Volonté générale vs Volonté de tous ; Légitimité = consentement. \\ \hline
\textbf{Montesquieu} (\emph{Esprit des lois}) & « Pour qu'on ne puisse abuser du pouvoir, il faut que le pouvoir arrête le pouvoir ». & Séparation stricte des 3 pouvoirs (Législatif, Exécutif, Judiciaire). \\ \hline
\textbf{Marx} (\emph{Manifeste, Critique droit}) & État = instrument de domination de classe (bourgeoisie). Doit « dépérir » en communisme. & Infrastructure/Superstructure ; Dictature du prolétariat (transition). \\ \hline
\textbf{Weber} (\emph{Économie et société}) & Définition sociologique : « Monopole de la violence physique légitime ». & 3 types de légitimité ; Bureaucratie rationnelle. \\ \hline
\end{tabularx}
\end{center}

\courssub{Méthodes express (Bac : Dissertation / Commentaire / Cas pratique)}

\begin{methode}[Analyser une situation concrète (Cas pratique / Explication de texte)]
\begin{enumerate}
    \item \textbf{Identifier} la notion centrale (Souveraineté, Légitimité, Séparation des pouvoirs, État de droit, etc.).
    \item \textbf{Qualifier} le régime : Démocratie / Autoritarisme / Totalitarisme ? Unitaire / Fédéral ? Présidentiel / Parlementaire ?
    \item \textbf{Mobiliser} un auteur pertinent (ex: Weber pour la légitimité, Montesquieu pour la séparation des pouvoirs).
    \item \textbf{Appliquer} au contexte camerounais (Constitution 1996 modifiée, décentralisation, rôle du Président, Conseil constitutionnel, élections).
    \item \textbf{Conclure} sur la conformité ou l'écart entre la norme juridique et la réalité politique.
\end{enumerate}
\end{methode}

\begin{methode}[Dissertation : Structure type « L'État et le pouvoir »]
\begin{enumerate}
    \item \textbf{Analyse du sujet} : Définir les termes, poser le problème (ex: « L'État est-il au service de la liberté ou de l'oppression ? »).
    \item \textbf{Plan dialectique} :
    \begin{itemize}
        \item \textbf{Thèse} : L'État garantit l'ordre, la sécurité, les droits (Hobbes, Locke, Kelsen).
        \item \textbf{Antithèse} : L'État peut devenir totalitaire, oppressif, instrument de classe (Marx, Arendt, Tocqueville - tyrannie de la majorité).
        \item \textbf{Synthèse} : L'État de droit démocratique, contrôlé, limité par la Constitution et le peuple (Rousseau, Montesquieu, Droit contemporain).
    \end{itemize}
    \item \textbf{Illustrer} : Exemples historiques (Révolution française, totalitarismes XX\textsuperscript{e}) + Cameroun (transition 1990, décentralisation 2019/2020).
\end{enumerate}
\end{methode}

\end{bilan}

\begin{aretenir}[Checklist avant le Bac : L'État et le Pouvoir]
\begin{itemize}
    \item $\square$ Je sais définir l'État (4 éléments : Population, Territoire, Gouvernement, Souveraineté) et le distinguer de la Nation, de la Société civile, du Gouvernement.
    \item $\square$ Je connais les 3 types de légitimité selon \textbf{Weber} (Traditionnelle, Charismatique, Légale-rationnelle) et je sais les illustrer.
    \item $\square$ Je maîtrise les 3 théories du \textbf{Contrat social} : \textbf{Hobbes} (absolutisme), \textbf{Locke} (libéralisme/propriété), \textbf{Rousseau} (souveraineté populaire/volonté générale).
    \item $\square$ Je sais expliquer la \textbf{Séparation des pouvoirs} (\textbf{Montesquieu}) : Législatif (vote loi), Exécutif (applique), Judiciaire (juge) ; et les régimes (Présidentiel, Parlementaire, Mixte).
    \item $\square$ Je distingue \textbf{Souveraineté interne} (puissance publique) et \textbf{externe} (indépendance) ; je connais les limites (Droit international, UE/UA/CEMAC, Droit humanitaire).
    \item $\square$ Je définis l'\textbf{État de droit} (Hiérarchie des normes, Contrôle de constitutionnalité, Indépendance justice, Droits de l'homme) vs État policier / État providence.
    \item $\square$ Je connais la critique \textbf{marxiste} : État = superstructure, instrument de classe, « comité d'affaires de la bourgeoisie ».
    \item $\square$ Je sais analyser une \textbf{forme de gouvernement} : Monarchie/République, Démocratie/Dictature, Unitaire/Fédéral, Décentralisation (cas du Cameroun : Régions, Communes).
    \item $\square$ Je suis capable d'argumenter sur la \textbf{finalité de l'État} : Sécurité (Hobbes), Liberté/Propriété (Locke), Bien commun/Justice (Aristote/Thomiste), Égalité (Socialisme).
    \item $\square$ Je sais rédiger une \textbf{introduction de dissertation} (Accroche, Définitions, Problématique, Annonce de plan) et une \textbf{conclusion} (Bilan, Ouverture).
    \item $\square$ Je mobilise des \textbf{exemples camerounais précis} : Constitution 1996, Président de la République (chef de l'État, élu au suffrage universel), Parlement (Assemblée + Sénat), Conseil Constitutionnel, CHAN 2020/COUPES D'AFRIQUE (exemple service public), Décentralisation effective (Conseils régionaux 2020).
\end{itemize}
\end{aretenir}

\findecours

\end{document}
